\quellseite{11}

\part{I. Auslegungen über das erste Buch Mosis}

(Fortsetzung.)

\chapter{B. D. M. Luthers Predigten über das 1. Buch Mosis}

sammt einem Unterricht, wie Moses zu lehren ist.\footnote{Luther begann, wie Aurifaber berichtet, am Sonntag Lätare (den 15. März) 1523 mit seinen fortlaufenden Predigten über das erste Buch Mosis. (Weim. Ausg., Bd. 12, S. 435 und 704.) Am Sonntag Judica, den 13. März 1524, stand er bei dem 31. Capitel (Buchwald, Andreas Poachs Sammlung, Band I, S. XXXIII und S. XX), und ging am 19. Sonntage nach Trinitatis, den 2. October 1524, Nachmittags, zur Auslegung des zweiten Buches Mosis über. Alles, was wir über die Veranlassung zu diesen Predigten etc. wissen, beschränkt sich auf das, was Luther in seiner kurzen Anrede an den Leser mitteilt, daher ist es nicht nötig, dasselbe hier noch einmal zu sagen. Unter den „Rottengeistern von Münzerischem Geist“ sind wohl hauptsächlich die Zwickauer Schwärmer zu verstehen. Daß Luther sie mit dem Namen Münzers bezeichnet, hat darin seinen Grund, daß diese Predigten, welche von anderen nachgeschrieben waren, erst im Jahre 1527 mit Luthers Bewilligung herausgegeben wurden, und es daher nahe lag, den Namen Münzers, dessen Geistesverwandte sie waren, auch auf sie anzuwenden. Die diesen Predigten eingefügte Schrift: „Eyn Unterrichtung wie sich die Christen ynn Mosen sollen schicken, gepredigt durch Mar. Luth.“ erschien (allein) schon im Jahre 1526 bei Hans Weiß zu Wittenberg. Dieselbe ist mehrfach nachgedruckt, unter andern zu Nürnberg durch Hans Hergot im Jahre 1527. Diese „Unterrichtung“ gehört nicht eigentlich zu den „Predigten über das erste Buch Mosis“ als ein integrierender Theil, sondern Luther hat sie, wie aus § 28 und § 30 hervorgeht, nach Beendigung des Bauernaufruhrs (am 15. Mai 1525 fand die Schlacht bei Mühlhausen statt) gepredigt. Deshalb ist sie in der Weimarer Ausgabe nicht den Predigten über das erste Buch Mosis, sondern besonders abgedruckt, unter den Schriften vom Jahre 1525. Eine lateinische Uebersetzung der Predigten erschien 1527 zu Ingolstadt bei Joh. Secreius unter dem Titel: In Genesin, Mosi librum sanctissimum, D. Martini Lutheri declamationes. Die erste deutsche Einzelausgabe des Ganzen kam unter dem Titel: „Uber das Erst buch Mose predigte Mart. Luther. sampt einer Unterricht, wie Moses zu leren ist“ bei Georg Rhau (Rhau) in Wittenberg 1527 heraus. Diese Ausgabe wurde noch in demselben Jahre durch Friedrich Peypus in Nürnberg nachgedruckt. Im Jahre 1528 erschien zu Wittenberg bei Georg Rhau eine andere Ausgabe. In den Sammlungen: In der Wittenberger (1553), Bd. V, Bl. 1; in der Jenaer (1566), Bd. IV, Bl. 1 (Die „Unterrichtung“, Bd. III, Bl. 176 b); in der Altenburger, Bd. IV, S. 1; in der Leipziger, Bd. I, S. 1, und in der Erlanger, Bd. 33 und Bd. 34. Wir geben den Text nach der Erlanger, welche den Urdruck bringt, unter Vergleichung der Wittenberger und der Jenaer. Ueber die Vorrede vergleiche man die Anmerkung in diesem Bande Col. 18.}

Gedruckt vom 15. März 1523 bis Herbst 1524; herausgegeben 1527.

Martinus Luther den christlichen Lesern Gnade und Frieden in Christo.

1. Ich habe eine Zeit daher den Propheten Mosen allhie zu Wittenberg gepredigt, allermeist um der Rottengeister willen, ihren Träumen vorzukommen,\footnote{In den alten Ausgaben: ihre Träume zu verspinnen.} auf daß sie nicht den gemeinen Mann durch Mosen verführeten. Wie denn der Münzerische Geist anfing, aus den Christen Juden zu machen, durch falschen Ruhm und Mißverstand des Alten Testaments. Welches Geistes noch viel Braunes und Schwärmens allenthalben übrig ist, und sich noch immer regt.

2. Solche Predigten sind durch andere Gelehrte aufgefangen, und allhie zusammenbracht, also, daß ich mir’s wohl gefallen lasse, und für meine Predigten erkenne, wiewohl, so meine Feder die Zeit gehabt hätte, möchten sie vielleicht völliger und stärker ausgegangen sein. Doch ist meines Sinnes und Verstands hie genug gegeben. Befehle dieselbigen, wem sie gelüstet zu lesen, welchen ich meinen Dienst damit wohl gönne. Besser zu machen ist niemand verboten. Gottes Gnade sei mit uns, Amen.
\quellseite{12}

\section[Eine Unterrichtung, wie sich die Christen in Mosen]{Eine Unterrichtung, wie sich die Christen in Mosen schicken sollen}

gepredigt durch Martin Luther

(Gepredigt nach dem 15. Mai 1525, ausgegangen 1526.)\footnote{Vergleiche die Anmerkung zu § 28.}

3. Lieben Freunde, ihr habt oft gehört, wie daß nie keine öffentliche Predigt vom Himmel herab geschehen sei, denn nur zweimal, wiewohl Gott sonst oft geredet hat durch und mit den Menschen auf Erden, als durch und mit den heiligen Erzvätern, Adam, Noah, Abraham, Isaak, Jakob und andern mehr, bis auf Mosen. Aber durch und mit diesen hat er nicht geredet mit solcher herrlichen Pracht und äußerlichem Wesen, oder öffentlichem Geschrei und Ausrufen, wie er diese zweimal gethan hat, sondern hat ihnen innerlich das Herz erleuchtet, und durch ihren Mund geredet, wie Zacharias, der Vater Johannis, in seinem Gesang anzeigt, da er [Luc. 1, 70.] spricht: „Als er vor Zeiten geredet hat durch den Mund seiner heiligen Propheten.“

4. Nun, die erste Predigt steht im andern Buch Mosis [Cap. 19 und 20], da sich Gott selber hat hören lassen vom Himmel herab, mit großer Pracht und herrlicher Gewalt, zu der Zeit, da er dem Volk von Israel das Gesetz gab mit Donnern und Blitzen,\footnote{In der Wittenberger und Jenaer: „pligen“. Erlanger: „Bligen“.} mit Rauch, Dampf und sehr starken Posaunen, welches das Volk alles hörte, und darüber zitterte und erschrak.

5. Zum andern hat Gott noch eine andere öffentliche Predigt [Apost. 2, 2—4.] lassen ausgehen durch den Heiligen Geist am Pfingstage. Denn daselbst kam der Heilige Geist auch mit großer Pracht und äußerlichem Ansehen, also, „daß ein schnell Brausen eines gewaltigen Windes vom Himmel kam, und erfüllete das ganze Haus, darinne die Jünger saßen; und man sahe an ihnen die\footnote{Erlanger: „ihre“. Dies ist nicht ein Versehen, sondern „die“ (welches auch die Wittenberger und die Jenaer bieten) ist als Variante Walchs angeführt.} Zungen zertheilt, und als wären sie feurig, und Er setzte sich auf einen jeglichen unter ihnen, und wurden alle voll des Heiligen Geistes, und fingen an zu predigen, und zu reden mit andern Zungen“. Dies geschah mit großer Pracht und herrlicher Gewalt, so, daß die Apostel darnach so gewaltig predigten, daß die Predigten, die jetztund in der Welt gehen, kaum ein Schatten sind gegen ihre Predigten, nämlich nach der äußerlichen Pracht und Wesen. Denn sie redeten mit allerlei Zungen, und thaten große Wunderzeichen, wie das Lucas in den Geschichten der Apostel beschreibt.

6. Aber durch die jetzigen Prediger läßt er sich weder hören noch sehen, es geht nicht öffentlich zu vom Himmel herab. Darum habe ich [§ 3] gesagt, es sind nur zwei sonderliche und öffentliche Predigten, die man gesehen und gehöret hat vom Himmel herab. Wiewohl Gott der Vater auch redete mit Christo vom Himmel, da er im Jordan getauft ward, und auf dem Berge Thabor [Matth. 3, 17. 17, 5.]; aber das geschah nicht vor der Gemeinde.

7. Die andere Predigt wollte er senden in die Welt, die zuvor durch den Mund und in den Büchern der heiligen Propheten angezeigt war [Röm. 1, 2.]. Er wird nicht mehr also öffentlich reden mit Predigten, sondern zum dritten wird er selber persönlich kommen mit göttlicher Herrlichkeit, daß alle Creaturen vor ihm werden zittern und beben, und er wird ihnen nicht mehr predigen, sondern sie werden ihn selber sehen und fühlen.

8. Die erste Predigt und Lehre ist das Gesetz Gottes, die andere das Evangelium. Diese zwei Predigten kommen nicht überein, darum muß man guten Verstand darüber haben, daß man sie wisse zu unterscheiden, und wisse, was das Gesetz sei, und was das Evangelium. Das Gesetz gebeut und fordert von uns, was wir thun sollen, ist allein auf unser Thun gerichtet, und steht im Fordern. Denn Gott spricht durch das Gesetz: Das thue, das lasse, das will ich von dir haben. Das Evangelium aber predigt nicht, was wir thun oder lassen sollen, fordert nichts von uns, sondern wendet es um, thut das Widerspiel, und sagt nicht: Thue dies, thue das; sondern heißt uns nur den Schooß herhalten und nehmen, und spricht: Siehe, lieber Mensch,
\quellseite{13}
das hat dir Gott gethan: er hat seinen Sohn für dich ins Fleisch gesteckt, hat ihn um deinetwillen erwürgen lassen, und dich von Sünde, Tod, Teufel und Hölle errettet; das glaube und nimm es an, so wirst du selig. Also sind zweierlei Lehren und zweierlei Werke, Gottes und des Menschen. Und wie wir und Gott voneinander geschieden sind, also sind auch die zwei Lehren weit voneinander geschieden. Denn das Evangelium lehrt allein, was uns von Gott geschenkt ist, nicht, was wir Gott geben und thun sollen, wie das Gesetz pflegt zu thun.

9. Sie wollen mir sehen, wie die erste Predigt sei erschollen, und mit was Pracht Gott das Gesetz habe gegeben auf dem Berge Sinai. Er hat ihm den Ort sonderlich dazu erwählt, daß er da hat wollen gesehen und gehört werden; nicht, daß Gott also geredet habe, denn er hat keinen Mund, Zunge oder Lippen, wie wir, aber der den Mund aller Menschen geschaffen und gemacht hat, kann auch die Sprache und Stimme machen. Denn niemand könnte ein Wort reden, Gott gebe es denn ihm zuvor, wie der Prophet sagt: es wäre unmöglich zu reden, Gott gebe es denn uns zuvor in den Mund. Also ist die Sprache, Rede und Stimme eine Gabe Gottes, wie andere Gaben, als, die Frucht an den Bäumen. Der nun den Mund geschaffen hat, und legt die Sprache darein, kann auch die Sprache machen, ob schon kein Mund vorhanden ist.

10. Nun\footnote{„Nun“ fehlt in der Wittenberger.} die Worte, die hier in Mose geschrieben stehen, sind geredet worden durch einen Engel; nicht, daß allein ein Engel da sei gewesen, sondern eine große Menge, und ein unzähliges Heer, die da Gott gedient haben, und gepredigt vor dem Volk Israel auf dem Berge Sinai. Der Engel aber, der hie geredet hat, und der die Worte führt, redet gleich als ob Gott selber redete und spräche: „Ich bin der Herr dein Gott, der ich dich aus Egyptenland, aus dem Diensthause, geführet habe“ etc., 2 Mos. 20, 2. Gleich als ob Petrus oder Paulus an der Statt Gottes\footnote{So die Wittenberger. Jenaer: an Gottes Statt; Erlanger: an der Gottes-Statt.} redeten und sprächen: „Ich bin euer Gott, der ich euch will selig machen durch meinen allerliebsten Sohn etc. Paulus zum Galatern Cap. 3, 19. spricht, daß das Gesetz geordnet sei durch die Engel, das ist, es sind Engel verordnet gewesen, daß sie an Statt Gottes das Gesetz Gottes dem Volk Israel gäben, und Moses, als ein Mittler, sollte es empfangen von den Engeln. Das sage ich darum, daß ihr wisset, wer das Gesetz gegeben habe. Er hat es aber alles darum gethan, daß er die Juden damit zwingen, fassen und eintreiben wollte.

11. Was das aber für eine Stimme gewesen sei, mögt ihr wohl bedenken. Es ist eine Stimme wie eines Menschen Stimme gewesen, also, daß man sie wohl gehört hat; die Silben und Buchstaben haben daher geklungen, daß es das leibliche Ohr hat mögen fassen. Es ist aber eine tapfere, herrliche und große Stimme gewesen, wie im fünften Buch Mosis Cap. 4, 12. stehet, da er spricht, daß sie die Stimme haben gehöret, und haben keinen Menschen gesehen; sondern sie haben eine starke Stimme gehört, denn er hat eine starke Stimme geführt, als, wenn wir im Dunkeln eine Stimme von einem hohen Thurm oder Dache hörten, und sähen doch niemand, sondern hörten allein eine starke Stimme eines Mannes; und darum wird es auch genannt eine Stimme Gottes, daß sie über eine menschliche Stimme gewesen ist.

12. Nun werdet ihr hören, wie sich Gott geschickt hat zu der Stimme, damit er sein Volk wollte bewegen und wacker machen. Denn er hatte im Sinn, das äußerliche geistliche Regiment anzufahen. Denn zuvor [2 Mos. 18, 14. ff.] hat der Text gesagt, wie Moses durch Rath seines Schwähers Jethro das weltliche Regiment eingesetzt, Hauptleute und Richter verordnet hat. Ueber das ist noch ein geistlich Regiment, in welchem Gott regiert in den Herzen der Menschen; und dasselbige Reich kann man nicht sehen, denn es steht allein im Glauben, und wird währen bis an jüngsten Tag. Das sind nun zwei Reiche: [das] weltliche, das mit dem Schwert regiert, und äußerlich gesehen wird; das geistliche regiert allein mit Gnaden und Vergebung der Sünden. Und dasselbige Reich sieht man nicht mit leiblichen Augen, sondern wird allein mit dem Glauben gefaßt. Zwischen denen beiden Reichen ist noch ein ander Reich in die Mitte gesetzt, halb geistlich und halb weltlich; das fasset die Juden mit Geboten und äußerlichen Ceremonien, wie sie sich halten sollen gegen Gott und den Menschen vor der Welt in äußerlichem Wesen.
\quellseite{14}

\subsection[Das Gesetz Mosis bindet die Heiden nicht, sondern]{Das Gesetz Mosis bindet die Heiden nicht, sondern allein die Juden}

13. Das Gesetz Mosis geht die Juden an, welches uns forthin nicht mehr bindet. Denn das Gesetz ist allein dem Volk Israel gegeben, und Israel hat es angenommen für sich und seine Nachkommen, und die Heiden sind hie ausgeschlossen; wiewohl die Heiden auch etliche Gesetze gemein haben mit den Juden, als, daß ein Gott sei, daß man niemand beleidige, daß man nicht ehebreche noch stehle, und dergleichen andere mehr; welches alles ist ihnen natürlich in das Herz geschrieben, und haben’s nicht vom Himmel herab gehört, wie die Juden. Darum dieser ganze Text geht die Heiden nicht an. Das sage ich um der Schwärmergeister willen. Denn ihr jetzt und hört, wie sie den Mosen lesen, ziehen ihn hoch an, und bringen hervor, wie Moses das Volk mit Geboten habe regiert, wollen klug sein, wollen etwas Weiteres wissen, denn in dem Evangelio begriffen ist, achten für klein den Glauben, bringen etwas Neues auf, rühmen sich und geben vor, es stehe im Alten Testament, wollen nach den Buchstaben des Gesetzes Mosis das Volk regieren, als ob man’s vor nie gelesen habe.

14. Das wollen aber wir nicht gestehen, wollte eher mein Leben lang nicht mehr predigen, ehe ich Mosen wieder einlassen wollte, und Christum lassen uns aus dem Herzen reißen. Wir wollen Mosen nicht für einen Regenten oder Gesetzgeber mehr haben, ja, Gott will es auch selber nicht haben. Moses ist ein Mittler und ein Gesetzgeber gewesen des jüdischen Volks allein, denen hat er das Gesetz gegeben. Man muß also den Rottengeistern das Maul stopfen, die da sagen: Also spricht Moses, da stehet’s im Mose geschrieben, und dergleichen. So sprich du: Moses geht uns nicht an. Wenn ich Mosen annehme in einem Gebot, so muß ich den ganzen Mosen annehmen; also würde daraus folgen, wenn ich Mosen zum Meister und Gesetzgeber annähme, so müßte ich mich lassen beschneiden, die Kleider waschen nach jüdischer Weise, und also essen und trinken, mich kleiden, und solches Wesen alles halten, wie den Juden im Gesetz geboten war. Also wollen wir Mosen nicht halten noch annehmen. Moses ist todt, sein Regiment ist aus gewesen, da Christus kam; er dient weiter hieher nicht.

15. Daß aber Moses die Heiden nicht binde, mag man aus dem Text zwingen im andern Buch Mosis, Cap. 20, 2., da Gott selber spricht: „Ich bin der Herr dein Gott, der ich dich aus Egyptenland, aus dem Diensthause, geführt habe.“ Aus dem Text haben wir klar, daß uns auch die zehn Gebote nicht angehen; denn er hat uns je nicht aus Egypten geführt, sondern allein die Juden. Die Rottengeister wollen uns Mosen auf den Hals legen mit allen Geboten; das wollen wir lassen. Mosen wollen wir halten für einen Lehrer, aber für unsern Gesetzgeber wollen wir ihn nicht halten, es sei denn, daß er gleichstimme mit dem Neuen Testament, und dem natürlichen Gesetze. Darum ist es hell genug, daß Moses der Juden Gesetzgeber ist, und nicht der Heiden. Denn in diesem Text hat Moses den Juden also ein Zeichen gegeben, daß sie Gott sollen erzeigen, wenn sie ihn anrufen, als den Gott, der sie aus Egypten geführt habe. Die Christen haben ein ander Zeichen, dabei sie Gott fassen, als den, „der ihnen seinen Sohn gemacht habe zur Weisheit, und zur Gerechtigkeit, und zur Heiligung, und zur Erlösung“, 1 Cor. 1, 30.

16. Item, man kann es bewähren aus dem dritten Gebot, daß Moses die Heiden, noch die Christen nicht angeht. Denn Paulus und das Neue Testament hebt den Sabbath auf, daß man greifen mag, daß der Sabbath den Juden allein gegeben ist, denen es ein streng Gebot war. Die Propheten haben das auch angezeigt, daß der Juden Sabbath sollte aufgehoben werden. Jesaias spricht Cap. 66, 23.: Wenn der Heiland wird kommen, so wird eine solche Zeit sein, ein Sabbath am andern, ein Neumond am andern etc. Als wollte er sprechen: Es wird alle Tage Sabbath sein, es wird ein solch Volk sein, die keinen Unterschied der Tage werden haben. Denn im neuen Testament liegt der Sabbath nieder, nach der groben, äußerlichen Weise, denn es ist alle Tage Feiertag etc.

17. Wenn nun dir einer Mosen vorhält mit seinen Geboten, und will dich bringen, die zu halten, so sprich: Gehe hin zu den Juden mit deinem Mose, ich bin kein Jude, laß mich unverworren mit Mose. Wenn ich Mosen annehme in einem Stück (spricht Paulus zum Galatern Cap. 5, 3.), so bin ich schuldig das ganze Gesetz zu halten. Denn kein Pünktlein geht uns an im Mose.
\quellseite{15}

18. Möchte nun einer sagen: Warum predigst du denn Mosen, so er uns nicht angeht? Antwort: Dazu will ich Mosen behalten und nicht unter die Bank stecken, denn ich finde dreierlei in Mose, die uns auch nützlich sein können. Zum ersten: Die Gebote, dem Volk Israel gegeben, die das äußerliche Wesen betreffen, laß ich fahren, sie zwingen noch dringen mich nicht, die Gesetze sind todt und ab, ohne solche ichs\footnote{ichs = ich sie.} gerne und willig annehmen will aus dem Mose. Also, wenn ich spräche: Also hat Moses regiert, es dünkt mich sein fein, ich will ihm in dem oder dem Stück folgen. Ich wollte wohl gerne, daß die Herren regierten nach dem Exempel Mosis, und wenn ich Kaiser wäre, wollte ich daraus ein Exempel nehmen der Satzungen; nicht, daß mich Mosis sollte zwingen, sondern daß mir’s frei wäre, ihm nachzuthun, und ein solch Regiment zu führen, wie er regieret hat.

19. Als, mit dem Zehntengeben, das ist ein recht fein Gebot. Denn mit dem Zehntengeben würden aufgehoben alle anderen Zinse, und wäre auch dem gemeinen Mann leidlicher zu geben den Zehnten, denn Rente und Gült.\footnote{Gült (Gülte) = Abgabe, Waarensteuer, Grundrente. (Fehlt bei Dietz.)} Als, wenn ich zehn Kühe hätte, gäbe ich eine; hätte ich fünf, gäbe ich nichts; wenn mir wenig auf dem Felde wüchse, gäbe ich wenig, wenn mir viel wüchse, gäbe ich viel; das stände in GOttes Gewalt. Aber also muß ich die heidnischen Zinse geben, und sollte gleich der Hagel alle Früchte erschlagen. Bin ich schuldig hundert Gulden Zinse, so muß ich’s geben, ob gleich keine Frucht auf dem Felde wüchse. Das ist auch des Papsts Decret und Regiment. Es ginge aber gleich zu, wenn es also geordnet wäre: wenn mir viel wüchse, daß ich viel gäbe, wenn wenig wüchse, daß ich wenig gäbe.

20. Item, in Mose ist auch gefasset, daß keiner seinen Acker sollte verkaufen für ein ewig Erbgut, sondern allein bis aufs Halljahr oder Jubeljahr; und wenn das Jahr kam, so kam ein jeglicher zu seinem Acker oder Gütern wieder, die er verkauft hatte, und blieben also die Güter bei der Freundschaft. Also sind andere aus der Maßen schöne Gebote in Mose mehr, die man möchte annehmen, brauchen und im Schwang lassen gehen, nicht, daß man dadurch sollte zwingen, oder gezwungen werden, sondern, wie ich zuvor [§ 18] gesagt habe, der Kaiser möchte ein Exempel daraus nehmen, ein fein Regiment aus dem Mose zu stellen. Wie auch die Römer ein fein Regiment geführt haben, und wie auch der Sachsenspiegel ist, darnach sich dies Land hält. Die Heiden sind dem Mose nicht schuldig, gehorsam zu sein; Mosis ist der Juden Sachsenspiegel.

21. Item, es steht im Mose [5 Mos. 25, 5. 6. Matth. 22, 24.], wenn einer starb ohne Kinder, so sollte der Bruder oder nächste Freund die Frau heimführen, und für eine Hausfrau haben, und dem gestorbenen Bruder oder Freund den Samen erwecken; und das erste Kind ward dem gestorbenen Bruder oder Freund zugerechnet. Und das ist auch ein fein Gebot. Dergleichen Gebote sind noch viel mehr in Mose, die man alle könnte zu einem feinen Regiment herausser glauben, und dadurch Land und Leute ordentlich und ehrbarlich regieren.

22. Wenn nun die Rottengeister kommen, und sprechen: Mosis hat es geboten, so laß du Mosen fahren, und sprich: Ich frage nichts nach dem, was Mosis geboten hat. Ja, sprechen sie, er hat geboten, man soll einen GOtt haben, dem trauen und glauben, nicht bei seinem Namen schwören, Vater und Mutter ehren, nicht tödten, nicht stehlen, nicht ehebrechen, nicht falsch Zeugniß geben, und seines\footnote{Erlanger: eines.} andern Weib noch Gut begehren. Soll man denn das nicht halten? Sprich also: Die Natur hat diese Gesetze auch; die Natur gibt, daß man GOtt soll anrufen, das zeigen auch die Heiden an. Denn es ist nie kein Heide gewesen, er hat seine Abgötter angerufen; wiewohl sie gefehlt haben des rechten GOttes, wie auch die Juden. Denn die Juden haben auch Abgötterei gehabt, wie die Heiden; allein, daß die Juden das Gesetz empfangen haben, die Heiden aber haben’s im Herzen geschrieben, und ist kein Unterschied, wie auch St. Paulus zum Römern Cap. 2, 15. anzeigt: Die Heiden, die kein Gesetz haben, die haben das Gesetz in ihrem Herzen geschrieben. Wie aber die Juden fehlen, also fehlen auch die Heiden.
\quellseite{16}

23. Und derhalben ist es natürlich, GOtt ehren, nicht stehlen, nicht ehebrechen, nicht falsch Zeugniß geben, nicht todtschlagen; und es ist nicht neu, das Moses gebeut. Denn was GOtt vom Himmel gegeben hat den Juden durch Mose, das hat er auch geschrieben in aller Menschen Herzen, beide der Juden und Heiden, allein, daß er’s den Juden, als seinem eigenen erwählten Volke, zu einem Ueberfluß, auch mit einer leiblichen Stimme und Schrift hat schreiben und verkündigen lassen. Also halte ich nun die Gebote, die Moses gegeben hat, nicht darum, daß sie Moses geboten hat, sondern daß sie mir von Natur eingepflanzt sind, und Moses allhie gleich mit der Natur übereinstimmt etc. Aber die andern Gebote in Mose, die allen Menschen von Natur nicht sind eingepflanzt, halten die Heiden nicht, gehen auch sie nicht an, als, von Zehnten und andern, die doch auch schön sind. Ich wollte, wir hätten sie auch, wie ich [S 18] gesagt habe. Das ist nun das Erste, daß ich in Mose sehen soll, nämlich die Gebote, zu welchen ich nichts verbunden bin, denn sofern sie einem jeglichen von Natur sind eingebildet und in sein Herz geschrieben.

\subsection{Das andere Stück in Mose zu merken}

24. Zum andern finde ich in Mose, welches ich aus der Natur nicht habe; das sind nun die Verheißungen und Zusagungen GOttes von Christo. Und das ist das Beste fast in dem ganzen Mose, welches da nicht natürlich in die Herzen der Menschen geschrieben ist, sondern kommt vom Himmel herab. Als, daß GOtt hat verheißen, daß sein Sohn uns Fleisch geboren sollte werden, daß verkündigt uns das Evangelium. Und das sind nun nicht Gebote, fordern auch nichts von uns, daß wir etwas thun oder lassen sollen, sondern es sind tröstliche, fröhliche Verheißungen GOttes, die wir sollen annehmen, und uns festlich darauf verlassen, wider alle Anfechtung der Sünde, des Todes, des Teufels und der Hölle. Und das ist das Vornehmste in Mose, welches uns Heiden auch angehört. Das Erste, nämlich die Gebote, gehen uns nicht an; aber des andern sollen wir mit Herzen wahrnehmen, und Mose darum lesen, daß so treffliche und tröstliche Zusagungen darinnen geschrieben stehen, damit ich meinen schwachen Glauben mag stärken. Denn also geht es im Reich Christi zu, wie ich’s in Mose lese, darinne ich auch den rechten Grund finde.

25. Und also, auf die Weise, soll ich Mosen annehmen, und nicht unter die Bank stecken. Zum ersten, daß er schöne Exempel der Gesetze gibt, die daraus mögen genommen werden, äußerlich Land und Leute sein ordentlich zu regieren. Zum andern sind darinne die Zusagungen GOttes, damit der Glaube gestärkt und erhalten wird. Als, da GOtt zu der Schlange sagt, wie im ersten Buch Mosis, Cap. 3, 15., geschrieben steht: „Ich will Feindschaft setzen zwischen dir und dem Weibe, und zwischen deinem Samen und ihrem Samen; derselbe soll dir den Kopf zertreten, und du wirst ihn in die Ferse beißen.“ Das ist das erste Evangelium und Verheißung von Christo, geschehen auf Erden, daß er sollte Sünde, Tod und Hölle überwinden, und uns von der Schlange Gewalt selig machen, daran Adam glaubte mit allen seinen Nachkömmlingen, davon er auch ein Christ,\footnote{In den alten Ausgaben „Christen“ statt: ein Christ.} und selig worden ist von seinem Falle.

26. Item, Abraham ward von GOtt diese Zusagung gegeben, wie auch im ersten Buch Mosis, Cap. 22, 18. 12, 3., steht, da er zu ihm also sagte: „Durch deinen Samen sollen alle Völker auf Erden gesegnet werden.“ Das war das andere Evangelium von Christo, daß durch den alle Menschen sollen gesegnet und selig werden; wie das St. Paulus zum Galatern, Cap. 3, 8., auslegt.

27. Item, im fünften Buch spricht Moses, Cap. 18, 15: 16., zu dem Volk Israel: „Einen Propheten, wie mich, wird der HErr dein GOtt dir erwecken, aus dir und aus deinen Brüdern, dem sollt ihr gehorchen. Wie du denn von dem HErrn deinem GOtt gebeten hast zu Horeb, am Tage der Versammlung.“ Und bald hernach, V. 18. 19., setzt Moses die Worte, die GOtt geredet hat zu ihm, also: „Ich will ihnen einen Propheten, wie du bist, erwecken aus ihren Brüdern, und meine Worte in seinen Mund geben, der soll zu ihnen reden alles, was ich ihm gebieten werde. Und wer meine Worte nicht hören wird, die er in meinem Namen reden wird, von dem will ich’s suchen.“ Das ist alles von Christo gesagt, daß der sollte eine neue Predigt auf Erden bringen. Der Sprüche sind viel im Alten
\quellseite{17}
Testament, daran sich die gläubigen Juden gehalten haben, welche die heiligen Apostel oft führten und anzogen.

28. Aber unsere Rottengeister fahren zu: Alles, das sie in Mose lesen, sprechen sie: Da redet Gott, das kann niemand leugnen; darum muß man’s halten. Da fällt denn der Pöbel zu: Hui, hat es Gott geredet, wer will dawider reden? Da werden sie denn eingetrieben, wie die Schweine über ein Trögel.\footnote{Trögel = Trog. So die Jenaer. Wittenberger und Erlanger: Trügel.} Unsere lieben\footnote{Wittenberger: „aufrührische“ statt: lieben.} Propheten haben also ins Volk geplaudert: Liebes Volk, Gott hat geheißen sein Volk, daß sie den Amalek zu Tode schlugen; und andere Sprüche mehr. Daraus ist Jammer und Noth kommen, da sind die Bauern aufgestanden, haben keinen Unterschied gemußt, sind also in diesen Irrthum geführt von den tollen Rottengeistern. Wenn da wären gewesen gelehrte Prediger, die da hätten können den falschen Propheten entgegen kommen und ihnen wehren, und also zu ihnen sprechen: Lieben Rottengeister, es ist wahr, Gott hat es Mosi geboten, und hat also zum Volk geredet; aber wir sind nicht das Volk, dazu es der Herr redet. Lieber, Gott hat auch mit Adam geredet, ich bin darum nicht Adam. Er hat Abraham geboten, er sollte seinen Sohn erwürgen; ich bin darum nicht Abraham, daß ich meinen Sohn erwürge. Also hat er auch mit David geredet. Es ist alles Gottes Wort, wahr ist es. Aber, Gottes Wort hin, Gottes Wort her, ich muß wissen und Acht haben, zu wem das Wort Gottes geredet wird. Es ist noch weit davon, daß du das Volk seiest, damit Gott geredet hat. Die falschen Propheten sprechen: Du bist das Volk, Gott redet mit dir. Das beweise mir. Also hätten sie mögen dawider gelegt werden. Aber sie wollten geschlagen sein, und ist also der Pöbel zum Teufel gefahren.\footnote{Das in diesem Paragraph Gesagte beweist, daß der Bauernaufruhr bereits zu Ende war, also diese „Unterrichtung“ nicht eigentlich zu den „Predigten über das erste Buch Mosis“ gehört. Vergl. § 30.}

29. Man muß mit der Schrift säuberlich handeln und fahren. Das Wort ist in mancherlei Weise geschehen von Anfang. Man muß nicht allein ansehen, ob es Gottes Wort sei, ob es Gott geredet habe, sondern viel mehr, zu wem es geredet sei, ob es dich treffe oder einen andern. Da scheidet sich’s denn, wie Sommer und Winter. Gott hat zu David viel geredet, hat ihn heißen dies und jenes thun; aber es geht mich nicht an, es ist auch zu mir nicht geredet. Er kann es wohl zu mir reden, will er’s haben. Du mußt auf das Wort sehen, das dich betrifft, das zu dir geredet wird, und nicht, was einen andern antrifft. Es ist zweierlei Wort in der Schrift: Eins geht mich nicht an, betrifft mich auch nicht. Das andere betrifft mich, und auf dasselbige, das mich angeht, mag ich’s kühnlich wagen, und mich darauf, als auf einen starken Felsen, verlassen; trifft es mich nicht, so soll ich still stehen. Die falschen Propheten fahren zu und sprechen: Liebes Volk, das ist das Wort Gottes. Es ist wahr, können’s auch nicht leugnen; wir sind aber das Volk nicht, zu dem er redet. Gott hat uns auch weder dies noch jenes geheißen, daß er ihm\footnote{So die Wittenberger. Jenaer: „in“ = ihnen. Erlanger: ihn.} zu thun befohlen hat.

30. Die Rottengeister führen zu, wollten etwas Neues aufstreiben, und sagten: Man muß das Alte Testament auch halten. Haben also die Bauern in einen Schweiß geführt, den sie nicht bald werden abwischen. Ja, sie haben das arme Volk, beide an Leib und Gut, an Weib und Kind, verderbt; wie wir, leider, erfahren und gesehen haben. Die tollen Leute meinten, man hätte ihnen solches Wort Gottes verhalten, es hätte ihnen niemand gesagt, daß sie sollten die Gottlosen zu Tode schlagen. Aber es geschieht ihnen recht, sie wollten niemand folgen noch hören. Ich habe es selber gesehen und erfahren, wie toll, rasig und unsinnig sie waren.

31. Darum spricht zu denselbigen Rottengeistern also: Den Mosen und sein Volk laß bei einander, es ist mit ihnen aus, er geht mich nicht an; ich höre das Wort, das mich betrifft. Wir haben das Evangelium. Christus spricht [Marc. 16, 15.]: „Gehet hin und prediget das Evangelium“, nicht allein den Juden, wie Moses, sondern allen Heiden, ja „allen Creaturen“. Mir ist gesagt [Marc. 16, 16.]: „Wer da glaubet und getauft wird, der wird selig.“ Item [Luc. 10, 37.]: Gehe hin und thue deinem Nächsten, wie dir geschehen ist. Die Worte treffen mich auch, denn ich bin eine aus allen Crea
\quellseite{18}
turen. Wenn Christus nicht hätte hingesetzt: „Prediget allen Creaturen“, so wollte ich mich nicht daran kehren, wollte nicht getauft werden, und mich also dagegen halten, wie ich mich jetzund gegen Mosen halte; an den kehre ich mich lauter nichts, er geht auch mich nicht an, denn er nicht mir, sondern allein den Juden gegeben ist. Dieweil aber Christus spricht, man soll das Evangelium: „Wer glaubt und getauft wird, der wird selig werden“, nicht Einem Volk allein, nicht an dem oder an diesem Ort der Welt, sondern allen Creaturen predigen, so ist niemand ausgenommen, sondern es sind alle Creaturen darunter verfaßt; daran darf niemand zweifeln, es soll ihm auch das Evangelium gepredigt werden. Also glaube ich denn dem Worte, es gehe mich auch an, ich gehöre auch unter das Evangelium und in das Neue Testament, darum wage ich’s auf das Wort, und sollte es hunderttausend Hälfe kosten.

32. Den Unterschied sollen wohl merken, lassen und zu Herzen nehmen die Prediger, so andere Leute wollen lehren, ja, alle Christen. Denn es ist die Macht ganz und gar daran gelegen. Wenn es die Bauern also verstanden hätten, wären ihrer viel erhalten worden, und nicht so jämmerlich verführt und verdorben. Und wo wir’s anders werden verstehen, so machen wir Secten und Rotten, wo wir also unter den Pöbel, in das tolle, unverständige Volk, ohn allen Unterschied speien und geifern: Gottes Wort, Gottes Wort. Ja, lieber Gesell, nicht also; es heißt, ob es dir gesagt sei, oder nicht. Gott redet auch wohl zu den Engeln, Holz, Fischen, Vögeln, Thieren und zu allen Creaturen, es geht darum mich nicht an; ich soll auf das sehen, das mich betrifft, das mir gesagt ist, damit er mich mahnt, treibt und fordert.

33. Deß nimm ein Exempel. Wenn ein Hausvater hätte eine Frau, Tochter, Sohn, Magd und Knechte; nun er spräche zum Knechte, und hieße ihn die Pferde anspannen, und ins Holz fahren, den Acker pflügen, und dergleichen Arbeit thun; zu der Magd spräche er, sie solle die Kühe melken, buttern, und dergleichen; zu der Frau aber, sie solle der Küche warten; zur Tochter, sie solle spinnen und das Bette machen. Das alles wären Worte eines Herrn, eines Hausvaters. Wenn nun die Magd zuführe, und wollte mit den Pferden umgehen, wollte gen Holz fahren; der Knecht setzte sich unter die Kühe und wollte melken; die Tochter wollte mit dem Wagen fahren, wollte pflügen; die Frau wollte das Bette machen, wollte spinnen, und die Küche versäumen; und wollten also sprechen: Der Herr hat es geheißen, es ist der Befehl des Hausvaters. Da sollte der Hausvater zufahren, und einen Knüttel nehmen, und sie allzumal auf einen Haufen schmeißen, und sprechen: Wiewohl es mein Befehl ist, so habe ich’s doch dir nicht befohlen, habe einem jeglichen seinen Bescheid gegeben; dabei solltet ihr geblieben sein.

34. Also hält es sich auch mit dem Worte Gottes. Wenn ich mich des wollte annehmen, das er einem andern befohlen hat, und wollte sprechen: Hast du es doch gesagt; sollte er sprechen: Wer weiß dir’s\footnote{Erlanger: dir.} Dank? ich hab es aber dir nicht gesagt. Man muß einen guten Unterschied machen, wenn das Wort Einen trifft, oder alle zumal. Wenn nun der Hausvater spräche: Am Freitag wollen wir Fleisch essen, das wäre ein gemein Wort allen im Hause. Also, was zu Mose durch Gott geredet ist der Gebote halben, trifft allein die Juden; aber das Evangelium geht durch die ganze Welt durch und durch, niemand wird ausgenommen, sondern allen Creaturen wird es vorgetragen. Darum soll sich deß alle Welt annehmen, und also annehmen, als ob es einem jeglichen insonderheit vorgetragen sei. Das Wort Joh. 13, 34.: Wir sollen einander lieb haben, geht mich an; denn es geht alle an, die zum Evangelio gehören.

35. Also lesen wir Mosen darum: nicht daß er uns betreffe, daß wir ihn müssen halten; sondern daß er gleichstimmt\footnote{In den Ausgaben: gleichstimme.} mit dem natürlichen Gesetz, und ist besser gefasset, denn die Heiden immer hätten mögen thun. Also sind die zehn Gebote ein Spiegel unseres Lebens, darinne wir sehen, woran es uns fehlt etc. Die Rottengeister haben auch den Mosen von den Bildern nicht recht verstanden, denn es geht auch allein die Juden an etc. Zum andern, wie jetzund gesagt ist, lesen wir Mosen um der Verheißungen willen, die von Christo lauten, der nicht allein den Juden, sondern auch den Heiden zugehört. Denn durch ihn sollten alle Heiden den Segen und die Gebenedeyung haben, wie Abraham verheißen war [1 Mos. 22, 18.].
\quellseite{19}

\subsection{Das dritte Stück, so in Mose wahrzunehmen ist}

36. Zum dritten lesen wir Mosen von wegen der schönen Exempel des Glaubens, der Liebe und des Kreuzes in den lieben heiligen Vätern, Adam, Abel, Noah, Abraham, Isaak, Jakob, Mose, und also durch und durch. Daran wir lernen sollen, GOtt zu vertrauen und ihn lieben. Herwiederum auch sehen wir die Exempel des Unglaubens der Gottlosen, und des Zorns GOttes, wie GOtt nicht schenkt den Ungläubigen ihren Unglauben; wie er gestraft hat den Kain, Ismael, Esau, die ganze Welt mit der Sündfluth, Sodoma und Gomorra; und dergleichen viel Strafen mehr, so er über die Gottlosen hat gehen lassen. Und die Exempel sind vonnöthen. Denn wiewohl ich nicht Kain bin; doch, wenn ich thue, wie Kain gethan hat, so werde ich gleiche Strafe mit Kain empfangen. Man findet an keinem andern Orte so schöne Exempel, beide vom Glauben und Unglauben, als eben in Mose. Darum soll man Mosen nicht unter die Bank stecken. Und also wird das Alte Testament recht verstanden, so man die schönen Sprüche von Christo aus den Propheten behält, und die schönen Exempel wohl faßt und merkt, und so wir die Gesetze nach unserm Wohlgefallen brauchen, und dieselbigen uns nütze machen.

\subsection{Beschluß}

37. Ich habe gesagt, daß alle Christen, und insonderheit, die andere Leute lehren wollen und das Wort Gottes handeln, daß die sich wohl vorsehen sollen, und den Mosen recht lernen; also, wo er Gebote gibt, daß wir ihn darinne nicht weiter annehmen, denn so ferne er sich mit dem natürlichen Gesetze reimt. Moses sei ein Meister und Doctor der Juden. Wir haben unsern Meister Christum, der uns vorgelegt hat, was wir wissen, halten, thun und lassen sollen. Aber das ist wahr: Moses schreibt, neben den Gesetzen, schöne Exempel des Glaubens und Unglaubens, Strafe der Gottlosen, Erhöhung der Frommen und Gläubigen, und auch die lieblichen und tröstlichen Zusagungen von Christo. Deren sollen wir uns annehmen, wie wir auch thun sollen in den Evangelisten. Als, wenn man liest von den zehn Aussätzigen, das geht mich nicht an, daß er sie heißt zum Priestern gehen, und ihr Opfer thun; das Exempel aber ihres Glaubens geht mich an, daß ich Christo auch, wie sie, glaube.

38. Davon ist nun genug geredet, und ist wohl zu merken, denn es liegt die Macht daran, und haben viel große, treffliche Leute darinne gefehlt, und stoßen sich jetzund viel große Prediger daran, wissen den Mosen nicht zu predigen, können sich nicht wohl darein schicken, sind unsinnig, toben, rasen und wüthen, plaudern ins Volk: Gottes Wort, Gottes Wort, Gottes Wort; verstören die armen Leute, und stoßen sie in die Grube. Es haben viel gelehrte Leute nicht gewußt, wie weit Moses soll gelehrt werden. Origenes, Hieronymus und ihres Gleichen haben nicht klar angezeigt, wie weit uns Moses diene.\footnote{Erlanger: dienet.}

39. Das habe ich wollen zu einem Eingang in Mosen reden, wie man sich darein schicken soll, und wie Moses verstanden und angenommen, und nicht gar unter die Bank soll gesteckt werden, darinne so eine schöne Ordnung und äußerlich Regiment begriffen wird, daß es [eine] Lust ist, ohne das, daß er viel treffliches, schönes Dinges beschreibt, wie ihr gehört habt, welches nicht allein nicht zu verwerfen ist, sondern auch hoch zu achten, und mit ernstem Herzen anzunehmen, als zur Förderung und Stärkung unseres christlichen Glaubens, durch welchen, wie wir, so auch die lieben heiligen Väter sind selig worden.
\quellseite{20}

\section[Vorrede auf die Predigten über das erste Buch Mosis]{Vorrede auf\footnote{Wittenberger und Erlanger: in.} die Predigten über das erste Buch Mosis.\footnote{Diese Vorrede und die Auslegung des ersten Capitels (bis § 19 incl.) erschien schon 1524 in mehreren Einzeldrucken (die Weimarische Ausgabe zählt deren vier auf) unter dem Titel: Ein Sermon und Eingang in das erste Buch Mosis, daß ist, in das Buch der Schöpfung, darinn gründlich angezeigt wird der Artikel des Glaubens, so wir sprechen: Ich glaube in Gott Vater allmächtigen, Schöpfer Himmel und Erdreichs. Martinus Luther. Wittenberg 1523. Dieser Abschnitt, welcher im Wesentlichen mit dem hier gegebenen Text übereinstimmt, ist abgedruckt in der Weimarischen Ausgabe, Bd. XII, S. 438. Unter der Zeile bringt diese Ausgabe aus der Zwickauer Rathsbüchlibibliothek eine Nachschrift Stephan Roths in lateinischer Sprache, welche (wiewohl sie die längere seiner zwei Recensionen ist) wohl kaum den fünften Theil des deutschen Textes in sich begreift. Die in unserer Ausgabe gegebene Redaction dieses Abschnittes ist unstreitig bedeutend besser als die in der Weimarischen Ausgabe gegebene.}}

1. Sintemal es in der Christenheit also geordnet sein soll, daß man immerdar Gottes Wort predige und treibe, und täglich im Schwange gehen lasse (weil an demjenigen allein alle Macht liegt, und ohne das die Seelen nicht leben können), auf daß unter den Christen die heilige Schrift, nämlich beide, das Alte und Neue Testament, jedermann bekannt und läuffig werde, daß wir durch Gottes Wort gerüstet, geharnischt, und gestärket werden im Glauben, zu bestehen wider allerlei Anfechtung und Unglück: so will ich demnach auch ein Buch vor mich nehmen, aus dem Alten Testament, und dasselbige auslegen und erklären, so viel Gott Gnade geben wird, wie wir bisher im Neuen Testament gethan haben, daraus ein jeglicher Christ sehen möge, wie die Schrift allenthalben übereinstimmt, und wie alle Exempel und Historien, ja, die ganze Schrift durch und durch, sich dahin lenken, daß man Christum erkenne, davon wir immer gepredigt haben, und noch predigen.

2. Denn, wiewohl es wahr ist, daß wir genug hätten allein an Einem Spruche aus der Schrift und Gottes Wort, den Glauben zu lehren, so hat uns doch Gott sonderliche Gnade gethan, und mit vielen Exempeln und Sprüchen reichlich überschüttet, auf daß wir ja klärlich sehen, wie es allein am Glauben gelegen ist, und durch so viel Sprüche unsern Glauben stärken, und bei uns gewiß werden, daß wir nicht wanken. Darum, wie wir bisher gesehen haben, wie alle Evangelia nicht mehr denn Ein Ding allenthalben lehren und treiben, also wollen wir dasselbige im Alten Testament auch sehen, ohne daß das Alte Testament etwas mehr darüber hat, daß es die Leute auch nach dem äußerlichen Regiment unterweist, für die, so da nicht Christen sind; wie man dieselbigen halten und regieren soll im äußerlichen Wandel, Regiment und Ordnung.

3. Wir wollen aber vorne anheben, und zum ersten das erste Buch Mosis vor uns nehmen, darinne vornehmlich beschrieben ist, wo alle Creaturen hergekommen sind, sonderlich der Mensch mit der Sünde, mit dem Tod, und der Frömmigkeit, davon alle Welt viel trachtet und dichtet, und dennoch nie hat mögen erlangen.

4. Ehe wir aber dazu greifen, ist vor allen Dingen noth zu wissen, daß das Wort Gottes auf zweierlei Weise gehandelt wird. Einmal von denen, die nicht rechtschaffen glauben, und dasselbe äußerlich allein mit dem Buchstaben, nach der Vernunft und menschlichem Verstand, fassen, einen Gedanken davon machen, wie die Worte lauten. Das andere von denen, die der heilige Geist lehrt, die es nicht mit äußerlichen Gedanken fassen, sondern denen Gott einen rechten Verstand und Erfahrung im Herzen gibt.

5. Der erste Haufe versteht es nimmer recht, wiewohl viel große Schwätzer daraus werden, die sich rühmen, sie habens auch gelesen und verstehen es, meinen, sie können es besser denn die andern alle. Daher auch die Juden verboten haben (wie Hieronymus schreibt), daß dies Buch niemand lesen sollte, so lange bis er wohl geübt, und zu seinen dreißig Jahren kommen wäre. Desgleichen sie auch etliche Bücher mehr ausgezogen haben, für junge Leute nicht zu lesen. Nun aber achte ich nicht, daß solches gesetzt und geordnet sei von groben Heuchlern, sondern halte es dafür, daß es freilich hohe Leute gethan haben. Denn es ist wohl also geschrieben, vorn im ersten Capitel, daß es keine Vernunft begreift.

6. Denn das ist ohne Zweifel der höchste Artikel des Glaubens, darinne wir sprechen: Ich
\quellseite{21}
glaube an GOtt Vater, allmächtigen Schöpfer Himmels und der Erde. Und welcher das rechtschaffen glaubt, dem ist schon geholfen, und ist wieder zurecht gebracht, und dahin kommen, da Adam von gefallen ist. Aber wenig sind ihrer, die so weit kommen, daß sie völliglich glauben, daß er der GOtt sei, der alle Dinge schafft und macht. Denn ein solcher Mensch muß allen Dingen gestorben sein, dem Guten und Bösen, dem Tod und Leben, der Hölle und dem Himmel, und von Herzen bekennen, daß er aus eigenen Kräften nichts vermag.

7. Das ist nun die Ursache, warum die lieben Väter dies Capitel sonderlich nicht haben jedermann lassen lesen und handeln. Darum wollen wir’s auch nicht aufs schärffte handeln, sondern, so viel GOtt Gnade wird geben, einfältig auslegen, und wollen sehen, daß wir mögen treffen den Kern, das ist, wie es Moses gemeint hat, ob wir gleich nicht erlangen, wie es alles zugegangen ist.

8. Denn es haben viel Leute so viel darüber geschrieben, und so mancherlei Verstand gegeben, daß ich nicht weiß, ob sie den rechten Verstand haben. Zum ersten haben etliche also davon geredet, daß alle Dinge, so GOtt geschaffen hat, in einem Augenblick sind vollkommen gewesen. Und nehmen solches aus diesem Spruch 5 Mos. 32, 4.:\footnote{Beide Sprüche sind nach der Vulgata angeführt.} Dei perfecta sunt opera, das ist, „GOttes Werke sind vollkommen und ganz.“ Und aus diesem, Sir. 18, 1.: Qui vivit in aeternum, creavit omnia simul, „der da ewig lebt, hat alle Dinge geschaffen mit einander.“ Und sind viel hohe Leute, als Augustinus und Hilarius, dieser Meinung, daß es in einem Augenblick alles und gar gestanden sei, wie wir’s jetzt sehen, der Himmel mit Sonne, Mond, Sternen, die Erde voll Früchte und lebendiger Thiere, also, daß Anfang und Ende zugleich gestanden seien. Aber wir wollen es nicht also verstehen.

9. Denn ich habe oft gesagt, daß, wer in der heiligen Schrift studiren will, soll je darauf sehen, daß er auf den einfältigen Worten bleibe, wie er immer kann, und sie nicht davon weiche, es zwinge denn irgend ein Artikel des Glaubens, daß man es müsse anders verstehen, denn die Worte lauten. Denn wir müssen deß sicher sein, daß keine einfältigere Rede auf Erden kommen sei, denn das GOtt geredet hat.

10. Darum, wenn Moses schreibt, daß GOtt in sechs Tagen Himmel und Erde, und was darinnen ist, geschaffen habe, so laß es bleiben, daß es sechs Tage gewesen sind, und darfst keine Glosse finden, wie sechs Tage ein Tag sind gewesen. Kannst du es aber nicht vernehmen, wie es sechs Tage sind gewesen, so thue dem Heiligen Geist die Ehre, daß er gelehrter sei, denn du. Denn du sollst also mit der Schrift handeln, daß du denkest, wie es GOtt selbst rede. Weil es aber GOtt redet, so gebühret dir nicht, sein Wort aus Frevel zu lenken, wo du hin willst, es zwinge denn die Noth, einen Text anders zu verstehen, denn wie die Worte lauten, nämlich, wenn der Glaube solchen Verstand, als die Worte geben, nicht leidet.

11. Wenn man nun sagt: GOttes Werke sind vollkommen, darum soll es folgen, daß er nicht so lange mit der Schöpfung umgegangen sei, nämlich sechs ganzer Tage. Antwort: Sie sind ja vollkommen; aber, wie wir sehen, daß wir Menschen nicht bald gemacht werden, sondern das Kind neun Monate lang im Mutterleibe liegt, und dazu nicht eher vollkommen wird, denn bis es groß und männlich ist, so er es doch wohl könnte in einem Augenblick vollkommen machen, desgleichen er auch mit andern Dingen umgeht, so müssen wir davon auch sagen, daß GOtt nicht solche vollkommenen Werke machete. Darum verstehe es also: Wenn es GOtt ausgerichtet und vollendet hat, sind seine Werke wohl vollkommen; aber weil er noch das Werk vor ihm hat, und daran macht, so ist es nicht vollkommen. Nun ist es also vor unsern Augen, daß er immerdar macht und schafft. Darum, gleich als wenn ein Mensch ein Bild macht, ist es noch nicht vollendet, dieweil er noch daran arbeitet: also ist auch GOttes Werk nicht eher vollkommen, denn wenn es gemacht ist. Denn also ist es in der ganzen Welt, daß GOtt von Anfang hat beschlossen, so viel Menschen zu schaffen, und doch täglich immerdar schafft, wiewohl er alle Menschen auf einmal könnte machen. Darum schließt dieser Spruch nicht, daß sie wollen damit bewähren.

12. Desgleichen folgt ihre Meinung auch aus diesem Spruch nicht: „GOtt hat alles mit einander geschaffen“ 2c. Denn er will nur so viel sagen: Es ist einer, der da lebt in Ewigkeit, der hat alle Dinge gemacht, alles, was da ist auf einem Haufen; und sonst niemand. Darum ist
\quellseite{22}
es eben das, daß hier Moses sagt und erzählt, wie GOtt die Welt habe angefangen zu schaffen, und sechs Tage darüber gemacht, und eben gethan, wie er noch thut, wenn er einen Menschen schafft in Mutterleibe. Und wiewohl er hier schwerer und mehr geeilet hat, ist es doch mit Muße zugegangen, und sein eines nach dem andern gemacht worden, bis es alles vollkommen ist worden. Darum soll man hier bei den Worten bleiben, deweil der Glaube nicht zwingt, daß man es anders verstehe. Wenn aber ein Spruch der Schrift wäre, der da klärlich zwänge, daß alle Dinge zugleich und auf einen Augenblick gemacht wären, so müßten wir diesen Text auch anders verstehen. Das ist Ein Stück, das wir zur Vorrede haben sagen wollen.

13. Das andere ist, wie wir [§ 4] gesagt haben, daß ein jeglicher dieser Worte soll mit dem Geist fassen und so viel daraus nehmen, daß GOtt alle Dinge im Himmel und Erden thue, schaffe und wirke, wie der Text auch will. Wer das verstehet, der wird so bald inne, daß er keine Ader regen, auch nicht Einen Gedanken haben kann, GOtt muß es wirken; daß sein Leben gar in seiner Hand nicht steht, sondern bloß in GOttes Hand. Denn, so ich das glaube, daß er habe die ganze Welt aus Nichts gemacht, sondern alles allein aus seinem Worte und Gebote gestanden sei, so muß ich ja bekennen, daß ich auch ein Stück von der Welt und seiner Schöpfung sei. Daraus muß folgen, daß in meiner Macht nicht stehe, eine Hand zu regen, sondern daß allein GOtt alles in mir thue und wirke. Da will es hinaus, und da muß man es hinlenken, so ist der Verstand recht.

14. Wenn du nun denn das also fühlst, so wirst du müssen erschrecken; denn die Natur kann es nicht leiden. Tröstlich aber ist es denen, die im Glauben stehen; denn da ist nichts, das sie stärken und trösten möge, denn daß sie wissen, wie sie gar in GOttes Hand stehen, also, daß er auch die geringsten Gedanken in ihnen wirke. Wo nun solcher Glaube ist, der kann sich gar vor nichts fürchten, und sich auch auf nichts verlassen, weder im Himmel noch auf Erden, weder im Leben noch im Tode, weder in Sünde noch Frömmigkeit, denn allein auf GOtt. Darum, wenn schon die ganze Welt wider mich stünde und mich angriffe, daß ich mitten in ihren Händen wäre, so weiß ich, daß sie doch nichts vermögen, denn so ferne GOtt will. Und wenn schon der Feinde so viel wären, als Sandes am Meer, so sind sie ja GOttes Creatur, so können sie ohne seinen Willen und Sorgen keinen Schaden haben; geschweige, daß sie mir Schaden thun können, er wolle denn. Will er aber, wohl mir! denn ich weiß, daß es sein gnädiger Wille und väterliche Liebe ist.

15. Darum steht ein solcher gläubiger Mensch in solcher Freude und Sicherheit, daß er sich vor keiner Creatur läßt erschrecken, ist aller Dinge ein Herr, fürchtet sich vor keinem Dinge, das ihm möchte zuhanden stehen, denn allein vor GOtt, der im Himmel ist. Wiederum, wenn er in der Welt ein großer Herr wäre, und gesetzt würde über alle Königreiche, daß man ihm gäbe alle Wollust und Freude auf Erden, so nimmt er sich kein nichts an; fragt auch nichts darnach, wenn ihm solches alles wieder genommen würde. Denn er setzt sein Vertrauen nicht auf irgend eine Creatur, sondern allein auf GOtt.

16. Wo aber der Glaube nicht ist, da der Mensch nach seiner Vernunft und Dünkel richtet, und dieser Verstand auch offenbar wird, da ist auch die Hölle selbst, und kann der Mensch nicht größere Marter haben. Denn wie er dort im Glauben ohne Furcht steht, und aller Dinge ein Herr ist, und sich allein in GOttes Hand gibt, also ist es hier wiederum; wenn er nicht glaubt, und doch sieht, daß alle Creaturen in GOttes Gewalt stehen, so ist keine Creatur, die ihn nicht erschrecke, daß er sich vor allen fürchten muß. Denn dieweil GOtt wider ihn ist, so müssen auch alle Dinge wider ihn sein.

17. So geht es denn, wie Moses [3 Mos. 26, 36.] schreibt, daß die Gottlosen auch ein rauschend Blatt, das vom Baum fällt, erschreckt. Da kann das Herz nicht so viel Muths\footnote{Weimarische: muß.} kriegen, daß es eine Mannheit fasse wider ein solch geringe rauschend Blatt: was sollte es thun, wenn der Tod kommt? Wo er hingeht oder sieht, so ist ihm GOtt entgegen, und denkt, er wolle ihn vor den Kopf schlagen. Darum sind ihm diese Worte nichts Anderes denn ein Blitz und Donnerschlag. Derhalben sagt abermal Moses [5 Mos. 28, 65. ff.]: „GOtt wird dir ein verzagt Herz geben, daß du vor Trauern verschmachten wirst, und deines Lebens nimmer sicher sein. Des Morgens wirst du sagen: Wer weiß, ob ich den
\quellseite{23}
Abend erlebe? des Abends wirst du sagen: Ach, möchte ich des Morgens leben!

18. Also ist hier auf beiden Seiten in Mose die Hölle und der Himmel. Die es mit dem Glauben fassen, die lernen, daß sie Gott vertrauen, und sich ihm ganz hingeben, und werden so beherzt,\footnote{In den Ausgaben: geherzt.} daß sie sich vor keinem Dinge fürchten, denn sie wissen, daß es Gott mit ihnen hält. Die es aber mit der Vernunft, ohne Glauben, fassen und fühlen, daß ihnen das Gewissen sagt, daß ihnen Gott feind ist, die können weder Ruhe noch Freude haben. Der andere Haufe, oder die solchen Verstand nicht fassen noch fühlen, können nicht mehr, denn daß sie sagen: Gott hat Himmel und Erde geschaffen; kommt ihnen aber nicht ins Herz, sondern behalten nur die Worte auf der Zunge.

19. Das sind die zwei Stücke, die wir haben zur Vorrede müssen sagen, nämlich zum ersten, daß man die Worte in dem einfältigen, schlechten Verstande bleiben lasse; zum andern, daß man die Worte und den Kern recht fasse, und fühle es im Herzen. Die das nicht können thun, denen ist es verboten zu lesen, daß sie es nicht angreifen; denn sie handeln es ohne Frucht. Dadurch sind nun schon unsere Werke, Verdienste, freier Wille und Vernunft zu Pulver gestoßen, weil keine Creatur nun ihr selbst das geringste Werk thun kann. Nun wollen wir zum Text greifen, und sehen, wie Moses die Schöpfung nach einander beschreibt.

\section{Das erste Capitel}

V. 1. 2. Am Anfang schuf Gott Himmel und Erde. Und die Erde war wüste und leer, und es war finster auf der Tiefe, und der Wind\footnote{Wind oder Geist (In den alten Ausgaben am Rande).} Gottes schwebete auf dem Wasser.

1. Das ist der Anfang dieses Buchs, und ist wahrlich hoch angefangen. Wir aber wollen nicht zu tief hinein greifen. Es ist genug, wenn wir den Glauben daraus können lehren und Gottes Werk verstehen. Aufs erste soll sich ein jeglicher gewöhnen, daß er die Worte nicht anders achte, denn als wären sie gestern geschrieben. Denn es ist vor Gott so kurze Zeit, eben als wären sie jetzund geschehen. Darum ist eitel neu Ding da, und geht dieser Anfang noch immerdar.

2. Aufs andere, den Anfang soll man auch aufs einfältigste verstehen, daß davor nichts gewesen ist, weder Stunde, noch Tage, noch Zeit. Das ist aber auch ein hoher Verstand, daß vor der Zeit keine Zeit gewesen, und doch Gott gewesen ist. Darum ist göttlich und ewiges Leben viel ein ander Ding, denn jemund dies Leben ist, welches immer geht von Tag zu Nacht. Darum, wenn man daran hanget, kann man nicht denken, was Ewigkeit ist. Vor Gott ist der Anfang der Welt ja so nahe, als das Ende; tausend Jahr als Ein Tag; und Adam, der am ersten geschaffen ist, als der letzte Mensch, der da wird geboren werden. Denn er sieht die Zeit also an, wie des Menschen Auge zwei Dinge, die weit von einander sind, in einem Augenblick zusammen bringt. Das sage ich darum, daß nicht jemand einen närrischen Verstand fasse, daß er also vom Anfang rede, als wäre etwas vorhin gewesen, sondern daß da angegangen ist Zeit und Creatur, welches vorhin nicht war.

3. Was ist es aber, daß er sagt: „Gott schuf Himmel und Erde, und die Erde war wüste und leer“? Das ist’s, daß ich vor gesagt habe, daß der allmächtige Gott nicht hat die Welt auf einen Hui\footnote{Weimarische: haw.} geschaffen, sondern Zeit und Weile dazu genommen, und damit umgegangen ist, eben wie er jetzund ein Kind macht. [Er] hat zum ersten das Geringste gemacht, Himmel und Erde, also, daß es noch ungeschaffen, wüst und leer gewesen ist, da noch nichts aufgewiesen noch gewachsen, und nicht geschickt oder gestaltet, noch in eine Form gebracht ist.

4. Hier soll man aber nicht denken, wie die Philosophi, Plato und Aristoteles, gethan haben mit ihren Ideis, sondern aufs einfältigste also, daß da rechtschaffen Himmel und Erde ge
\quellseite{24}
wesen sei, wie er es selbst nennt, und die erste Creatur; aber der beiden keines geschickt sei gewesen, wie es sein sollte. Gleich als ein Kind im Mutterleibe zum ersten nicht gar nichts ist; aber doch nicht zugerichtet, wie ein völlig Kind sein soll. Wie auch ein Rauch nicht nichts ist, doch also hinführt, daß er weder Licht noch Schein hat: also ist auch die Erde zu keinem Dinge geschickt gewesen, und hatte noch keine Form, weder in die Breite noch in die Länge, und war weder Korn, noch Bäume, noch Gras darauf, sondern schlechte, unfruchtbare Erde, als ein Land oder Wüste, da niemand ist, und nichts wächst. Also ist der Himmel auch ungestalt gewesen, und war doch nicht gar nichts.

5. Das zeigt Moses an mit den Worten, als er sagt, daß die Erde wüste und leer sei gewesen, und Finsterniß auf der Tiefe, und der Wind Gottes schwebete auf dem Wasser. Die Finsterniß, Tiefe, Wasser, wüste und leere Erde, ist nichts Anderes (wie es Petrus [2. Ep. 3, 4. 5.] erklärt), denn daß die Erde sei durch Gottes Wort bestanden im Wasser. Um die Erde herum ist es gewesen wie eine finstere Nacht oder Nebel, und feuchte Luft oder Gewässer, das gleichsam geraucht hat, und kein Licht da ist gewesen. Also ist die Welt angegangen.

6. Das soll nun der allereinfältigste Verstand sein, wie die Worte lauten, daß Gott da geschaffen habe Himmel und Erde, aber nicht also, wie es jetzund ist. Denn, daß man die Erde nicht also soll nehmen, bringt der Artikel des Glaubens, von dem man müßte also sagen, daß sie gezieret wäre gewesen mit Baum, Gras und allerlei Gewächse. Darum erklärt er sich selber, daß weder Gewächse noch lebendige Thiere darauf waren, sondern wüst und leer war, und nichts trug. Darnach aber, am dritten Tage und hernach, richtete er sie zu, daß sie allerlei trug. Darüber war sie bedeckt mit der Tiefe, welche er hernach heißt Wasser; das ist eine feuchte, dicke Luft gewesen, gleich als ob es eine Wolke wäre; in dem lag die Erde. Das aber über der Erde war, dasselbige war der Himmel, und war die Tiefe, darinne die Erde lag, und damit sie umringt war; und wo die Tiefe war, da war noch kein Licht, sondern der Wind, oder Geist Gottes, schwebte auf dem Wasser.

7. In der hebräischen Sprache ist „Wind“ und „Geist“ gleich Ein Name, und magst es hier nehmen, wie du willst. Wenn es ein Wind heißt, so ist es das, daß die Luft untereinander her webet\footnote{Weimarische: webet.} auf der Tiefe, wie sie pflegt; willst du es aber einen Geist heißen, so magst du es auch thun. Denn ich weiß es nicht eben zu örtern; aber kein wäre es, daß es Geist hieße, so könnte man es also verstehen, daß Gott die Creatur, die er geschaffen hatte, unter sich genommen habe, wie eine Henne ein Ei unter sich nimmt und das Hühnlein ausbrütet. Doch, ich will es lieber also, daß es ein Wind heiße, lassen bleiben. Denn ich wollte gerne, daß die drei Personen in der Gottheit hier ordentlich nach einander angezeigt würden, zum ersten der Vater; zum andern der Sohn; darnach der Heilige Geist, davon wir darnach sagen wollen. Also haben wir nun, wie zum ersten die Welt geschaffen ist, und wie es Gott angegriffen hat. Nun folgt nach einander, wie er alle Dinge unterschieden hat.

\subsection{Der erste Tag}

\begin{quote}
V. 3—5. Und Gott sprach: Es werde Licht; und es ward Licht. Und Gott sahe das Licht für gut an. Da scheidete Gott das Licht vom Finsterniß, und nannte das Licht Tag, und die Finsterniß Nacht. Da ward aus Abend und Morgen der erste Tag.
\end{quote}

8. Hier hebt sich eine große Frage: ob man das, das vorher, vor dieser Zeit geht, auch rechnen solle zum ersten Tage, oder ob man das Vorige nur als eine Vorrede verstehen solle? Es klingt also, als habe der erste Tag angefangen an dem, da Gott sprach: „Es werde Licht.“ Wir\footnote{In den Ausgaben: „Und wir“. Wir haben das „Und“ getilgt.} haben vor [§ 2] gesagt, daß nicht hat etwas können sein vor dem Anfang; so muß\footnote{„muß“ von uns gesetzt nach der Weimarischen Ausgabe. In den andern Ausgaben: müßte.} je der Anfang angefangen haben an der Nacht, und der Morgen angegangen sein, da Gott das Licht schuf. Denn er hat gesagt, daß es finster ist gewesen. Also ist die Nacht der erste Theil des Tages, welches die Schrift auch sonst hält. Denn man rechnet nach der jüdischen Sprache also, daß der Tag an der Nacht angeht, wenn es finster ist worden, und währt bis wieder auf den Abend; welches wir nun umkehren. Darum, dieweil wir klar haben aus dem Text, daß am ersten kein Licht ist da gewesen, sondern
\quellseite{25}
eitel Nacht und Finsterniß, welche bis auf den Morgen gewährt hat, und GOTT die Nacht vom Tage hat geschieden, muß man es also rechnen, wie der Text sagt, daß aus Abend und Morgen, das ist, aus Nacht und Tag, Ein Tag worden ist.

9.\footnote{Dieser ganze Paragraph fehlt in der Weimarischen.} Das sind nun die ersten Worte Mosis, damit er dies Buch ansehet; die soll man je, wie [§ 2] gesagt ist, auf das einfältigste verstehen, wie sie gesetzt sind, daß es sechs Tage unterschiedlich nach einander gewesen sind, und daß Moses am ersten Himmel und Erde mit Namen beschrieben hat, und dieselbigen noch taub und rohe gewesen sind und nicht geschmückt, wie hernach, wie er denn sich selbst erklärt, als er spricht: „Die Erde war wüste und leer“; also, daß noch keine Kraft noch Saft darinnen gewesen ist, daß etwas hätte mögen heraus wachsen, sondern im Wasser ist gelegen, welches um sie hergangen ist; daß er auch hier nennt die Tiefe, das ist, ein finsterer, dicker Nebel oder Luft, daraus das Wasser ist worden.

10. Weiter sind in diesen Worten auch berührt und heimlich angezeigt die drei Personen in der Gottheit; sonderlich aber sind die ersten zwei aufs klärlichste ausgedrückt, darob sich zwar viel hoher Leute, aber St. Augustinus am meisten, gebrochen haben. Dies ist aber die Summa davon: Aufs erste, da Moses spricht: „Am Anfang schuf GOTT Himmel und Erde“ etc., ist noch keine Person sonderlich genannt oder ausgedrückt. Aber so bald sagt er weiter: Und GOTT sprach: „Es werde Licht“; [da] drückt er aus, daß bei GOTT ein Wort war, ehe denn das Licht ist worden. Nun konnte dasselbe Wort, das GOTT da redet, nicht der Dinge etwas sein, die da geschaffen wurden, weder Himmel noch Erde, sintemal GOTT eben durch das Sprechen, das er that, Himmel und Erde sammt dem Licht und allen andern Creaturen machte, also, daß er nichts mehr zum Schaffen gethan hat, denn sein Wort. Darum muß es vor allen Creaturen gewesen sein. Ist es denn zuvor gewesen, ehe sich Zeit und die Creaturen anfangen, so muß es ewig sein, und ein ander und höher Wesen denn alle Creaturen; daraus denn folgt, daß es GOTT sei.

11. Wenn aber GOTT redet, und das Wort gefället, so ist er nicht allein; so kann er auch das Wort nicht selbst persönlich sein, das er redet. Darum, weil das Wort auch GOTT ist, muß es eine andere Person sein. Also sind die zwei Personen ausgedrückt: der Vater, der das Wort spricht und das Wesen von ihm selbst hat; der Sohn, der das Wort ist und vom Vater kommt, und ewig bei ihm ist.\footnote{Hier findet sich in der Weimarischen Ausgabe eine weitere Ausführung (von zwei Paragraphen), in welcher im Wesentlichen dieselben Gedanken gegeben sind, denen wir in § 14 begegnen.} Daher hat nun St. Johannes sein Evangelium genommen, welches er also anfäh(e)t: „Im Anfang war das Wort, und das Wort war bei GOTT, und GOTT war das Wort“ etc.

12. Nun muß hier ein fleischlicher Mensch seine Schuhe abziehen. Denn diese hohen, trefflichen Worte sind nicht Kindern gesetzt noch geschrieben, sondern wollen haben verständige Leute in der Schrift, die ihrer gewohnt sind; sonst lautet es kindisch, daß der Herr ein Wort spricht, und gleich durch dasselbige alle Creaturen schafft. So soll man nun halten, wie der Text zwingt, daß es ein solches Wort sei, das vor dem Licht und allen Creaturen gewesen sei, dieweil dadurch alles gemacht wird. Das Wort aber wird nicht gemacht, denn du hörst hier nicht, daß er das Wort mache, sondern spricht es nur. Damit will Moses lehren, daß es unaussprechlich und unergründlich sei, wie GOTT die Creaturen schafft. Wir können wohl sagen, daß GOTT alle Dinge hat geschaffen frei, aus Nichts; es ist uns aber unbegreiflich, ohne daß wir die Worte also können nachsprechen.

13. Also viel will nun Moses hier sagen von dem Licht: GOTT hat bei sich ein Wort, und durch dasselbige machte er am ersten das Licht, und bedurfte nicht mehr dazu denn das Wort. Damit beschloß er, daß das Licht sollte am ersten Tage gemacht werden; und weil er es also beschlossen hatte, geschah es auch also. Darum schließt sich’s, daß das Wort GOTT gleich sei, und muß GOTT selbst sein, dieweil es keine Creatur ist, sondern das, dadurch alle Creaturen entspringen. Wenn aber der, der da spricht, und das, das da gesprochen wird, zweierlei sein muß, so kann es nicht beides Eine Person sein, und muß dennoch Ein GOTT sein. Also wird hieraus gegründet der Artikel des Glaubens, daß zwei Personen in der Gottheit sind, doch Ein einiger\footnote{So das Original, die Wittenberger und die Jenaer. Die Erlanger hat, ebenso wie Walch: ein einiger. Letztere Lesart findet sich auch in der Weimarischen.} GOTT sei. Denn Moses hat mit
\quellseite{26}
diesen kurzen Worten sehr viel gegeben, und redet also, daß es unmöglich ist einem Menschen also zu reden, und mit solchen kurzen Worten so viel zu sagen, wenn er schon den Verstand hätte. Es ist so hoch, so eigentlich, kurz und tief geredet, daß es alle Vernunft übertrifft.

14. Die dritte Person aber, nämlich der Heilige Geist, wird heimlich in den Worten angezeigt und begriffen, als Moses spricht: „Gott sähe das Licht für gut an“, das ist, er hatte ein Wohlgefallen daran. Denn es ist ein Wohlgefallen an dem gewesen, das er durch das Wort gemacht hatte; welches Wohlgefallen nichts Anderes ist, denn daß Gott die Creaturen erhält,\footnote{Weimarische: „erhebt“ und gleich folgend: „erhebung“ statt: Erhaltung.} wie er sie gemacht hat, und ihnen beisteht. Das wird eigentlich dem Heiligen Geist zugeeignet, daß er sei das Leben und Erhaltung aller Dinge. Darum die Schrift also von ihm redet, daß er das Band sei, das da alle Creaturen halte, und allen ihre Uebung und Wirkung gebe. So heißt nun der Heilige Geist das Wohlgefallen Gottes des Vaters, wie das Wort sein ewiger Rath ist.

15. Hier ist aber\footnote{aber = wiederum.} eine Frage: warum Moses nicht vorne an also gesetzt habe: Am Anfang sprach Gott: Es werde Himmel und Erde, sondern spricht also: „Am Anfang schuf Gott Himmel und Erde“, und allererst hernach spricht: „Gott sprach: Es werde Licht.“ Antwort, aufs erste: Ob wir die Ursache nicht treffen, wollen wir dem Heiligen Geist die Ehre thun, daß er es besser wisse denn wir. Aber, so viel ich kann denken, ist's der Ursache halben also gesetzt. Denn, wenn er gesagt hätte: Am Anfang sprach Gott: Es werde Himmel und Erde, würde es fast also klingen, daß das Wort nicht wäre vor dem Anfang gewesen, daß man nicht könnte wissen, ob es angefangen hätte, oder von Ewigkeit gewesen wäre. Darum konnte nicht wohl dabei stehen, daß er hätte gesagt: Am Anfang sprach Gott; daß nicht jemand meinte, es wäre allererst angegangen, da die\footnote{„die“ fehlt in der Erlanger.} Zeit und die Creaturen angingen, wie die Ketzer Arriani gesagt haben. Also ist dem Irrthum zuvor damit begegnet. Das ist Eine Ursache, so viel ich kann ermessen.

16. Die andere ist, daß er die drei Personen hat müssen ordentlich nach einander anzeigen. Zum ersten den Vater, da er sagt: „Gott schuf“; zum andern den Sohn, da er sagt: „Gott sprach“; darnach den Heiligen Geist, da er sprach: „Gott sahe es für gut an.“ Das Sehen, das Gott gethan hat, ist nicht nach dem Sprechen, sondern zugleich; und sein Sprechen und Sehen ist nicht da angegangen, sondern sein Schöpfen ist da angegangen. Wiewohl nun die Dreie zugleich gewesen sind, so mußte er es dennoch nach einander schreiben, denn er konnte es nicht zugleich mit Einem Worte fassen. Wie nun die drei Personen nach einander gehen, und wir je die Worte auch nach einander also lesen müssen, wenn wir davon wollen reden, und nicht auf einmal, mit Einem Worte, aussprechen können, also hat er es auch nicht zugleich können sagen: Gott schuf, sprach, und sahe; und will damit nicht verleugnet, sondern bekräftigt haben, daß alle drei Personen gleich seien, und alle auf den ersten Tag da gewesen.

17. Denn dieweil er sagt, daß Gott durch das Wort das Licht macht, und scheidet von der Finsterniß, so ist gewißlich das Wort vor dem Licht und Finsterniß gewesen. Denn wie könnte er es durchs Wort machen und von einander scheiden, wenn es nicht vorhin da wäre? Also geht es in einander, daß der Vater das Werk thut durch den Sohn, der das Wort ist, und in dem Heiligen Geist, der das göttliche Wohlgefallen ist, und ist jeglicher Person ihre Eigenschaft gegeben, daß man kann einen Unterschied fassen, daß drei Personen und Ein Gott bleiben; und wiewohl sie zugleich sind, haben sie dennoch ihre unterschiedliche Beschreibung müssen haben.

18. Was er nun nicht zum erstenmal konnte thun, daß er das Wort vorne an setze, das kehrt er darnach wieder um, und setzt das Wort vorne an, als er sagt: „Gott sprach: Es werde eine Feste“, darauf denn folgt: „Da machte Gott die Feste; und es geschah also.“ Solche Veränderung hat er eigentlich darum gethan, daß er wollte wehren, daß man das Wort nicht nach, sondern immer vor den Creaturen stellte; wiewohl er es am ersten Tage, aus Ursache, nicht hat vorne an können setzen, wie er doch darnach thut. Und ist alles mit Fleiß und bedachtem Rath also geschehen, daß der Artikel des Glaubens da recht gegründet würde, daß das Wort zugleich ist da gewesen, da Gott alle Dinge schuf. Also haben wir den ersten Tag, da die Erde sammt dem Himmel rauh und roh ist gewesen,
\quellseite{27}
und das Licht geschaffen; welches er nennt den Tag, als die Nacht zuvor her gegangen ist. Was nun dieser Tag und diese Finsterniß seinen gewiesen, werden wir weiter hören am vierten Tage, da wird sich’s wider einander stoßen. Nun folgt

\subsection{Der andere Tag}

B. 6—8. Und Gott sprach: Es werde eine Feste zwischen den Wassern; und die sei ein Unterschied zwischen den Wassern. Da machte Gott die Feste, und scheidete das Wasser unter der Feste von dem Wasser über der Feste. Und es geschah also. Und Gott nannte die Feste Himmel. Da ward aus Abend und Morgen der andere Tag.

19. Hier soll man immerdar Achtung haben, daß man im einfältigen Verstande bleibe. Darum, wie er am ersten Tage macht eine wüste und leere Erde, und den Himmel mit einem finstern Nebel oder finstern Gewässer: also nimmt er sonderlich hier am andern Tage den Himmel vor sich, und richtet eine Feste zu; das ist, das rohe, finstere Gewässer, das vor nichts denn ein schwebend und webend Ding war, machte er hier fest, und machte es also, daß er die Feste setzte zwischen zweien Gewässern. Denn vorhin waren eitel Wasser rings um die Erde herum. Jetztund aber greift er mitten drein, und scheidet die Wasser in zwei Theile, behält eines oben, das andere unten. Mitten hinein macht er einen Ring, den wir den Himmel heißen, und nennt ihn eine Feste darum, daß er nicht so wankt, hin und her fährt und unstät ist, wie die Luft und Wasser, und wie er auch am Anfang war, sondern steht fest, und geht in strackem Lauf, und hat doch unter sich die Luft, Wasser und das Meer. Zuvor aber, ehe er also zugerichtet ward, ging er auch unstät.

20. Was aber für Wasser droben über der Feste sei, können wir nicht wohl wissen. Darum müssen wir, wie ich [§ 15] gesagt habe, dem Heiligen Geist Raum geben, und sagen, daß er es besser wisse, weder wir es verstehen. So kann Gott freilich wohl das Wasser über dem Himmel erhalten. Ich wollte wohl aus dem Wasser über der Feste die Luft machen, so bleibt es aber dennoch immerdar unter dem Himmel. Darum müssen wir uns gefangen geben, und dabei bleiben, daß der Himmel mitten zwischen den Wassern sei gemacht.

21. Da haben wir nun, wie er es zum ersten angreift, daß er dem\footnote{Wittenberger und Erlanger: den.} Himmel nimmt das unstäte Wesen, und gibt ihm ein sicher Wesen, daß er also recht wird zugerichtet und in die zwei Haufen getheilt, daß das Finsterniß und Schweben hinweg kommt, das ist, daß es nicht mehr finster ist auf der Tiefe, und der Wind nimmer schwebt auf dem Wasser. Und als er nun also zugerichtet ist, gibt ihm erst Gott an dem andern Tage seinen rechten Namen, und nennt es den Himmel. Am ersten Tage ist es auch wohl Himmel und Erde gewesen, aber es hatte noch nicht einen rechten Namen. Denn es war noch nicht recht vollkommen zugerichtet und befestigt, daß es also konnte heißen.

22. Also haben wir, meine ich, den aller einfältigsten und rechten Verstand des Texts. Doch wollen wir dem Heiligen Geist kein Ziel stecken. Gibt er uns ein Weiteres, so wollen wir es auch gern haben. So beschließt nun Moses, daß wieder aus Abend und Morgen ist worden der andere Tag. Also hat er über diesem Werk auch einen Tag und Nacht gemacht.

23. Aber hier entspringt eine Frage: Warum Moses nicht auch hier dazu setzt: „Gott sahe es für gut an“, und da den Segen hinweg nimmt, und nicht schreibt, daß Gott ein Wohlgefallen daran gehabt habe, so er doch sonst desgleichen gedenkt durch die andern Tage alle? Da hat man sich auch sehr bekümmert, was er damit gemeint habe, und mancherlei darüber erdacht. Ich halte aber, es bleibe wohl verborgen. Darum weiß ich auch keine gewisse Ursach darauf zu geben. Sie sagen, daß die zweite Zahl sei eine Theilung aus der Einigkeit, welche\footnote{Erlanger: „und“ statt: welche.} bedeute die Tiefe; wie denn dieselbe Zahl in der Schrift ein groß Sacrament und Heimlichkeit hat. Aber das wollen wir fahren lassen. Es sei, wie ihm wolle, so ist’s ohne Zweifel, daß es Moses nicht vergessen hat dazu zu setzen, sondern ist aus bedacht Rath geschehen. So ist nun der Himmel zubereitet, daß er nicht mehr finster ist, noch schlottert und wankt, sondern fest steht, und seinen Raum und Statt hat. Folgt

\subsection{Der dritte Tag}

B. 11—13. Und Gott sprach: Es lasse die Erde aufgehen Gras und Kraut, das sich besame; und fruchtbare Bäume, da ein jeglicher
\quellseite{28}
nach seiner Art Frucht trage, und habe seinen eigenen Samen bei ihm selbst auf Erden. Und es geschah also. Und die Erde ließ aufgehen Gras und Kraut, das sich besamet, ein jegliches nach seiner Art, und Bäume, die da Früchte trugen,\footnote{Erlanger: trügen.} und ihren eigenen Samen bei sich, selbst hatten, ein jeglicher nach seiner Art. Und Gott sahe es für gut an. Da ward aus Abend und Morgen der dritte Tag.

24. Nachdem in den zweien Tagen der Himmel zugerichtet ist, nimmt Gott hier am dritten Tage die Erde vor sich, und richtet sie auch recht zu. Denn zuvor war sie bedeckt und lag im Wasser, war auch niemand darauf, sondern ganz wüste und nichts nütze. Jetzund aber, als er das Wort spricht: Die Erde komme hervor, und die Wasser unter dem Himmel theilen sich, hier ins Meer, da in Teiche oder Seen, daß man das Land und Inseln sehe, also, daß hier ein Stück, und dort ein Stück von der Erde hervorkomme: da hört es nun auf, daß die Erde nimmer so wüst und ungeschickt ist, und kommt nun ans Licht und gewinnt eine rechte Gestalt, daß sie etwas Anderes ist, denn vor; so ist sie auch nimmer leer, sondern es wird ihr eine Kraft gegeben, die sie haben soll in ihr selbst, daß sie heraus gebe und von sich bringe Gras, Kraut und Bäume, und allerlei Gewächse. Das ist also geschehen, und gefiel Gott wohl, und gibt ihr alsdann einen rechten Namen. Darum wohl zu nehmen ist, daß er sie eben am ersten Tage nicht recht hat zugerichtet. Denn jetzund wird sie erst ihres Namens werth, weil sie nicht mehr wüst noch leer ist, sondern zugerichtet und gearbeitet, wie sie sein soll, und ihre Kraft kriegt, daß sie allerlei Gewächs von sich bringt.

25. Das sind nun alles treffliche Worte, damit Mosis schon das Maul gestopft hat allen Philosophis, und die bisher studirt haben in den hohen Schulen, und sich genannt Doctores und Magistros, und haben gelehrt die natürliche Kunst, ja, natürliche Thorheit. Die haben viel disputirt und mancherlei Wahn gehalten, wovon die Erde und alle Dinge hergekommen sind, und hat es doch keiner getroffen. Darum schlägt hier Moses solches alles mit kurzen Worten nieder, und sagt, daß durch das ewige Wort Gottes ist beschlossen gewesen, daß am dritten Tage die Erde sollte hervor kommen und zugerichtet werden.

26. Wenn wir nun alles gegen einander halten, so sehen wir, daß die Creaturen ihr Wesen nicht von ihnen selbst haben; und wenn ihnen gleich das Wesen gegeben ist, haben sie doch auch keine Kraft von ihnen selbst. Denn die Erde und der Himmel, weil sie wüst, leer und finster waren, konnten sie nichts von sich geben noch bringen, sondern waren ein arm, elend, bloß Wesen. Sollte es aber Licht werden, so mußte es Gott geben und schaffen, von ihnen selbst vermochten sie es nicht. Also war es auch nicht des Himmels Kraft noch Vermögen, daß er fest ward. Desgleichen hat auch die Erde deren keines von ihr selbst, sondern mußte alles von neuem in sie gepflanzt werden, daß sie Korn, Gras, Kraut und Bäume trüge; wie wir sehen, wenn man sie aufgräbt, daß sie hervorbringt allerlei Gewächs, auch daß man weder säet noch pflanzt. Es ist wohl ihr Vermögen; sie hat es aber von ihr selbst nicht, sondern Gott hat es ihr durch das Wort eingepflanzt, daß [es] niemand sehen noch begreifen kann, wie es zugeht, und doch zusehends sehen, wie es hervor wächst; wie wir auch nicht begreifen können, wie sie geschaffen sei. Darum müssen wir die Vernunft zuschließen, und nur glauben; wie Moses hier lehrt, daß es Gott also durchs Wort schaffe, daß es ohne Unterlaß wächst, und wir doch nicht wissen wie. Also muß man auch sagen von dem Himmel, woher er das Licht hat und fest steht, welcher er keines von ihm selbst hat. Denn Gott könnte wohl machen, daß er schwerer wäre, weder die Erde, wie die Wolken oben schweben, die doch schwerer sind weder die Luft. Daß es nichts ist, das Aristoteles sagt: Was leicht ist, das geht oben. Der Dinge kannst du keines verstehen, es sei denn, daß du glaubest.

27. Und hüte dich, wenn du diese Worte hörst: „Und Gott sprach“, daß du sie nicht denkest, als seien es vergängliche Worte, wie wir Menschen sprechen, sondern wisse, daß es ein ewig Wort sei, das von Ewigkeit gesprochen ist, und immer gesprochen wird. So wenig als Gottes Wesen aufhört, so wenig hört auch das Sprechen auf, ohne daß zeitlich die Creaturen durch dasselbige angefangen haben. Aber er spricht noch immerdar, und geht ohne Unterlaß im Schwange, denn keine Creatur vermag ihr Wesen von sich selbst zu haben. Darum, so lange eine Creatur währt,
\quellseite{29}
so lange währt das Wort auch; so lange die Erde trägt oder vermag zu tragen, so geht immer das Sprechen, ohne Aufhören. Also verstehe Moses, wie er uns Gott vorhält in allen Creaturen, und durch dieselbigen zu Gott führe. So bald wir die Creaturen ansehen, daß wir denken: Siehe, da ist Gott also, daß alle Creaturen in ihrem Wesen und Werken ohne Unterlaß getrieben und gehandhabt werden durch das Wort.

28. Denn es ist auch wohl daran zu prüfen, daß die Erde allein ihre Kraft von Gottes Wort habe, daß man noch eine Erde sieht, darauf nichts wächst; noch ist es Erde und trocken Land, sie ist aber leer, denn Gott gibt dahin sein Wort oder Gebot nicht, daß sie trage und etwas da wachsen könnte. Darum, daß alles Land nicht gleich Früchte trägt, ist daher, daß solch Vermögen nicht ist vom Lande, sondern von Gottes Wort. Wo das ist, da folget auch die Kraft, fruchtbar zu sein, also, daß die ganze Welt voll Worts ist, das alle Dinge treibt, Kraft gibt und erhält. Darum sagt Gott hernach [Cap. 3, 17. 18.] zu Adam: „Die Erde soll dir ihr Vermögen nicht geben“, da er einen Fluch über sie gibt; zeucht das Wort ab, und läßt sie in ihrer Natur. Also haben wir die Erde, zugerichtet, daß sie nimmer wüst und leer ist, sondern voll Saft und Kraft und Vermögens, nicht von ihrer Natur, sondern durch Gottes Eingeben und Sprechen.

29. Daneben mußt du auch der Sprache gewohnen, daß der Text sagt: „Die Sammlung der Wasser hieß er Meer.“ Denn die hebräische Sprache heißt alles Meer, wo Wasser ist, nämlich auch Seen und Teiche; darum sind viel Meere in der Schrift hin und wieder; wo aber trocken Land ist, heißt sie alles Erde.

30. So ist nun die Erde ausgeschmückt, und das Wasser davon geschieden, welches vorhin rings um sie ging, daß sie mitten darinnen schwebte. Da siehe, wie leicht es Gott zu thun wäre, daß er die Welt ersäufte, denn es ihre Natur ist, daß das Wasser um und über sie geht; aber durch Gottes Wort und Befehl wird sie erhalten wider ihre Natur, sonst bliebe sie wohl ungehalten, und ginge alle auf einen Klumpen. Also sagt er im Hiob [Cap. 38, 8—11.]: „Wer hat das Meer mit seinen Thüren verschlossen, da es heraus brach, wie aus Mutterleibe? Da ich es mit Wolken kleidete, und in Dunkel einwickelte, wie in Windeln. Da ich ihm den Lauf brach\footnote{Wittenberger und Erlanger: bracht.} mit meinem Damm, und setzte ihm Riegel und Thür, und sprach: Bis hieher sollst du kommen und nicht weiter. Hier sollen sich legen deine stolzen Wellen.“ Wenn das Wasser herschlägt, so schützet niemand, daß das Meer nicht alles ersäuft, denn Gottes Wort, das\footnote{So die Jenaer. Erlanger: „daß es“. Wittenberger: welches.} alles gewaltig in der Hand hat.

31. Also malt uns die heilige Schrift Gott und die Creaturen ab, daß sie die Vernunft niederstoße, die da meint, er sei von den Creaturen hinweg, und habe nichts damit zu schaffen. Es ist nicht ein Stüblein noch Tröpflein, damit Gott nicht zu schaffen habe, und daselbig treibe.

32. Aber hier ist uns vorzusehen, daß wir nicht anlaufen. Denn es auch gefährlich ist, wenn eine unreine Seele drauf fällt und fähret an, darnach hoch zu trachten, daß sie nicht verzage. Darum gehört das Stück nicht für Kinder. Denjenigen, so das schmecken, das Moses sagt: „Gott habe es für gut an“ etc., denen ist es\footnote{„es“ fehlt in der Erlanger.} nöthig, darnach zu trachten und dichten. Die sind es, die Gott erkennen in den Creaturen, daß sie sagen müssen: Ei, wie ein freundlicher, feiner Gott ist das? ist er doch nichts denn eitel Süßigkeit und Güte, daß er uns also speist, erhält und ernährt. Diese sehen nichts Anderes an den Creaturen, denn das Gott daran sieht. Er sieht aber nichts daran, denn Gutes. Darum haben sie alle Lust und Freude davon. Die andern aber, so da böse Gewissen haben und nicht glauben, sollen wohl vor dem geringsten rauschenden Blatt erschrecken, wenn sie fühlen, daß Gott so nahe ist, und mögen keine Creatur leiden, wie gut sie immer ist. Darum ist es ein hoch und groß Ding, wer dahin zu diesem Kern kommt, und das versteht. Davon auch oben [§ 14 in der Vorrede] gesagt ist, und habe es darum hier wiederholt, daß man es allemal brauche, und sich erinnere durch dies ganze Capitel.

33. So hat er nun in den dreien Tagen Himmel und Erde zum ersten geschmückt und geziert; den Himmel mit der Feste und Licht, die Erde mit Gras und Kraut, und allerlei Holz oder Bäumen. Wird nun folgen, wie Gott in den andern dreien Tagen beide, Himmel und Erde,
\quellseite{30}
noch mehr, weiter, besser schmückt und reich macht. Und hält gleich die Ordnung, wie in den vorigen Tagen, daß er zum ersten\footnote{In den alten Ausgaben: von ersten.} den Himmel schmückt mit Lichtern und Sternen; darnach die Erde, daß er hineinsetzt allerlei Thiere, Vögel, Fische und Menschen.

\subsection{Der vierte Tag}

B. 14—19. Und Gott sprach: Es werden Lichter an der Feste des Himmels, und scheiden Tag und Nacht, und seien zu Zeichen, Zeitungen, Tagen und Jahren, und seien Lichter an der Feste des Himmels, daß sie scheinen auf Erden. Und es geschah also. Und Gott machte zwei große Lichter, ein groß Licht, daß dem Tag vorstände, und ein klein Licht, daß der Nacht vorstände, dazu auch Sterne. Und Gott setzte sie an die Feste des Himmels, daß sie scheinen auf die Erde, und dem Tage und der Nacht vorständen, und scheideten Licht und Finsterniß. Und Gott sahe es für gut an. Da ward aus Abend und Morgen der vierte Tag.

34. Das ist das Werk, am vierten Tage gemacht, daß Sonne und Mond sammt allen Sternen, klein und groß, geschaffen und an Himmel geheftet sind, daß sie vorständen dem Tage und Nacht und die Erde erleuchteten, und wir daran hätten Stunden, Tage und Jahrzeiten. Aber das soll man nicht also verstehen, daß Gott zuvor Sonne und Mond habe geschaffen, und hernach an Himmel geheftet, wie man eine Spange anheftet, sondern weil unsere Sprache also gespannt ist, daß sie ein Wort nach dem andern machen muß, und zwei Dinge, die zugleich geschehen sind, nicht zugleich mit Einem Worte ausreden und fassen mag, so hat er vor müssen sagen, wie sie geschaffen sind, und hernach, wie sie an Himmel geheftet sind, wiewohl es beides zugleich geschehen ist, daß, indem es gemacht ward, auch an Himmel ist gesetzt worden. Also, daß man wisse, daß Gott einen Tag daran gemacht habe; wie lang oder kurz, viel oder wenig Stunden, steht bei ihm. Steht wohl darauf, daß er langsam mit ungezwungenen sei; nicht darum, als könnte er es nicht in einem Augenblick ausrichten, sondern hat auch einen Tag wollen dazu nehmen, wie er oben an dem dritten Tage alles Gras, Kraut und Bäume, alle Haiden, Wiesen und Feld gemacht hat. Und hat wahrlich viel gemacht, ob er gleich nicht eine Stunde daran geschaffen hat.

35. Also mag er auch hier von Abend bis zu Morgen einen ganzen Tag an dem Gestirne gemacht haben, deß so viel ist, daß sie niemand zählen kann, als die Schrift [1 Mos. 15, 5.] sagt. Gleich als noch jetzund: wenn er die Sonne aufgehen läßt, treibt er sie also, daß sie von Morgen an bis auf den Abend läuft, so sie doch wohl in einer Stunde vom Morgen bis zum Abend gehen könnte, ja, in einem Hui an beiden Orten sein. Er thut es aber nicht, sondern gibt Raum und Weile dazu. Wie er nun hier mit der Sonne thut, als wir vor Augen sehen, also thut er mit allen Creaturen. So läßt er im Lenzen Gras und Kraut aufgehen, daß es alles grünt und blüht, so lange es Sommer ist; welches alles wohl in einer Stunde geschehen möchte, oder, so er wollte, in einem Jahr oder zweien. Er macht es nach seinem Willen und Wohlgefallen so lang, und nicht kürzer noch länger.

36. Das sage ich alles darum, daß man bei den einfältigen Worten und Verstand bleibe, und nicht empor fahre und hoch speculire. Wenn wir gleich daran fehlen, so ist es doch ohne Fahr. Das sollen wir aber hieraus lernen, wie ich auch oben [§ 26] gesagt habe, wie keine Creatur weder ihr Wesen noch Kraft von ihr selbst habe, und nicht in ihrer Macht ist, wie lange sie währen und bleiben soll, sondern hat alles seine Ordnung von Gott, wie lange er will, daß es währen soll, daß wir je sehen, wie Gott die ganze Welt in seiner Hand gewaltiglich hält, daß sich nichts regen kann, denn was und wann er will. Das sei nun das erste Stück.

37. Weiter sagt Moses, Gott habe Sonne und Mond dazu gemacht und geordnet, daß sie vorstehen sollten dem Tag und der Nacht. Die Sonne ist ein Herr des Tages, der Mond ist ein Meister der Nacht. Der Mond macht nicht die Nacht, sondern sein Licht geht und leuchtet dennoch in der Nacht. Also auch macht die Sonne den Tag nicht, sie geht aber und scheint am Tage. Denn wir haben zuvor gehört, daß Gott das Licht schuf, und nannte es den Tag; die Finsterniß aber nannte er Nacht. Da war noch keine Sonne noch Mond, und war doch Tag und Nacht, und wäre auch also blieben. Also, wiewohl die Sonne den Tag nicht macht und nichts dazu thut, regiert sie doch am Tage, wie der Mond in der Nacht.
\quellseite{31}

38. Derhalben müssen wir den Text aufs einfältigste lassen bleiben, daß die zwei Lichter darum geordnet sind, daß sie Tag und Nacht regieren, nicht, daß sie Tag oder Nacht machen sollen. Wie wir von einem Menschen sagen, er soll über das Land regieren; daß nicht so viel gesagt ist, daß er das Land mache, sondern daß er es findet, und sein Regiment darauf ist. Wenn gleich Sonne und Mond nicht geschaffen wären, so wäre gleichwohl Tag und Nacht, Licht und Finsterniß gewesen; es wäre aber kein Regiment gewesen. Also liest man auch im Psalter hin und wieder, als [Ps. 136, 7—9.]: „Er hat große Lichter gemacht, die Sonne, den Tag vorzustehen; den Mond und Sterne, der Nacht vorzustehen.“

39. Also sehen wir, daß es allenthalben geht aus Gottes Ordnung, und kein Ding sein Wesen noch Regiment von ihm selbst hat, sondern alles von Gottes Hand, Rath und Willen, daß man in allen Creaturen Gott sehen soll, wir thun die Augen oder Ohren auf, und ihm also danken. Und, als ich oben [Vorrede, §§ 11. 12] genug gesagt habe, sollen wir nicht denken, daß Gott die Creaturen alsohin geschaffen habe, und sei hernach davon gegangen, und lasse sie forthin selbst handeln; sondern daß er sie gemacht habe, und noch immer mache, und erhalte sie durch das Wort. Also, das Wort treibt die Sonne noch heute und immerdar vom Morgen bis zum Abend. Denn daß sie am Himmel alle Tage einmal herum läuft, ist nicht ihre eigene Natur oder Kraft, sondern des Wortes, durch welches sie geschaffen ist.

40. Item, daneben spricht Gott, daß diese zwei Lichter dazu dienen sollen, daß sie seien zu Zeichen, Zeitungen, Jahren und Tagen, das ist, daß man dabei die Zeit könnte zählen und wissen. Das zu thun, haben wir aus unserer Kunst noch Gedanken nicht, sondern ist alles also von Gott geordnet. Also mißt man nun bei diesen Lichtern beide Tag und Nacht, daß man der Sonne Schein in zwölf Stunden theilt, und lehrt auch das Jahr bei der Sonne Lauf rechnen, weil sie alle Jahre einmal herum kommt, und mit ihrem Lauf die Jahrzeit, Sommer und Winter, mit sich bringt.

41. Sie sollen auch dazu dienen, daß man die Zeitungen dabei wisse zu rechnen. Das Wort Zeitungen heißt auf Hebräisch מועד. Also nennt die Schrift die sonderlichen Feste im Jahre, und ist gleich wie man einen Tag theilt über die zwölf Stunden in vier Theile, Morgen oder Vormittag, Mittag, Nachmittag und Abend; das in unserer Willkür steht, wie wir es theilen wollen. Wie man nun den Tag theilt, also kann man auch das Jahr theilen, in wie viel Stücke man will. Als, zum ersten in vier Theile; darnach in zwölf Monat, dabei man eine gewisse Zeit kann bestimmen; als, wenn ich sage, im vierten oder fünften Monat. Item, weiter kann man einen jeglichen Monat theilen ins halbe, oder vierte Theil; das heißt Hebräisch: מועד, festa vel tempora, das ist, solche Zeitungen, die im Jahre bestimmt werden. Als nach unserer Rechnung sind Weihnachten, Ostern, Pfingsten, oder Michaelis etc. Wie man nun das Jahr mag scheiden und theilen, das heißen Zeitungen, dazu sind Sonne und Mond geschaffen.

42. Weiter, zum ersten spricht Gott, sie sollen „Zeichen“ sein. Da sind die Sternkucker und natürliche Meister hinauf in Himmel gefahren, und haben das, daß er hie von Zeichen sagt, auf ihre Lügen gezogen, daß sie sagen: Wer in dem oder diesem Zeichen der Gestirne geboren wird, der soll so oder also geschickt werden. Welcher unter der Sonne geboren wird, der müsse ein Buhler oder weiser Mann werden. Wer in Mercurio geboren wird, der werde ein guter Hantierer werden, und so fortan werde es ihm sonst oder so gehen. Aber diese groben Lügen lassen wir fahren, und bleiben bei dem einfältigen Verstande, daß sie Zeichen sind, als wie ihrer brauchen die Schiffsleute, und sich darnach richten auf dem Meer. Item, wie ein Zeichen ist, wenn die Sonne oder der Mond ihren Schein verliert, damit Gott etwas läßt anzeigen, nämlich ein zukünftiges Unglück und Strafe über die Welt.

43. Darum sagen wir, daß sie nur Zeichen sind, und nicht sonderliche Kraft oder Wirkung haben. Daß sie nun sagen, daß die Lichter und Gestirne dreierlei wirken auf Erden: zum ersten, daß sie leuchten, zum andern eine Bewegung, zum dritten einen Einfluß haben, ist eitel Narrenwerk. Daß das Licht herab wirkt, das greift man wohl, wenn die Sonne scheint. So ist es auch wahr mit der Bewegung; denn wir sehen, je weiter sie hinweg kommt, je kälter es wird. Aber das Dritte, daß sie sagen, daß ein jegliches Zeichen und Gestirn seinen Einfluß habe, sonderlich auf die Menschen, daß, wer unter einem
\quellseite{32}
solchen Zeichen geboren wird, der müsse also genaturt werden, so ein Leben führen, eines solchen Todes sterben, damit sie alles wollen rathen, wie es einem jeglichen gehen soll, das ist falsch und erdichtet. Denn sie sind nicht darum geschaffen, daß sie mich meistern, sondern mir zu Nutz und zu Dienst. Ueber Tag und Nacht sollen sie regieren, aber über meine Seele sollen sie kein Regiment noch Gewalt haben. Der Himmel ist dazu gemacht, daß er Licht und Zeit gebe; die Erde, daß sie uns trage und speise. Mehr können sie nicht von sich geben noch wirken.

44. Ueber das wollen die Narren auch ausfechten, warum ein jeglich Land etwas Sonderliches trägt, als, daß dieses Gold und Silber trägt, ein anderes Korn etc., und lügen mit Gewalt von solchen Dingen, wie sie wollen, wie die alten Leute und Landfahrer\footnote{Alte Leute lügen von dem, was der Zeit nach, Landfahrer von dem, was dem Raume nach fern liegt. Landfahrer (Landstreter) = Reisende, die weit über Land kommen (vgl. Walch, St. Louiser Ausg., Bd. XI, 301, § 18); auch zu Wasser (Walch, alte Ausg., Bd. III, 961, § 34).}, daß sie keinen Grund noch Ursache haben. Wenn sie aber uns fragen: warum an einem Orte Korn, am andern Erz, Safran, Ingwer und ander Ding wächse, das man sonst nicht an mehr Orten findet? so wollen wir antworten: Also stehet geschrieben in diesem ersten Capitel, V. 12., von dem dritten Tage, daß Gott an demselben hat von der Erde, durch sein Wort, lassen aufgehen\footnote{Wittenberger und Erlanger: aufgehen.} Gras, Kraut, Bäume, und allerlei Gewächse, ein jegliches nach seiner Art.

45. Darum ist der keines, was da aus der Erde wächst, von des Himmels Einfluß, sondern durch Gottes Wort also geschaffen. Und ob sie gleich das nicht wollen zulassen, so müssen sie dennoch bekennen, daß alles Gras, Kraut und Bäume ehe geschaffen ist denn Sonne und Mond mit den Sternen. Und ist freilich auch darum geschehen und also geschrieben, daß man den Narren nicht glaube, als komme es von des Himmels Einfluß, sondern sei eine sonderliche göttliche Ordnung. Da bleibe bei, und behalte deinen Glauben rein. Sonne und Mond soll nicht mehr schaffen noch Kraft haben, denn hier geschrieben ist. Hätte ihnen Gott mehr gegeben, so hätte er mehr lassen beschreiben, sintemal in diesem Capitel alle Creaturen, fast mit allen ihren Werken und Vermögen, begriffen sind. Das ist der vierte Tag. Spricht nun Mosis weiter:

\subsection{Der fünfte Tag}

\begin{quote}
V. 20—23. Und Gott sprach: Es errege sich das Wasser mit webenden und lebendigen Thieren, und mit Gewögel, das auf Erden unter der Feste des Himmels fliegt. Und Gott schuf große Wallfische und allerlei Thier, das da lebt und webt und das Wasser erregt, ein jegliches nach seiner Art, und allerlei gefiedertes Gewögel, ein jegliches nach seiner Art. Und Gott sahe es für gut an, und segnete sie, und sprach: Seid fruchtbar und mehret euch, und erfüllet das Wasser im Meer,\footnote{Erlanger: der Meere.} und das Gewögel mehre sich auf Erden. Da ward aus Abend und Morgen der fünfte Tag.
\end{quote}

46. Da nimmt er abermal den Himmel vor, und gibt ihm noch ein Geschmuck. Bisher hat er ihm gegeben Licht, Feste, Sonne, Mond und Sterne, daß er alles hat, was er haben soll. Hier gibt er ihm zu die Vögel, welcher Reich in der Luft ist. Und hier siehst du, daß die Schrift Himmel heißt alles, was oben ist, auch die Luft, darinne die Vögel schweben; item die Wolken; also, daß Wasser und Himmel, oder Luft, fast Ein Ding ist in der Schrift, denn der Himmel aus dem Wasser kommen ist. Und wir sehen, daß auch etliche Vögel im Wasser leben. Darum faßt er auch hier zusammen, beide die Vögel am Himmel und Fische im Wasser, und hebt wieder am Wasser an.

47. Hier merke abermal, daß keine Creatur von ihr selbst lebendig sein kann; daß auch kein Vogel, der doch klein ist, Fittige oder Federn von ihm selbst könnte haben, sondern Gott gibt es einem jeglichen, also, daß es unmöglich ist, die Fittige zu zählen, und nachzusehen, wie es gemacht ist. Daneben hörst du, daß beide, die Vögel und Fische, aus wässeriger Natur gemacht sind, wie auch die Sonne und Mond, das ist, aus dem Himmel, der von Wasser gemacht ist, wie das Gras und Bäume aus der Erde. Denn die zwei Creaturen haben die andern alle gegeben.

48. Da geht nun das erstmal an, daß Gott der Herr den Segen gibt, welches er bisher nicht gethan hat. Denn der Segen geht nur über das, so da lebt, dazu, daß es sich mehren oder fruchtbar sein soll. Denn er hat allen die Natur und Geschicklichkeit in ihrem Leibe gegeben, daß sie können Früchte zeugen und viel tragen; das noch immer anzeigt, daß ich [§ 26] gesagt habe,
\quellseite{33}
daß keines solches von ihm selbst habe. Denn der Text sagt zum ersten, daß sie GOtt habe heißen aus dem Wasser kommen, und sich erregen, und sie also geschaffen, und ihnen das Leben gegeben; aber fruchtbar konnten sie noch nicht sein, bis er weiter das Wort über sie spricht und damit er sie segnet, und heißt sie, sich mehren. Das ist nun eine eingepflanzte Natur, wie er oben [V. 11. 12.] in die Erde gepflanzt hat die Kraft, das Gras und Bäume zu tragen. Wo er sie nur hervorgebracht hätte, würden sie die Kraft nicht gehabt haben, sich zu mehren. Darum ist es ihre Natur nicht, daß sie Junge zeugen, sondern ein sonderlich Werk göttlicher Majestät und seines Wortes oder Segens, in die Natur gepflanzt.

49. Von den Wallfischen weiß ich nichts Sonderliches zu sagen, ohne daß dabei zu wissen ist, daß er auch ein Herr des Meers ist, und darinne viel mächtige, wunderliche und ungeheure Thiere macht. Ich halte auch, daß sich's viel mehr mehrt, und eine größere Menge von Fischen im Wasser ist, denn anderer Thiere auf dem Lande, und freilich das Wasser das reichste Element ist. Also haben wir alles, was das Wasser gegeben hat. Folgt nun:

\subsection{Der sechste Tag}

\begin{quote}
V. 24—27. Und GOtt sprach: Die Erde bringe hervor lebendige Thiere, ein jegliches nach seiner Art, Vieh, Gewürm und Thiere auf Erden, ein jegliches nach seiner Art. Und es geschah also. Und GOtt machte die Thiere auf Erden, ein jegliches nach seiner Art, und das Vieh nach seiner Art, und allerlei Gewürm auf Erden nach seiner Art. Und GOtt sahe es für gut an. Und GOtt sprach: Lasset uns Menschen machen, ein Bild, das uns gleich sei, die da herrschen über die Fische im Meer, und über die Vögel unter dem Himmel, und über das Vieh, und über die ganze Erde, und über alles Gewürm, das auf Erden kreucht. Und GOtt schuf den Menschen ihm zum Bilde, zum Bilde Gottes schuf er ihn, und er schuf sie ein Männlein und Fräulein.
\end{quote}

50. Bisher haben wir gehört alles, was GOtt in den fünf Tagen gemacht hat. Zum ersten Himmel und Erde; darnach alles Gewächs auf Erden, Sterne am Himmel, Thiere im Wasser und in der Luft. Hier schafft er nun am sechsten Tage allerlei Thiere, so da leben auf Erden, wilde und zahme, und auch den Menschen; und daneben wird ihm gesagt, mit andern Thieren, daß sie fruchtbar sein sollen, und sich unter einander mehren.

51. Aufs erste, wie ich oben [§ 3] gesagt habe, so sage ich auch hier, daß am sechsten Tage nicht auf einen Hauf alles gemacht ist, denn Moses schreibt hier mit wenig Worten davon, und überläuft es auf das allerkürzeste, daß GOtt habe den Menschen gemacht, ein Bild, das ihm gleich war, und sie geschaffen ein Männlein und Fräulein. Wie aber dasselbe zugegangen ist, schweigt er hier stille, wird es aber hernach erklären im andern Capitel, und dieses Tages Geschäfte weiter ausstreichen. Wie er auch sonst gemeiniglich von ersten ein Ding pflegt grob und obenhin zu schreiben, und rauh überhin zu gehen, das er hernach besser herausstreicht. Nun, davon wäre viel zu reden, was GOtt mit diesem Werk gemeint habe. Wir wollen ein wenig davon sehen.

52. Dies ist hier am meisten zu bedenken, daß der Mensch nicht geschaffen wird wie die andern Creaturen, wiewohl er darunter gezählt und auch auf die Erde gesetzt wird. Denn GOtt zuvor mit bedachtem Rath beschließt und spricht: „Lasset uns Menschen machen.“ Denn „lasset uns machen“ ist ein Wort eines bedachten Raths, welches er zuvor nicht gesagt hat. Damit er anzeigen will, daß er mit sonderlichem Fleiß den Menschen habe wollen schaffen. Das ist bereit eine große Ehre menschlicher Creatur, daß GOtt so viel mehr Fleiß an ihn gewendet hat, denn an andere Creaturen.

53. Und aus den Worten erzwinget sich's abermal, daß in der Gottheit mehr denn Eine Person sein müsse, weil er spricht: „Lasset uns machen“; item: „Ein Bild, das Uns gleich sei.“ Denn das Wörtlein „uns“ zeigt an, daß [der], der da redet, nicht allein sei, wiewohl die Juden den Text verspotten, damit, daß also eine Weise sei zu reden, auch wo nicht mehr denn Eine Person ist. Die lassen wir fahren, das Wort wird vor ihnen wohl bleiben.

54. Hierbei haben sich nun unsere Lehrer, beide alte und neue, sehr gebrochen, daß sie auslegen, was das Bild GOttes sei, darnach der Mensch geschaffen ist, und gesagt, daß dreierlei Kräfte in der Seele seien, nämlich, Gedächtnis, Verstand und Wille, damit sie ähnlich sei der heiligen Dreifaltigkeit, dem Vater, Sohn und Heiligen Geist. Darüber haben sie sich wunder
\quellseite{34}
lich müssen brechen, und sind in so viel Fragen kommen, daß man sie nimmer kann ausarbeiten. Den Sinn lassen wir gehen, und bleiben bei den einfältigen Sprüchen und Rede der Schrift, und wollen also von dem Bilde, das GOtt ähnlich ist, reden, wie Paulus davon redet, da er [1 Cor. 15, 48. 49.] spricht: „Welcherlei der irdische Mensch ist, solcherlei sind auch die irdischen; und welcherlei der himmlische ist, solcherlei sind auch die himmlischen. Und wie wir getragen haben das Bild des irdischen, also werden wir auch tragen das Bild des himmlischen.“ Und zum Ephesern Cap. 4, 22—24.: „So leget nun von euch ab, nach dem vorigen Wandel, den alten Menschen, der durch Lüste des Irrthums verdirbt; erneuert euch aber im Geist eures Gemüths, und ziehet den neuen Menschen an, der nach GOtt geschaffen ist, in rechtschaffener Gerechtigkeit und Heiligkeit.“

55. Da seht Paulus zweierlei Bilde, irdisch und himmlisch. Adam ist von Erde und irdisch, das Bild tragen alle, die von ihm kommen. Was ist es aber für ein Bild? Adam, von dem wir kommen, ist ein Sünder, seine Vernunft war verblendet, sein Fleisch verderbet mit böser Lust und Liebe zu groben Sünden und Unglauben, Mißglauben und Verzweiflung. Also ist er worden nach dem Fall, und sind also alle, die das Fleisch von ihm bringen. Er war aber nicht also geschaffen. Wiederum, das himmlische Bild ist Christus, der war ein Mensch voll Liebe, Barmherzigkeit und Gnade, Demuth, Geduld, Weisheit, Licht und alles Gutes, also, daß all sein Wesen dahin gerichtet war, daß er jedermann dienete, und niemand schädlich wäre. Deß Bild müssen wir auch tragen, und ihm gleichförmig werden. In das Bild gehört auch, wie er gestorben ist und gelitten hat, und alles, was an ihm ist; seine Auferstehung, Leben, Gnade und Tugend, ist alles dahin gerichtet, daß wir dasselbe Bild auch anziehen. Der Art sind nun alle himmlischen Menschen, das ist, alle Gläubigen. Also siehst du, was „das Bild“ heiße und „ähnlich sein“. Darum mußt du jenen Sinn aus den Augen setzen, den die Doctores gegeben haben, sonst wird dieser Verstand nur finsterer.

56. Im Hebräischen lauten die Worte also: „Lasset uns Menschen machen in unserm Bilde, nach unserm Gleichniß“, das ist, ein Bild, das uns ähnlich und gleich sei. Der Mensch muß ein Bild sein, entweder GOttes, oder des Teufels; denn nach welchem er sich richtet, dem ist er ähnlich. So ist’s nun hier so viel gesagt, daß der Mensch am Anfang geschaffen ist ein Bild, das GOtt ähnlich war, voll Weisheit, Tugend und Liebe etc., und kurzum, im Guten ersoffen, und ohne alle böse Lüste, gleichwie GOtt, also, daß er voll GOttes war. Das meint auch der weise Mann [Pred. Sal. 7, 30.]: „GOtt hat den Menschen aufrichtig geschaffen.“ Diese Aufrichtigkeit war natürlich an Leib und Seele, und wo Adam darinne blieben wäre, hätte er auch solche Kinder gezeugt, in welchen keine böse Lust gewesen wäre, sondern wären jedermann freundlich und dienstlich gewesen, wie denn GOtt ist. Also wären wir alle GOtt ähnlich gewesen. Das hätte man denn geheißen eine Erbgerechtigkeit, als die da gebracht und geerbt wäre von Adam durch die Geburt. Nun ist er nicht also geblieben, und ist das Bild umkommen, und wir sind dem Teufel ähnlich worden durch diese Geburt, wie wir hernach hören werden. Folgt nun am Ende dieses Capitels:

\begin{quote}
V. 28. Und GOtt segnete sie, und sprach zu ihnen: Seid fruchtbar und mehret euch, und erfüllet die Erde, und bringet sie unter euch. Und herrschet über Fische im Meer, und über Vögel unter dem Himmel, und über alles Thier, das auf Erden kreucht.
\end{quote}

57. Also hat GOtt die lebendigen Thiere und den Menschen gesegnet. Wenn nun Adam also geblieben wäre, wie er da ist geschaffen worden, so wäre die Geburt auch sein zugegangen, wären mit voller Freude und ohne alle böse Lust empfangen, daß man nicht gewahr würde einiger bösen Neigung. Gleichwie noch jetzund: wo irgend ein Mensch [ist], der voll Glaubens und Geistes ist, scheint es, als [ob] er trunken wäre, daß ihm seine Werke abgehen, ehe er es bedenkt, als [ob] ihn seine Natur trüge zu guten Werken; wie man in allen Werken Christi sieht. Sonst, wo solcher Geist nicht ist, da bedenkt und bricht man sich, wie man es wolle machen; da geht viel Rathschlags, daß man thut eines mit Lust, das andere mit Unlust, und ist so darinnen ersoffen, daß man nimmer kein rechtschaffen gutes Werk kann thun.

58. Aber der rechte Glaube thut immerdar gute Werke, also, daß er auch oftmals nicht darauf denkt noch gewahr wird, was er thut; so
\quellseite{35}
\subsection{Predigt über 1 Mose 1, 28}

\subsection{B. III, 63–66}

gar ist er im Geist ersäuft. Und solche sind auch die Besten. Denn sonst, wenn sie es empfinden und sehen, werden sie gemeiniglich hoffärtig. Also wäre es gegangen, wenn wir wären rein geblieben, daß wir so tief versenkt wären im Guten, durch welches wir GOtt ähnlich und sein Bild sind. Aber es ist nun aus, der Segen ist dahin; daß das Bild sollte sein, ist zerrissen; aber die Fruchtbarkeit ist geblieben.

59. Also verstehe, was GOttes Bild ist; nicht ein todt gemalt Ding, sondern das lebendig und rechtschaffen sei, wie GOtt ist, der es rechtschaffen macht, das so vernünftig und voll Weisheit ist, daß es regieren kann Viehe, Vögel und alle Thiere auf Erden, wie GOtt regiert mit rechter Frömmigkeit. Da wäre das Werk gangen in voller Liebe, Freude, Demuth und allem Guten, auch in den Kindern, die der Mensch gezeugt und augezogen hätte.

60. Hier wäre nun wohl viel vom ehelichen Stande zu reden, und wäre wohl gut, daß einer wohl erfahren wäre, der davon handeln sollte. Doch wollen wir etwas davon sagen, und [ich] meine, ich wolle nicht weit fehlen, denn die Schrift wird uns nicht lügen.

61. Aufs erste haben wir gehört, wie GOtt den Menschen in zwei Stücke theilt. Daher schließen wir: wie andere GOttes Werke nicht in des Menschen Kraft stehen, also auch nicht, daß ein Mensch Mann oder Weib sei. Die Sonne kann nicht sagen, ich will Mond sein; wiederum, der Mond kann nicht machen, daß er die Sonne sei, sondern ein jegliches muß also bleiben, wie es von GOtt geschaffen ist. Also muß auch ein Mann bleiben, das er ist, und kann nicht ein Weib sein; wiederum, muß das Weib ein Weib bleiben, wie es gemacht ist, und steht auch nicht in ihrer Gewalt, solches zu wandeln.

62. Aufs andere hat GOtt zum Menschen gesagt, als er geschaffen war: „Seid fruchtbar und mehret euch.“ Dieser Spruch ist ein Donnerschlag wider des Papstes Gesetz, und gibt Urlaub allen Pfaffen, Mönchen und Nonnen, ehelich zu werden. Denn wie die Sonne leuchten muß, und sich nicht enthalten kann, denn es also eingepflanzt ist in ihre Natur durch GOttes Wort und Gebot: also ist es auch in des Menschen Natur eingepflanzt, daß er muß fruchtbar sein, es sei Männlein oder Fräulein.

63. Das nun GOtt etliche ausgesucht, als, die verdorbenen und untüchtigen Leiber, und etliche hohe Geister, das gehört zu andern Wunderzeichen; aber die Natur insgemeinhin steht in keines Menschen Kraft zu wehren oder halten, sondern muß ihren Gang haben und sich züchten. Darum steht es nicht in meinem Willen, solches zu verhalten, denn es uns nicht möglich ist zu halten. Darum gilt kein Gelübde dawider nichts. Denn das ist stracks beschlossen, das Werk kann niemand wehren, das GOtt gemacht hat.

64. Was wäre es, daß die Sonne wollte geloben, nicht mehr zu scheinen? So viel ist es auch, wenn du gelobst, nicht fruchtbar zu sein, noch Kinder zu zeugen oder [zu] tragen. Gelobe oder gelobe nicht, so kannst du dich nicht anders machen, denn wie dich GOtt geschaffen hat. Du wollest oder wollest nicht, so mußt du thun, wie die Natur ist, oder geht doch andere Wege, daß solcher Jammer daraus folgt, der nicht zu nennen ist. Und es soll auch also gehen, wo man GOttes Werk wehren und hindern will. Möchtest du doch auch wohl geloben, du wollest kein Mensch sein, und mußt es doch sein, und dein Gelübde mit Füßen treten.

65. Also siehst du, wie die Gelübde und Papsts Gebote stracks wider GOttes Ordnung und Einsetzung streben. Darum fast kein greulicher Ding auf Erden ist, denn das man heißt coelibatum, das ist, unserer Geistlichen Keuschheit. Willst du nun dem Greuel entlaufen, so tritt nur das Gelübde und den schändlichen Stand mit Füßen, so du darinne bist, oder hüte dich davor, daß du nicht hinein gerathest. Hast du gelobt, so wisse, daß es nicht gilt noch bindet. Denn wie kannst du das geloben, das nicht dein ist, noch in deiner Macht steht zu gewinnen noch zu halten? Es ist also von GOtt eingesetzt, daß du sollst ehelich sein; da\footnote{Wittenberger und Jenaer: daß.} ist kein freier Wille. Kannst du doch nicht eine Elle, ja, nicht einen Finger breit zu deinem Maß legen, wie Christus [Matth. 6, 27.] sagt; so wenig kannst du Fleisch und Blut wehren seine Kraft und Natur, sich zu besamen, die ihm von GOtt eingepflanzt ist.

66. Aber damit will ich der hohen Tugend der Jungfrauschaft nichts abgebrochen haben, denn GOtt der Allmächtige hat ihm seine Macht inne behalten, über die Natur zu wirken. Darum sollen wir ihm dasjenige lassen; wo wir aber sehen, daß er nicht die übernatürliche Kraft wirkt,
\quellseite{36}
das natürliche Werk gehen lassen, sich zu ehelichen, wie er es geschaffen hat. Wo Mann nicht gefallen wäre, so wäre kein Mann noch Weib unfruchtbar geblieben; nicht daß es geboten wäre, sondern daß es in die Natur gepflanzt ist, und also gehen muß ungewehrt und unaufgehalten. Nun aber hat Gott sein Werk darein geschlagen. Daß etliche untüchtig sind, gibt er ihnen eine sonderliche, hohe Gnade, daß sie ohne das leben [Matth. 19, 11.].

67. Welchen nun solche Gnade befreiet, der danke Gott und folge ihr; welchem sie aber nicht gegeben ist, der begebe sich in das gemeine eheliche Leben. Thut er es nicht, so wird er es viel greulicher machen. Wie es denn jetzund in der Welt vor Augen geht, die jetzund alle voll Hurerei und Buhlerei worden ist, und nun eingerissen hat der Pfaffenstand, den der Papst aufgerichtet hat zu Gottesdienst, und in die Leute gebracht: wer mit Gott und geistlichen Sachen wolle umgehen, der dürfte kein ehelicher Mann oder Weib sein; und also die jungen Leute von dem Ehestande geschreckt, daß sie nur in Hurerei ersäuft würden. Das sind rechte Teufelslehren [1 Tim. 4, 1. 3.], daß auch keine schädlichere auf Erden hat kommen mögen.

68. Daher ist es kommen, daß sie das eheliche Leben nicht für einen christlichen Stand noch für ein gut Werk gehalten, und haben nicht gesehen, daß im alten Testament die höchsten Patriarchen, die Gott am höchsten gedienet haben, ehelich gewesen sind, und oft viel Weiber gehabt haben. Und den Priestern war sonderlich von Gott geboten, daß sie mußten Weiber haben, und dennoch Gott dienen. Welches er darum gethan hat, daß er die Creaturen in ihrem Schwange und Erben ließe gehen, wie er sie gemacht hatte, daß man nicht sollt ander Ding anfangen wider sein Wort und Ordnung.

69. Also haben wir einen klaren Text, der alle Mönche, Nonnen und Pfaffen absolvirt von ihren Gelübden. Das sind zwei Gottes Sprüche, da ziemt nichts wider zu thun, reden, noch geloben und anzusagen. Denn Gott soll in seinen Worten und Werken recht haben. Darum, bist du ein Fräulein oder Männlein, und bist fruchtbar, so greife nur frisch und fröhlich zur Ehe, auf Gottes Wort, der diesen Stand gesegnet hat.

70. Weiter schließen wir auch hieraus, daß nicht möglich ist zu erkennen, was ein Weib oder Mann sei, denn in dem Glauben. Denn Gottes Worte und Werke sind hier beschrieben; aber es kann weder Wort noch Werk niemand verstehen, denn durch den Glauben und Geist. Das geschieht wohl, daß man böse Lust zu einem Mann oder Weib habe, aber das heißt nicht Mann oder Weib erkannt. Denn wer da soll wissen, was ein Weib sei, der muß also geschickt sein, daß er sie halte für Gottes Werk; welches allein der Glaube thut. Denn Vernunft vermag es nicht, sie ist zu blind, denkt nicht mehr denn: Hätte ich nur diesen, oder diese; sieht nicht, ob es Gottes Werk oder Creatur sei, sondern fährt nur in ihrer tollen Lust hin, hebt die Augen nimmer so hoch, daß sie Gott erkenne in seinen Werken.

71. Daher kommt es, daß sie den ehelichen Stand verachtet, denn sie nichts darinne findet und sieht denn Jammer und Unglück, welches von Noth wegen folgt, wo eines das andere nicht kennt. Denn wo der Glaube und Erkenntniß Gottes und seines Werks nicht ist, da kann nimmer kein Friede noch Ruhe sein. Ehe man hinein kommt, meint man, es soll eitel Lust darinnen sein. Wenn man nun darinne ist, und meint zu finden, was man gesucht hat, und der Vorwitz ein wenig gestillt ist, so ist darnach nichts mehr, was überbleibt, denn eitel Unlust. Das ist denn die Schuld und Ursache, daß du einen Mann oder Weib nicht genommen hast wie einen Mann und Weib, sondern wie dich die alte böse Natur und böse Lüste getrieben hat. Wenn nun aber geistlich wären geblieben in unverrückter Natur, wie Adam am ersten war, so erkennete auch jedermann einen Mann wie einen Mann, und wiederum; so wäre es voll Gutes, und folgete keine Unlust.

72. Also ist Mann und Weib, Kinderzeugen und Früchte tragen gewißlich Gottes Werk. Warum ist es denn uns so schwer, daß wir es nicht können leiden? Darum, daß das Gottes Werk nicht kann getragen werden, denn von denen, so da glauben. Wenn wir Christen wären, so könnten wir es recht erkennen. Nun ist die ganze Welt eitel Unglaube, darum greift sie es also an, daß nichts Gutes folgt. Derhalben ist die Schuld nicht der Creatur, noch des Standes, wie die Heiden klagen, sondern des Unglaubens. Es ist an ihm selbst alles gut und köstlich; weil du aber im Unglauben bist, so geht über dich der Spruch im Psalm [Ps. 18, 27.], den Gott sagt: „Bei den Verkehrten wirst
\quellseite{37}
du verzehret sein.“ Weil du verzehrt bist, so verzehrt Gott alles mit dir, daß es eitel Jammer und Noth ist; wie wir hören werden, daß es Adam wird gehen.

73. So merke nun wohl, daß diese Worte eitel Worte des Glaubens sind, daß sie niemand verstehen kann denn im Glauben; der will überall in allen Sachen und Geschäften sein. Adam, da er noch stand, war er voll in allem Guten; alles, was er angefangen hätte, wäre mit Lust und Freuden abgegangen. Als er aber fiel, mußte er dennoch ein Weib haben. Da kehrte sich’s nun um, daß ihm alles Werck eitel Jammer und Unlust war. Daher folgt, daß alles Unglück auf Erden allein des Unglaubens Schuld ist. Der Unglaube aber ist, daß man nicht verstehet, was Gottes Werke sind. Das ist das Stück, daß Gott anfühet Mann und Weib, Kinder zu zeugen, daß man lerne, daß es Gottes Werk sei, und also kein muß, und die Natur unverbittert sein will. Folgt nun das letzte Stück in dem ersten Capitel:

\begin{quote}
V. 29. 30. Und Gott sprach: Sehet da, ich habe euch gegeben allerlei Kraut, das sich besamet auf der ganzen Erde, und allerlei fruchtbare Bäume, und Bäume, die sich besamen, zu eurer Speise und aller Thiere auf Erden; und allen Vögeln unter dem Himmel, und allem Gewürm, das das Leben hat, daß sie allerlei grün Kraut essen. Und es geschah also.
\end{quote}

74. Das ist je ein feiner, freundlicher Vater. Als er alle lebendigen Thiere und Menschen geschaffen hat, versorgt er sie auch mit allem, was sie haben sollen. Wir sehen jeztund die gemeinen Fälle in der Welt, warum der eheliche Stand schwer wird, und niemand gerne daran will, daß jedermann sich fürchtet, er könnte sich nicht ernähren. Darum bleibt es so voll Auben und Hurerei. Die Schuld ist allein des Unglaubens, der uns allen natürlich anhängt nach Adams Fall. Darinne müssen wir bleiben so lange, bis uns Gott andere Menschen macht, und den Glauben in das Herz gibt.

75. Dies sind nun aber des Glaubens Worte. Denn ehe Adam ein Wort sagt noch gebeten hatte, ja, ehe er denkt, wo er sich ernähren soll, kommt Gott zuvor, und gibt ihm alles Kraut, das sich besamt, das ist, allerlei Getreide, Korn, Weizen, Roggen, Gerste, Hafer, Hirse, Reis etc., daß er sich davon nähre. Daher wir auch unsere Speise haben. Ob es Adam nicht also zugerichtet hat, wie wir, das gehe seinen Weg.

76. Und hier siehst du abermal, wie uns kein Körnlein noch Blättlein wächst, denn Gott gebe es, sintemal der Text klar spricht: „Ich habe euch gegeben allerlei Kraut“ etc. Das versteht aber niemand denn der Glaube. Item, zum andern gibt er nicht allein das, sondern auch allerlei Bäume und Früchte, Aepfel, Birnen, Weintrauben, Feigen, Oel etc., daraus man nicht allein Speise, sondern auch Getränke macht. Also hat er den Menschen versorget mit Essen und Trinken; wo hat er aber die Kleider gelassen? Deren hat er keine Nothdurft gehabt, denn er wäre also nacket gegangen; aber nach dem Fall mußte er sie haben um der Schande willen, wie wir hören werden.

77. Nun siehe, ob Moses in diesem ersten Capitel nicht meisterlich den Glauben lehrt. Am letzten Tage macht er den Menschen, aber zuvor bauet er ihm ein Haus; macht ihm das Licht am Himmel, daß er sehen kann; scheidet das Wasser von der Erde, daß er Raum hat, darauf zu wohnen, und schmückt sie um seinetwillen mit allerlei Gewächse; gibt ihm dazu das Regiment über alles, was da lebt; daß wir ja sehen sollen, daß Gott unser nicht vergesse, sondern als unser lieber Vater mit allen Gütern versorgt, und alles vorbereitet, ehe er den Menschen macht; und als er ihn gemacht hat, befiehlt er ihm auch, und ordnet ihn dazu, daß er sich von der Erde nähre.

78. Und wahrlich, wer dieses Capitel allein wohl faßte, der hätte große Lehre genug vom Glauben, daraus wir Gott vertrauen, und sehen, wie er alles gibt, was wir haben sollen; noch hat der Unglaube alles Herzeleid. Daß wir uns nicht ernähren, fehlt an Creaturen nicht; denn es ist alles voll, und haben alle genug; sondern es fehlt allein am Glauben. Darum scharren und kratzen wir so viel, auf daß wir ja nicht dürfen glauben; davon haben wir auch also viel Unlust und Mühe zu Lohn. Das sind die Werke Gottes, die er hat geschaffen in den sechs Tagen, darum beschließt er nun, und spricht:

\begin{quote}
V. 31. Und Gott sahe an alles, was er gemacht hatte, und siehe da, es war alles sehr gut. Da ward aus Abend und Morgen der sechste Tag.
\end{quote}
\quellseite{38}

79. Das Ansehen, haben wir oben [§ 14] gehört, ist Gottes Wohlgefallen und Lust, so er an [den] Creaturen hat. Denn er hat es fort bestätigt, wie er es angefangen hat, also, daß nicht allein Adam befohlen ist, Kinder zu zeugen, sondern auch bei GOtt für gut angesehen, daß es ihm wohl gefällt, und noch\footnote{Erlanger: nach.} immer so gehen muß. Darum sollen wir es auch als Gottes Werk ansehen, und für gut halten. Das wird aber niemand thun denn der Glaube.

80. Bisher ist nun kürzlich beschrieben, woher alle Creaturen kommen im Himmel und [auf] Erden; und ist endlich das der Beschluß gewesen, daß alle Creaturen geschaffen sind durch GOttes Wort, also, daß sie von ihnen selber noch von eigenen Kräften nichts sind noch vermögen, sondern wie GOtt gebeut, also geht es alles frei unverhindert.

\section{Das andere Capitel}

B. 1—3. Also ward vollendet Himmel und Erde mit ihrem ganzen Heer. Und also vollendete GOtt am siebenten Tage alle seine Werke, die er machte. Und ruhte am siebenten Tage von allen seinen Werken, die er machte. Und segnete den siebenten Tag, und heiligte ihn, darum, daß er an demselbigen geruhet hatte von allen seinen Werken, die GOtt schuf und machte.

1. Aufs erste, da unser lateinischer Text hat: Perfecti sunt coeli et terra, et omnis ornatus eorum, „also ward vollendet Himmel und Erde, mit allem ihrem Schmuck“, heißt [es] nach dem Hebräischen also: Et omnis militia eorum, das ist, „mit allem ihrem Heer“, und ist besser also verdeutscht, um der Propheten willen, die dies Wort Mosis einführen, und daher gründen, von welchem man sonst nicht wüßte, woher es käme, welches man auch täglich singt in der Messe: Sanctus, Sanctus, Sanctus Dominus Deus Zebaoth, das ist, Heilig, Heilig, Heilig ist der Herr, ein GOtt der Heerscharen, Jes. 6, 3. Desselben Wortes braucht Moses hier, und nennt Sonne, Mond und Sterne ein Heer des Himmels; und wiederum, Menschen, Thiere, Vögel, Fische, und was auf Erden wächst, nennt er das Heer der Erde. Gleich als wenn ein Fürst mit seinem Heer zeucht, so zeucht er mit ganzer Gewalt, hat seine Macht und Kraft alle bei einander.

2. Warum GOtt nun eben solches Wortes brauche, müssen wir ihm lassen befohlen sein; denn es ist nicht ohne Ursach. Denn freilich ist es also, daß alle Creaturen, die GOtt geschaffen hat, ein Heer sind, damit sie ihre Kraft und Macht ausrichten, und täglich im Streit stehen, also, daß sie den Frommen dienen zum Besten, den Bösen zum Aergersten. Wo ein gut Gewissen ist, dem steht es alles bei, wie ein gut Heer; wiederum, wo das Gewissen böse ist, ficht alles dawider; also, daß sich die Worte in der Schrift auf des Menschen Gewissen ziehen, daß die demselben dienen, oder ihm wider sind. Und kommen also aus dem Spruch viel Sprüche der Schrift, die man sonst nicht weiß, wohin sie zu ziehen sind.

3. Aufs andere ist hier zu merken, daß er spricht, GOtt hat geruhet am siebenten Tage von allen seinen Werken, die er gemacht hat; und nicht allein das, sondern auch denselben Tag gesegnet und geheiligt. Das ist auch ein Großes, davon auch viele geredet und geschrieben, aber wenige verstanden haben. Sabbath, das hebräische Wort, heißt Ruhe oder Feier, also, daß man stillehalte und stehe von allen Worten und Werken, und allein an GOttes Werken hange. Wie aber das zugegangen sei, daß GOtt also ruht von allen seinen Werken, ist freilich eine hohe Frage, wiewohl es Moses schlecht und einfältig dahin beschreibt. Wie denn seine Weise ist, daß er oft so herein fährt, und also redet, daß es jedermann versteht, und wiederum, die Worte setzt, daß es die ganze Welt nicht versteht.

4. Wir sehen täglich vor Augen, daß noch immerdar allerlei geschaffen wird, alles, was sich begibt auf Erden, alle Frucht und alle Thiere, und ist doch ein Werk, das GOtt eigentlich zugehört; wie Christus selbst sagt Joh. 5, 17.: „Mein Vater wirket bis hieher, und ich wirke
\quellseite{39}
auch.“ Wie sollen wir denn die zwei zusammenfügen, daß die Schrift zeuget und wir vor Augen sehen, daß GOtt ohne Unterlaß schafft und wirkt, bis an den jüngsten Tag, und hier Moses dagegen sagt, er habe geruht am siebenten Tage von allen Werken? Ich habe Sorge, es sei höher, denn man es geben könnte für den gemeinen Mann; doch müssen wir ein wenig davon reden.

5. Wenn man die Creatur ansieht, eine gegen die andere, so hat es gar viel einen andern Blick, denn wenn man sie vor GOtt ansieht. Eben [Cap. 1, § 2] ist gesagt, daß vor GOtt der Anfang der Welt und das Ende gleich als auf Einen Augenblick sind, und gleich gilt der erste Augenblick und der letzte am Ende der Welt. Aber wenn wir es unter uns ansehen, so kommt immerdar eines nach dem andern: der Sohn nach dem Vater, ein Jahr, ein Tag nach dem andern; aber das alles, wie es nach einander geht, ist vor GOtt als ein Augenblick. Welches die Vernunft nicht wohl begreifen kann, denn es ist eine Vergleichung zeitlicher und ewiger Dinge. Weil GOtt ewig ist, so muß es bei ihm gleich sein, vor Anfang und nach. Denn was ewig ist, dem ist alle Creatur gegenwärtig, keine zuvörderst noch zuletzt, und kann je nichts vor ihm oder nach ihm sein. Wenn man nun die Welt ansieht, vom Anfang bis zum Ende, so geht vor den Leuten eines nach dem andern, vor GOtt aber alles mit einander zugleich. Also fasse diese zwei Anblicke.

6. Daß nun Moses sechs Tage beschrieben hat, ist eigentlich also geschehen; aber das ist auch wahr: daß vor unsern Augen sechs Tage macht, ist vor GOtt alles Ein Tag, ja ein Augenblick. Also wird sich’s nun reimen, daß, da der siebente Tag angefangen ist, immer währet bis an der Welt Ende, ja, darnach wird es erst recht angehen. Ist aber also geschrieben um unsertwillen, daß es sich auf unsern Verstand ließe fassen, daß es also angefangen sei.

7. Es ist aber geschrieben allein um unsertwillen, nicht der Engel oder anderer Creaturen. Denn die Sonne feiert gar keinen Tag nicht, sondern leuchtet einen Tag eben so hell als den andern, und geht immerdar ihren Lauf vor sich hin; uns aber ist etwas hiemit angezeigt, nämlich auf das einfältigste, daß GOtt damit den Menschen hat wollen regieren, welchen er geschaffen hat zum Bilde, das ihm gleich wäre. Weil nun GOtt solche Tage schafft und wirkt, und am siebenten feiert, hat er es darum lassen schreiben, daß wir thun sollen, wie er. Denn das ist GOttes Bild, daß eben also gesinnt ist, solchen Verstand und Licht hat, und solche Werke thut wie GOtt, und sich immer nach ihm omet.\footnote{In den alten Ausgaben: „sich nach ihm omet“, das heißt, sich nach ihm richtet und bildet, ihm nachahmt.}

8. Wie aber? Sintemal dies geschrieben ist, ehe der Mensch in die Sünde fiel, da dies denn bereits vorgeschrieben ist, und die Schrift den Sabbath viel eher meldet, denn wie Adam in die Sünde fiel. Ist denn dasselbe Mal auch bereits geordnet, daß er sechs Tage arbeiten sollte, und am siebenten feiern? Da ist kein Zweifel an, wie wir hören werden, daß er sollte arbeiten im Paradies, und vorstehen Fischen, Vögeln und Thieren auf Erden. Darum hat ihn GOtt nicht wollen müßig gehen lassen, und wäre gleichwohl geblieben, daß er arbeitete und regierte, ohne daß er es nicht mit Unlust gethan hätte, und wäre ihm nicht sauer worden, mit Schweiß und Angst; auch hätte die Erde nicht Disteln getragen, sondern wäre eitel köstlich Gewächs heraus kommen, wie er gewollt hätte; daß auch kein Thier nichts mit Ungehorsam gethan hätte, und alles gegangen wäre, wie er gewollt hätte.

9. Also wäre die Arbeit und das Regiment geblieben, doch ohne Mühe und Unglück; wie auch die Weiber Kinder getragen hätten, aber nicht mit Angst, Jammer und Noth. Aber am siebenten Tage wäre es auch alles stille gethan, und hätte geruht. Das sei einfältig vom Sabbath oder siebenten Tage gesagt; was er aber bedeutet, will ich lassen anstehen. Denn wir müssen hier bei dem einfältigen Text bleiben. Es ist nun ein anderer Sabbath, denn es ge
\quellseite{40}
wesen wäre, wo Adam blieben wäre; und ist nun auf Christum gedeutet, der hat den rechten Sabbath gehalten und gefeiert im Grabe; wie wir anderswo auch gesagt haben. Spricht nun Mosis weiter:

B. 4—6. Dies ist die Geburt Himmels und Erden, da sie geschaffen sind, zu der Zeit, da Gott Erde und Himmel machte, ehe denn irgend ein Sträuchlein war auf Erden, oder irgend ein Kraut auf dem Felde wuchs. Denn Gott der Herr hatte noch nicht regnen lassen auf Erden, und war kein Mensch, der das Land bauete. Aber ein Nebel ging auf von der Erde, und feuchtete alles Land.

10. Das ist eine wunderliche Rede. Denn es fehlt dem Capitel nichts, denn daß man der Sprache nicht kundig ist; wer die nicht weiß, der wird sich fühlen\footnote{„sich fühlen“ wird hier wohl so viel sein als sich abarbeiten. So mag es auch in unserer Ausgabe Bd. VII, 1847 genommen werden.} und brechen, daß er nicht weiß, wo er ein oder aus soll; wie auch Augustinus widerfahren ist. Das ist Mosis Weise, daß er ein Ding oft wiederholt und noch einmal sagt, daß es auch schier verdrießlich ist. Darum, was hiernach folgt, das soll alles nur eine Repetition sein, und eine Erklärung deß, daß er zuvor gesagt hat. Denn im ersten Capitel, v. 27., hat er es alles mit kurzen Worten geredet: „Gott schuf einen Menschen nach seinem Bilde.“ Item: „Er schuf sie ein Männlein und Fräulein“; damit ist noch nicht ausgedrückt, wie dasselbe alles zugegangen ist. Darum holt er es in diesem Capitel wieder, thut viel Worte dazu, daß er es erkläre, wie es nach einander gegangen sei.

11. Das nun Moses hier redet, wird also leichter zu vernehmen, wenn man die Worte im lateinischen Text alle macht in praeterito plusquamperfecto; das er aber nennt die Geburt Himmels und Erde, ist so viel geredet: Da Himmel und Erde geschaffen ist. Warum aber Mosis eben also redet, daß gehe seinen Weg, er hat seine eigenen mysteria. So ist nun alles so viel gesagt: Gott schaffte und richtete zu Himmel und Erde, wie gesagt ist; aber also, daß keine Creatur etwas dazu thäte, noch hätte mögen dazu thun, sondern es ging nur ein Nebel oder dunckel Gewölke auf, wie ein Thau, und machte die Erde feucht. Das ist alles geschehen ohne den Menschen und andere Creatur. Das ist Summa Summarum davon: Es ward geschaffen, ehe je ein Mensch da war, ja, ehe es noch einmal regnete, oder ein Sträuchlein gestanden wäre. Und sagt weiter:

B. 7. Und Gott der Herr machte den Menschen aus Staub von der Erde.

12. Ich habe droben gesagt (will er sprechen), daß Gott den Menschen geschaffen hat ein Männlein und Fräulein; nun will ich sagen, wie es zugegangen ist. De limo terrae, haben wir im lateinischen Text, das ist, vom Schlamm, heißt auf Hebräisch Aphar, und ist eben das Wort, das er hernach verdolmetscht pulvis, da er sagt: Pulvis es, et in pulverem reverteris. Es heißt aber eigentlich eine solche Erde, die aufgegraben ist, und ein wenig aufgeworfen, wie ein Land, das gepflügt ist, oder von einem Grabe; aber noch nicht Staub, welcher in die Luft fliegt. Von solcher loser Erde hat er genommen einen Schrollen, und den Menschen davon gemacht. Was mehr?

\subsection{Und blies in sein Angesicht einen lebendigen Odem}

13. Dieser Worte müssen wir gewohnen, denn es ist nicht deutsch geredet. Man kann es aber nicht besser verstehen, denn wenn man einen Menschen ansieht. Den Odem, den ihm Gott gegeben hat, findet man allein im Angesichte, und sonst an keinem Orte, und vornehmlich in der Nase. Darum auch das Wort Aph, das hier steht, heißt beide das Angesicht und die Nase. Daß wir aber in unserm Text haben spiritum vitae, soll nicht ein Geist heißen, sondern so viel, daß ihm Gott ein solches Leben gegeben hat, nicht wie die Fische, sondern wie die Thiere haben, als eine Kuh, Pferd, Hirsch, und alles solches, was Odem hat. Und wird damit angezeigt, daß unser Odem auch nicht in unserer Gewalt steht, noch daß wir von uns selbst schnauben, noch Odem holen können; und von diesem Odem leben wir. Denn wenn der verstopft wird, so sind wir todt. Darum heißt er es einen lebendigen Odem, daß er den Menschen bei dem Leben erhält, und ein Zeichen ist, daß der Mensch lebe.

\subsection{Und also ward der Mensch eine lebendige Seele}

14. Den Spruch zieht Paulus an in der ersten Epistel an die Corinthier, Cap. 15, 45.: „Also
\quellseite{41}
ist's geschrieben (spricht er): Der erste Mensch, Adam, ist gemacht ins natürliche Leben.“ Das Wort „Seele“ müssen wir auch lernen recht verstehen. Das heißen wir in unserer Sprache eine Seele, die, so bald der Mensch stirbt, vom Leibe fährt. Aber Moses und die Schrift heißen das Seele, alles, was da lebt in den fünf Sinnen; als, das heißt sie auch Seele, daß ein Fisch im Wasser lebt, wie er es oben im ersten Capitel, V. 20., genannt hat: Producant aquae animam viventem, und ich gedeutscht habe „lebendige Thiere“. Item, als die Vögel in der Luft, und die Thiere auf Erden leben, daß es auf das eigentlichste heißt ein Leibesleben oder ein lebendiger Leib. Also, daß der Spruch eigentlich auf den Verstand gehe, daß der Mensch geschaffen ist in das leibliche Leben, das wir heißen das natürliche Leben.

15. Also verstehe den Spruch Pauli: „Der erste Mensch ist gemacht ins natürliche Leben.“ Denn daselbst setzt er gegen einander ein leibliches und geistliches Leben. Das leibliche Leben ist, daß man hört und sieht, riecht, greift, schmeckt, däuet, zu sich nimmt und auswirft, Kinder zeugt, und was der Leib für natürlich Wesen und Werk hat; das heißt die hebräische Sprache „Seele“. Also lesen wir im andern Buch Mosis [Cap. 1, 5.]: „Aller Seelen, die aus den Lenden Jakobs kommen waren, der waren siebenzig“, das ist, siebenzig Kinder, die von Jakob\footnote{Erlanger: „ihm“ statt Jakob.} geboren waren. Das ist nun fast gemein in der Schrift durch und durch.

16. Item, also verstehe auch, daß Christus sagt Matth. 10, 39. 16, 25.: „Wer seine Seele findet, der wird sie verlieren“, das ist, sein natürlich Leben. Denn er meint nicht allein das Abscheiden der Seele, sondern will also sagen: Daß man sich um das ewige Leben in Fahr Leibes und Lebens\footnote{„in Fahr Leibes und Lebens“ fehlt in der Erlanger. In der Jenaer stehen diese Worte in Klammern.} müsse setzen mit dem leiblichen und natürlichen Leben. Darum kann man das Wort Seele nicht besser deuten, denn das leibliche Leben, oder einen Menschen, der da lebt im leiblichen Leben. Also ist Adam ins natürliche Leben gemacht; Christus aber, der der letzte Adam ist, spricht Paulus [1. Cor. 15, 45.], ins geistliche Leben; das ist, er hat einen geistlichen Leib, also, daß er nicht mehr isset noch trinkt, sieht noch hört, wie wir, thut kein

leiblich Ding noch Werth, sondern ist gar ein ander Wesen, und doch wahrer Mensch; wie wir auch in jenem Leben sein werden.

\subsection[V. 8—14]{V. 8—14}

\begin{quote}
Und GOtt der HErr pflanzte einen Garten in Eden, gegen dem Morgen, und setzte den Menschen darein, den er gemacht hatte. Und GOtt der HErr ließ aufwachsen allerlei Bäume, lustig anzusehen, und gut zu essen, und den Baum des Lebens mitten im Garten, und den Baum des Erkenntniß Gutes und Böses. Und es ging aus\footnote{Wittenberger und Erlanger: an.} von Eden ein Strom, zu tränken den Garten, und theilte sich daselbst in vier Hauptwasser. Das erste heißt\footnote{Am Rande: Pison ist das große Wasser in India, das man Ganges heißt, denn Hevila ist Indienland. Gihon ist das Wasser in Egypten, das man Nilus heißt. Hiddekel ist das Wasser in Assyria, das man Tigris heißt. Phrath aber ist das nächste Wasser in Syria, das man Euphrates heißt.} Pison, das fleußt um das ganze Land Hevila, und daselbst findet man Gold, und das Gold des Landes ist köstlich; und da findet man Bedellion, und den edlen Stein Onyx. Das andere Wasser heißt Gihon, das fleußt um das ganze Mohrenland. Das dritte Wasser heißt Hiddekel, das fleußt vor Assyrien. Das vierte Wasser ist der Phrath.
\end{quote}

17. Der lateinische Text hat hier recht gemacht: Plantaverat, er hatte gepflanzet. Also sollte wohl zuvor auch stehen, formaverat und inspiraverat, also, daß man müßte, daß es eine Repetitio wäre, wie ich gesagt habe. Da haben wir aber einen wunderlichen Text: Der Herr, spricht er, hatte lassen wachsen allerlei hübsche Bäume in dem Garten, den er geschaffen hatte, sonderlich aber den Baum, der da heißt ein Baum zu wissen, was gut und böse ist. Im lateinischen Text haben wir also: Plantaverat Paradisum voluptatis a principio. Ich halte aber nicht, daß es recht sei. Denn wenn er hätte wollen sagen a principio, am Anfang, so hätte er freilich ein ander Wort gebraucht. Darum klingt es im Hebräischen also, wie wir es gemacht haben: „Ein Garten Eden, gegen dem Morgen“, oder zuvor.

18. Das Wort Eden heißt ja Lust. Darum hat man gemacht Paradisum voluptatis, das ist, wie wir sagen, einen hübschen Lustgarten, darinne allerlei Bäume waren, lieblich zu sehen und lustig zu essen; und unter diesen zwei sonder
\quellseite{42}
liche Bäume mitten im Garten gegen dem Morgen: ein Baum des Lebens, und ein Baum, an dem man lernte, was gut und böse wäre.

19. Hier ist nun die Frage: wo das Paradies in der Welt sei. Daß es auf Erden ist, muß man zulassen. Denn da steht der Text: „Gott hat gepflanzet einen Garten in Eden, gegen dem Morgen“, so müssen es auch natürliche Bäume sein, wie unsere. Darum es nichts ist, daß unsere Sophisten gesagt haben, wie es hoch droben über der Erde liege, hart unter dem Monde. Es muß hier auf Erden sein, und müssen auch die Bäume sein, die Gott im ersten Capitel, V. 11., geschaffen hat. Zum andern ist ja Adam auf der Erde geschaffen, und dazu geordnet, daß er darauf sein sollte, und wird darum ins Paradies gesetzt, daß er es bauen und bewahren soll. Zum dritten werden hier vier Wasser genannt, welche man noch weiß, die aus dem Garten fließen. Dies alles beweist genugsam, daß es müsse auf Erden\footnote{Erlanger: der Erden.} sein.

20. Das will nun wunderlich Ding werden. Origenes und andere haben sich hin und her damit verworren; aber Augustinus hat klüglich gehandelt und gesagt: Wer dies und anderes nicht begreifen kann, der soll Gott die Ehre geben, und ihm befehlen. Das ist aber der Behelf auf jener Seite, daß Moses nicht sagt, daß die vier Wasser im Garten gewesen seien, sondern nur Ein Wasserstrom, davon sich jene Wasser leiten.

21. Was wollen wir denn nun daraus machen? Also habe ich mehr gesagt, und sage es noch: Möglich ist es, daß es zu der Zeit also gewesen ist, daß Gott einen Garten gemacht, oder ein Land beschränkt habe; aber nach meinem Dünken wollte ich gerne, daß es also verstanden möchte werden, daß es der ganze Erdboden wäre. Mir liegt aber das im Wege, daß der Text also klingt, daß es etwas Anderes, nämlich ein sonderlicher Ort und Raum sei, wie auch noch lustige Gärten nicht ein ganzes Land begreifen. Darum weiß ich es nicht zu örtern, wie es gewesen sei. Ich muß mich gefangen geben, weil man die vier Wasser, wie [§ 1.] gesagt, noch wohl weiß, die daraus kommen.

22. Darum wollte ich also sagen, daß der Lustgarten irgend ein Ort sei gegen dem Morgen, der nun verborgen oder vielleicht zerrissen ist, das Gott wohl weiß. Es muß aber fast ein weiter Raum gewesen sein, denn die Wasser liegen mächtig weit von einander, ja schier gegen einander. Darum will ich meine Vernunft gefangen geben, und dabei bleiben, daß es ein rechter, natürlicher Garten sei gewesen, wie noch möchte ein Lustgarten sein.

23. Denn weil Adam ein leiblicher Mensch war, und ein leiblich Wesen führte, muß auch der Garten also gewesen sein, daß man es gefühlt, geschmeckt und gerochen habe. Denn hernach klärlich steht, daß er von dem Baum gebrochen und gegessen habe, der da lehrte, was gut und böse war. Also, daß es muß eine natürliche leibliche Frucht gewesen sein, oder Adam hätte müssen ein geistlicher Mensch sein. Es sei nun der Garten, wie und wo er wolle, geben wir Gott heim; wiewohl die Schrift hernach einen geistlichen Verstand daraus macht, wie sie auch dergleichen mehr thut, als da Christus [Luc. 23, 43.] sagt zum Schächer: „Heute wirst du bei mir sein im Paradies.“

24. Was ist aber das für ein Baum mitten im Garten, daß er ihn heißt einen Baum, daran man lernt Gutes und Böses? Er hat ihn freilich um viel zukünftiges Dinges willen also genannt, und sonderlich um des Falles willen, daß Adam dabei lernen würde; wie auch die Schrift an andern Orten pflegt etlichen Dingen den Namen zu geben um zukünftiges Geschichts\footnote{Das Geschicht — das Geschehene, Ereigniß.} willen, per anticipationem. Als, 1. Sam. 4, 1. sagt der Text: „Israel\footnote{„Israel“ von uns gesetzt statt: „Die Philister“ in den Ausgaben.} lagerte sich bei dem Helfenstein“, da nennt er den Ort Helfenstein, der doch zur selben Zeit nicht den Namen hatte, und erst im siebenten Capitel, V. 12., gesagt wird, wie es den Namen kriegt von dem Geschicht, das sich darnach begab.

\subsection[B. 15—17]{B. 15—17}

\begin{quote}
Und Gott der HErr nahm den Menschen, und setzte ihn in den Garten Eden, daß er ihn bauete und bewahrete. Und Gott der HErr gebot dem Menschen, und sprach: Du sollst essen von allerlei Bäumen im Garten, aber von dem Baum des Erkenntniß Gutes und Böses sollst du nicht essen; denn welches Tages du davon issest, wirst du des Todes sterben.
\end{quote}

25. Hier siehe, warum Gott Adam das Gebot gibt, ehe denn er Evam schafft. Er hat es
\quellseite{43}
ohne Zweifel aus dem Heiligen Geiste geschrieben, und St. Paulus [1 Tim. 2, 13.] rührt es auch. Das Weib mußte nicht Gottes Wort ohne Mittel hören, sondern von Adam lernen, also, daß auch vor dem Fall das Regiment und Obrigkeit bei der männlichen Person gewesen ist.

26. Das ist nun das allererste Gebot, ist aber Adam nicht dazu gegeben, daß er dadurch fromm würde; denn das Gebot muß man weit sondern von allen Geboten, die hernach gegeben sind; denn Adam war noch ohne alle Sünde. Dies Gebot aber gab ihm Gott nur zum Zeichen, denn er mußte es also gegen den Menschen halten, daß er wüßte und gedächte, daß er einen Oberherrn hätte.

27. Also konnte er nicht fromm dadurch werden, wenn er es gleich gehalten hätte, konnte aber wohl ein Sünder werden. Welches bereit an eine große Beweisung ist, daß kein Gesetz nicht fromm machen kann, sondern wohl ärger; wird ihm aber dazu gegeben, daß er sich daran übete und beweisete, daß er fromm wäre und in Gottes Gehorsam ginge. Also gibt das Gesetz nicht die Frömmigkeit, sondern die da fromm sind, die thun das Gesetz. Das ist die Ursache des Gebots des Menschen halben. Warum es aber Gott für seine Person gegeben habe, so doch Adam also geschaffen war, daß er keines Gebots bedurfte, können und wollen wir nicht ausgründen.

28. Das andere Stück ist, daß Gott uns damit anzeigte, welche Narren es sind, die sich unterstehen und unterstanden haben, den Leuten zu helfen mit vielen Gesetzen und Lehren. Darum siehe dies Exempel wohl an, ob es gleich einfältig ist. Was will man doch uns arme Menschen treiben mit Gesetzen, Gelübden und Menschentand, so wir sehen, daß es nichts geholfen hat, daß Gott dem unschuldigen Menschen ein Gebot gegeben hat? Und wir, die nun verderbt sind, daß nichts Gutes in uns ist, sollen so viel Gesetze halten? So es Adam, der voll Gerechtigkeit war, fehlte, daß er es nicht hielt in solcher Fülle und Vorrath, daß alle Bäume voll Früchte waren, und er haben konnte, was sein Herz begehrte, ließ sich dennoch das Weib betrügen und bereden, daß er das einige Gebot nicht hielt: was wollen wir denn halten, wenn uns so viel unzählige Gebote, immer eines über das andere, aufgelegt werden? Darum laß dies einen großen, gewaltigen Puff sein wider die Narren, die sich wollen unterwinden, mit Gesetzen vor Gott zu kommen. Denn er hat es noch nie im Sinn gehabt, daß er jemand damit wolle fromm machen; sondern eigentlich, daß er uns dadurch will lernen abtreten von unserm Vornehmen. Als sollte er sagen: Hatte doch der erste Mensch, der da fromm war, nicht gehalten; was solltest du denn halten, da nichts Gutes in der Natur ist?

29. Spricht du: Warum hat er denn hernach durch Mosen die Leute mit Gebieten und Geboten\footnote{„Gesetzen und Geboten“ fehlt in der Erlanger.} überschüttet? Freilich nicht darum, daß er uns dadurch wollte fromm machen, denn er kennt unser Herz besser denn wir selbst, sondern darum, weil wir meinen und unterstehen uns mit Geboten fromm zu machen, daß wir zuletzt davon abtreten, und erkennen müssen, wie es nichts hilft. Hat uns eben darum so viel gegeben, daß wir bekennen sollen unsere Bosheit und der menschlichen Natur Gebrechen,\footnote{„Gebrechen“ fehlt in der Erlanger.} daß sie nichts Gutes will, noch gedenkt zu thun, und von der Scheitel an bis auf den Fuß böse ist [Röm. 7, 18. 19. 2 Cor. 3, 5. Jes. 1, 6.].

30. Darum sollen diesen Text billig feiern und bedeuten alle Prälaten und Bischöfe, welche die Gewissen wollen regieren. Der weltlichen Obrigkeit gebühret, mit dem Schwert zu wehren, daß niemand stehle, noch\footnote{„noch“ fehlt in der Jenaer.} morde, noch ehebreche etc. Aber daß sie die Welt wollen mit Gesetzen fromm machen, und die Gewissen dadurch regieren, da wollen wir ihr Gesetz nicht haben, sintemal sie doch nichts ausrichten, und niemand diese Gesetze hält.

31. Das dritte Stück ist, daß Gott hier anzeigt seine Barmherzigkeit, daß er Adam läßt fallen, aber bald wieder aufrichtet. Denn es steht fest, daß er fromm und gerecht geschaffen war, wie gesagt. Es war alles gut, daß Gott einen Wohlgefallen, und keinen Mangel daran hatte. Nun andern muß das auch stehen, daß er, als er von dem verbotenen Baum aß, von dem Stande fiel, der Gott am höchsten gefallen hat, in den allerärgsten Stand, der ihm nicht gefiel; denn er hat keine Lust am Tode [Ezech. 18, 32.]. Doch ist das auch klar, daß er ihn wieder zu Gnaden nimmt, erbarmt sich über ihn, und gab ihm eine gnädige Absolution, und brachte ihn wieder in den Gnadenstand, obwohl nicht so vollkommen als zuvor; doch war es
\quellseite{44}
gleich, wie wir hören werden im folgenden Capitel, V. 15., da GOtt zu der Schlange sagt: „Des Weibes Same soll dir den Kopf zertreten.“ Der Spruch ist die Absolution, damit er ihn ledig gesprochen hat und uns alle. Denn, ist der Same so stark, daß er der Schlange den Kopf zertritt, so zertritt er auch alle ihre Gewalt; so ist der Teufel überwunden, und aller Schade hinweg, den Adam hatte, und kommt in den Stand, da er vorhin inne war. Den Trost und die Gnade hat er ihm gegeben, daß er ihn wieder aufrichtete, und dahin geholfen, davon er gefallen war.

32. Dies ist nun das erste Exempel [von] GOttes Barmherzigkeit, daß uns GOtt nicht will lassen verzagen, noch verzweifeln heißen, ob jemand gefallen ist, weil wir sehen, daß Adam dahin fällt von dem höchsten Stande, und doch nicht so bleibt. Dadurch beschlossen ist, daß GOtt läßt fallen, und auch wieder aufhilft. Das ist kürzlich von dem Stück geredet.

33. Aber hier ist eine Frage: wie es zugegangen sei, weil GOtt spricht: An welchem Tage du von dem Baum mitten im Garten wirst essen, wirst du des Todes sterben, daß er doch nicht gestorben ist; denn er nach derselben Sünde je noch neunhundert und dreißig Jahre gelebt hat. Wie ist denn GOttes Wort wahrhaftig? Antwort: Das müssen wir lassen bleiben, daß Adam in diesem Gebot, wie wir hören werden, uns allzumal verkündigt hat [Röm. 5, 12. 1 Cor. 15, 21.], denn wir sind allzumal in ihm gepflanzet, und sein Blut und Fleisch, daß es uns muß gehen, wie es ihm gegangen ist. Denn GOtt hat es geordnet, daß von diesem einigen Menschen alle Menschen kommen müssen, und wir alle seine Kinder sind. Darum, wie er gethan hat, und was ihm zum Fluch auferlegt ist, dasselbe begegnet uns allen, also, daß wir müssen mit ihm für Einen Kuchen und Teig gerechnet werden; kürzlich, alles, was Menschen sind. Darum auch die Schrift den ersten Menschen und uns allen Einen Namen gibt, daß, was Mensch ist, alles Adam heißt, von diesem ersten Adam, das ist, vom ersten Menschen, von dem wir alle gemacht sind. Darum kein Esel den andern (wie man spricht) darf einen Sackträger heißen. Darum ist hier niedergelegt mit ganzer Gewalt aller Hochmuth, daß sich ein Mensch über den andern erheben will. Wir sind Adam, und bleiben Adam.

34. Warum hat nun GOtt gesagt: „Welche Stunde du von dem Baume essen wirst, wirst du des Todes sterben“? Etliche haben es also ausgelegt, daß es so viel sei: Du wirst sterblich sein; oder, die es schärfer machen: Du wirst des Todes schuldig sein. Das halte ich auch für den rechten Verstand. Denn, sintemal wir alle in der Sünde sind, so sind wir auch in dem Fluch und Strafe, bis so lange der Same kommt, der es beides hinweg nimmt. Wahr ist es, daß Adam und Eva nicht so bald gestorben sind, und auch den Tod nicht so bald fühlten, so sie doch so bald fühlten, daß sie nacket waren, und ihnen Schürze machten; aber darnach, da er die Stimme hörte, fühlte er den Tod [1 Mos. 3, 8.]. Denn das ist auch wahr, da GOtt mit Adam redete, stand er in Todesnöthen. Darum wollen wir es so einfältiglich verstehen, „du wirst des Todes sterben“, daß es so viel sei, du mußt sterben, und wirst den Tod fühlen; wiewohl die Zeit nicht eigentlich bestimmt ist. Nun sagt Moses weiter, wie das Weib gemacht ist.

\subsection[B. 18—20]{B. 18—20}

\begin{quote}
Und GOtt der HErr sprach: Es ist nicht gut, daß der Mensch allein sei; ich will ihm einen Gehülfen um ihn machen. Denn als GOtt der HErr gemacht hatte von der Erde allerlei Thiere auf dem Felde, und allerlei Vögel unter dem Himmel, brachte er sie zu dem Menschen, daß er sähe, wie er sie nennete; denn wie der Mensch allerlei lebendige Thiere nennen würde, so sollten sie heißen. Und der Mensch gab einem jeglichen Vieh, und Vogel unter dem Himmel, und Thiere auf dem Felde seinen Namen; aber es fand sich dem Menschen kein Gehülfe um ihn.
\end{quote}

35. Da macht Moses aber mal eine Repetition, und ist das die Summa: Da nun alles, was da lebt, geschaffen war, brachte sie GOtt zu Adam, daß er sie nennete; aber unter den allen fand er keinen Gehülfen um ihn. Und ist so viel gesagt: GOtt sahe Adam an, daß er allein ein Mannsbild war. Nun hatte er alle Thiere geschaffen beide Sie und Er, da brachte er alle Thiere, Sie und Er, zum Adam, aber seine Sie oder Gefährten\footnote{Wittenberger: „gegaten“.} fand er nicht.

36. Unser Text liest: Adjutori um simile ei; es sollte aber heißen: coram eo, id est, adjutorium ad generationem. Es war kein Thier, das sich zu ihm gethan hätte; sie gingen hin und ließen sich zählen, wie sie noch thun, aber
\quellseite{45}
keines hielt sich zu ihm als seine Gesellin. Er nannte allein ein jegliches, wie er wollte; wie auch noch, wie man die Thiere nennt, so müssen sie heißen. Welches ein Zeichen ist, daß der Mensch ein Herr ist aller Thiere.

V. 21—23. Da ließ GOtt der HErr einen tiefen Schlaf fallen auf den Menschen, und er entschlief. Und nahm seiner Rippen eine, und schloß die Stätte zu mit Fleisch. Und GOtt der HErr baute ein Weib aus der Rippe, die er von dem Menschen nahm, und brachte sie zu ihm. Da sprach der Mensch: Das wäre einmal Bein von meinen Beinen, und Fleisch von meinem Fleisch. Man wird sie Männin heißen, darum, daß sie vom Manne genommen ist.

37. Diese Worte sind nicht Menschenworte, darum wollte ich, daß man sie nicht für gering und spöttlich hielte. GOtt redet und thut es selbst. Darum ist es ein Ernst, und lästern GOtt, die ein Gaukenspiel daraus gemacht haben. Also will er aber sagen: GOtt ließ Adam entschlafen; das ist, wenn er wirkt, wo nicht Sünde ist, so wirkt er so säuberlich, daß es die Creatur nicht fühlt. Aber mit den Sündern geht er nicht also um, die müssen es fühlen, daß [es] ihnen wehe thut. Wenn Adam gewacht hätte, würde es ihm auch wehe gethan haben; aber als er schlief, fühlte er es nicht, und da er erwachte, sieht er sie vor ihm stehen, hat des Bildes zuvor nie gesehen unter den andern. Nun war der Geist in ihm, wie es Christus [Matth. 19, 5.] auslegt, daß dies Gottes Worte sind, die hier Adam redet. Daher man schließen kann, daß Adam zu der Zeit voll Gottes gewesen sei, oder GOtt durch Adams Mund geredet habe; darum soll man es keinen Spott sein lassen. Und merkte das: als GOtt das Weib aus der Rippe des Mannes machte, braucht der Text eben das Wort „bauen“: „Er bauete ein Weib“, gerade, als sollte es ein Haus werden. Davon wir hernach hören werden.

38. Das ist nun die Meinung, wie es Paulus [Eph. 5, 28. ff.] auslegt, daß keine größere Vereinigung ist denn Mann und Weib, und wäre allenthalben also geblieben, wenn Adam in der Unschuld geblieben wäre. Nun ist es verderbt, daß auch selten Einigkeit ist unter den Eheleuten. Darum spricht Moses, Adam erkannte, daß dies sein Bild ihm gleich wäre; denn zuvor war kein Thier da, das sich gegen ihm stellete, als wollte es ihm helfen zur Geburt, nach dem Wort Gottes [1 Mos. 1, 28.]: „Seid fruchtbar und mehret euch.“ Das hat nun GOtt also gepflanzt, daß der Mensch muß ein Mann und Weib sein, und keines ohne das andere Frucht zeugen kann.

39. Und ist beschlossen, daß das Weib darum geschaffen ist, daß es des Menschen Gehülfe sei, nicht zur Lust noch Büberei, sondern daß der Spruch fortgehe. Wie nun Adam hier nicht hat mögen wehren, daß er ein Mann wäre, also auch, daß GOtt aus seiner Seite eine Rippe nahm, und machte das Weib daraus; und wie es nicht in ihrer Gewalt ist, daß sie aus der Rippe kommt, und anders nirgend her, also auch nicht, daß sie nicht fruchtbar vom Manne werde.

40. Aber das hat nichts geholfen; man hat alle Welt voll gepredigt, daß man Jungfrauschaft geloben soll. Ist wohl gut und recht, wer es halten kann, den GOtt sonderlich auszeucht, daß er so bleibt. Hat er es doch Macht, zu machen, wie er will, könnte auch wohl aus Mann Weib machen, und wiederum. Wenn wir nicht gefallen wären, so wäre es also gegangen, daß alle hätten Frucht müssen tragen; denn der Segen war gesprochen, daß es ohne Wehe und böse Lust sollte zugehen. Der ist aber nun gebrochen; so hat GOtt gemacht, daß er etliche auszeucht, daß sie nicht fruchtbar sollen sein. Darinne lasse man dem Heiligen Geist sein Werk gehen. Welche aber nicht also ausgezogen sind, da lasse man es gehen, wie es gehen soll. Denn wie wir alle hätten müssen Frucht zeugen, wo die Natur unerrückt wäre, wie viel mehr müssen wir es thun, so sie jetzt verderbt, und viel geiler ist? Wer gesund ist, kann essen, wenn er will; einen Kranken muß man dazu zwingen. Darum kann man es nicht wehren; will man es aber wehren, so macht man es ärger. Die Weiber sind nirgend um geschaffen, denn daß sie dem Manne dienen, und ein Gehülfe seien, Frucht zu zeugen.

41. Nun siehe, wie ihr Adam einen Namen gibt, und heißt sie Männin, „darum“, sagt er, „daß sie vom Manne genommen ist“. Auf Hebräisch heißt das Wörtlein Isch eigentlich ein Mann unter den Menschen. Denn Sachar heißt es auch unter andern Thieren. So nennt er sie nun von seinem Namen Ischa, daß sie den Namen von und nach ihm hat, wie es noch
\quellseite{46}
bisher bleibt, daß man das Weib nach des Mannes Namen nennt. So ist nun das Wort Ischa so viel, als wir sagen mögen eine Männin, darum, daß sie vom Manne kommt und gemacht ist. Gleichwie man das, [was] von Holz gemacht ist, hölzern\footnote{In den alten Ausgaben: hülzen.} heißt, also muß sie von ihm den Namen haben, daß er ihr ihn gebe, und das Regiment behalte.

\begin{quote}
V. 24. Darum wird ein Mann seinen Vater und Mutter verlassen, und seinem Weibe anhangen, und werden sein Zwei Ein Fleisch.
\end{quote}

42. Diese Worte sind nicht also zu verstehen, daß sie allein leiblich Ein Fleisch und Blut seien, sondern von allem, was gehört zum äußerlichen, leiblichen Leben. So heißt die Schrift „Fleisch“ alles, was zum Fleisch gehört, daß man muß haben, Hausgesinde, Kinder, Geld, Acker, Wiesen, Gut, Ehre, oder Armuth, Schande, Krankheit und Gesundheit, und so fortan, was dem Fleisch mag zufallen, also, daß Fleisch heiße ein äußerlich Leben im Fleisch. So soll es nun gehen, daß es alles beider zugleich sei, und sie sich alles zugleich annehmen, und eines dem andern bringe Leib, Gut, Ehre, Schande, Armuth, Krankheit, und was es mehr ist. Das ist ein solch Leben, das im Fleische geht, das ist, im fleischlichen Wesen, und was dazu gehört, soll alles gemein sein; ohne daß der Mann das Regiment führen soll, und sie von ihm den Namen habe. Wenn er gesagt hätte, es soll Ein Geist sein, so wäre es besser worden; nun ist es wohl Ein Fleisch und Blut, aber mancherlei Seele und Geist.

\begin{quote}
V. 25. Und sie waren beide nacket, Adam und sein Weib, und schämten sich nicht.
\end{quote}

43. Das sehen wir an allen Thieren, daß sie nicht in dem Stande sind, daß sie sündigen können; also war es auch mit dem Menschen. Aber nun ist es aus, wie wir hören werden, daß wir uns scheuen und schämen müssen. Die Schande müssen wir leiden bis an den jüngsten Tag, darnach wird es wieder also sein, daß keines sich vor dem andern wird schämen. So viel Freude wird im Himmel sein. So auch wiederum, in der Hölle wird die Qual so groß sein, daß man vergessen wird, was Mann oder Weib sei. Das Unglück wird den Kitzel wohl vertreiben. So wird es in jenem Leben wieder also werden, wie zum ersten, daß man nicht sagen wird: Siehe, das ist eine Sie, das ist ein Er; jetzt aber ist die Natur also verderbt, daß eines das andere ohne Scham nicht kann ansehen.

44. Nun sollte ich hier die geistliche Deutung auch einführen: wie Adam Christus, und am Kreuze entschlafen, und daselbst aus seiner Seite die Christenheit, seine Braut, genommen ist; welches uns ein großer Trost, Trotz und Ruhm ist. Es würde aber zu viel; darum wollen wir es den Gelehrten und müßigen Geistern befehlen weiter zu handeln.

45. So haben wir nun bisher in den zweien Capiteln gehört die Werke, die Gott geschaffen hat in den sechs Tagen, und sonderlich, wie am letzten der Mensch geschaffen ist, beide Mann und Weib, und wie sie Gott zusammen gegeben hat, und ehe er das Weib schuf, Adam geboten hat, daß er von allerlei Bäumen möchte essen im Garten, ohne allein von der Frucht des Baums, daran man lernete, was gut und böse ist. Folget nun das dritte Capitel, darinne beschrieben ist der Jammer und Herzeleid, das bald hernach ergangen ist über die menschliche Natur, und noch geht.

\section{Das dritte Capitel}

\begin{quote}
V. 1—6. Und die Schlange war listiger, denn alle Thiere auf dem Erdboden, die Gott der Herr gemacht hatte, und sprach zu dem Weibe: Ja, sollte Gott gesagt haben, ihr sollt nicht essen von allerlei Bäumen im Garten? Da sprach das Weib zu der Schlange: Wir essen von den Früchten der Bäume im Garten; aber von den Früchten des Baums mitten im Garten hat Gott gesagt: Esset nicht davon, rühret es auch nicht an, daß ihr nicht sterbet. Da sprach die Schlange zum Weibe: Ihr werdet mit nichten des Todes sterben, denn Gott weiß, daß, welches Tages
\end{quote}
\quellseite{47}
ihr davon esset, so werden eure Augen wacker, und werdet sein wie Gott, und wissen, was gut und böse ist. Und das Weib schaute an, daß der Baum fein war davon zu essen, und lieblich anzusehen, daß ein lustiger Baum wäre, weil er klug machte.

1. Das Erste, das Moses hier beschreibt, ist, wie die Schlange mit dem Weibe geredet hat. Da können wir nicht vorüber, wir müssen, wie ich immerdar sage, die Schrift lassen bleiben in einem einfältigen, bloßen Sinne, wie die Worte geben, und keine Glosse machen. Denn es gebührt uns nicht, Gottes Wort zu deuten, wie wir wollen; wir sollen es nicht lenken, sondern uns nach ihm lassen lenken, und ihm die Ehre geben, daß es besser gesetzt sei, denn wir es können machen. Darum müssen wir es lassen stehen, daß es eine rechte, natürliche Schlange gewesen sei, die das Weib mit Augen gesehen hat; und es ist darum geschrieben, daß die Historie gefaßt würde in einem leichten Verstande. Denn, sollte er so schreiben, daß der Teufel hätte mit ihr geredet in eigener Person, würde sich nicht schicken. Darum hat er müssen durch die Schlange reden, und wird also beschrieben, als habe hier die Schlange selbst geredet.

2. Er hat aber genug damit angezeigt, daß die Schlange natürlich ist gewesen, aber der Teufel hat darinne gewohnt, darum, daß er spricht, sie habe mit ihr geredet. Denn reden ist keinem Thiere gegeben, ohne dem Menschen. Darum gibt es klar genug zu verstehen, daß der Teufel in der Schlange durch ihre Zunge geredet hat. Und das soll niemand wundern, denn der Teufel ist ein gewaltiger Geist [Luc. 11, 21. Eph. 6, 12.]. Gott hat ihm nicht gewehret, daß er mit leiblichem Ding umgeht; wie wir noch sehen, daß er der Welt Herr und Fürst ist [Joh. 12, 31. 14, 30. 16, 11.], und nicht allein durch die Thiere, sondern jetzt am meisten durch die Menschen redet.

3. Zum andern ist es\footnote{„es“ fehlt in der Wittenberger und in der Erlanger.} auch ein Zeichen, daß es der böse Geist gewesen ist, daß er so hoch daher redet von Gottes Geboten. Denn je kein Thier so klug ist, daß es wisse, was Gottes Gebot sei, oder nicht. Darum muß ein solcher Verstand in ihr gewesen sein, der über der Schlangen Natur ist, ja, auch über menschliche Natur: es muß ein Engel gewesen sein; weil er aber wider Gottes Gebot handelt, so kann es kein guter Engel sein.

4. Item, es ist Eva eine Frau der Welt gewesen. Denn sie da steht, redet mit der Schlange, und fürchtet sich nichts vor ihr, achtet sie eben wie ein ander Thier, denn sie war ein Herr über alle, wie wir gehört haben, daß Gott [1 Mos. 1, 28.] zum Menschen sagt: „Herrschet über die Fische im Meer, und über Vögel unter dem Himmel, und über alle Thiere, die auf Erden kriechen.“ Es stund da noch wohl um sie, daß sie sich vor keinem Gift fürchten durfte, noch vor dem Tode, und war nichts, das ihr schaden konnte.

5. Aber deß versähe sie sich nicht, daß der Teufel da wäre. Denn also sagt Paulus 1 Tim. 2, 14.: „Adam ward\footnote{In den Ausgaben: war.} nicht verführet, das Weib aber ward verführet, und hat die Uebertretung eingeführt.“ Das ist so viel gesagt: Eva war nicht so verständig als Adam; wie auch oben [Cap. 2, § 25] gesagt ist, daß Gott mit Adam selbst geredet hat, und ihm ein Gebot gegeben, daß er Evam sollte lehren. Darum will er sprechen: Adam hat es wohl gewußt und verstanden, sie aber war einfältiger, und dem listigen Teufel zu schwach, und versähe sichs nicht; aber Adam versähe sichs wohl, der hätte wohl sollen und können wehren, wenn er es hätte wollen thun.

6. Also siehst du hier, wie der Teufel sich verstellen kann, daß man ihn nicht kennt, und wie er alle Anfechtung treibt. Das laßt uns sonderlich merken. Denn wie er da gethan hat, so thut er noch immerdar in denen, die Christen sind. Darum ist es noth, daß sie gar fleißig machen, und Achtung darauf haben, daß er sie nicht betrüge. Denn die Seinen, die er gar besessen hat, ficht er nicht an.

7. Aufs erste greift er den Menschen an, da er am schwächsten ist, nämlich die weibliche Person, das ist, Evam, und nicht Adam. Denn alle seine Anfechtungen also gerichtet sind, daß er einbricht, wo wir schwach und nicht wohl verwahrt sind. Wo er Adam angegriffen hätte, würde er ihm wohl eine andere Antwort gegeben haben. Davor fürchtete er sich, und dachte also: Ich will die Sie zum ersten angreifen, vielleicht möchte ich ihn auch durch sie hernach zu Fall bringen.

8. Da ist uns nun seine eigene Farbe abge
\quellseite{48}
malt. Wo er sieht, daß du am schwächsten bist, dahin du am meisten geneigt bist, zur Hoffahrt, Geiz, Zorn oder Unkeuschheit, oder dergleichen, da tastet er dich an; und kitzelt dich am meisten, da du hitzig bist; greift dich an, da du weich bist. Darum sollten wir gerüstet und wacker sein, daß wir seine Schalkheit lernen kennen, und vor ihm [uns] hüteten, wie Paulus [2 Cor. 2, 11.] sagt: „daß wir nicht übervortheilet werden von dem Teufel; denn uns ist nicht unbewußt, was er im Sinn hat.“ Also hat uns ihn GOtt verrathen, und uns gewarnet, daß wir uns der Tücke und Schalkheit wohl zu ihm versehen sollen.

9. Zum andern. Weiter siehe, wie schalkhaftig er das Weib angreift. Hebt also an wider sie zu reden: „Ja, sollte GOtt gesagt haben, ihr sollt nicht essen von allerlei Bäumen im Garten?“ Mit diesen Worten schlägt er GOttes Gebot in Wind dahin, und redet so leichtfertig davon, als sollte er sagen: Meinst du, daß GOtt so närrisch sei, daß er das sollte verboten haben? Ich kann das Hebräische nicht wohl geben, weder Deutsch noch Lateinisch, es lautet eben das Wort Aphti, als wenn einer die Nase rümpft und einen verlacht und verspottet. Damit will er sie herum rücken, daß sie denken soll: Ei, es muß dennoch wahr sein. Also rückt und reißt er sie von dem Wort GOttes, welches straks da steht. So ist es denn verloren. So lange das Wort im Herzen stand, da lebte sie, und blieb auch stehen.

10. Darum denkt er: das Wort muß ich ihr vor allen Dingen nehmen; und braucht die Tücke, daß er ihr einen solchen Wahn macht, daß sie meint, es müsse nicht so sein. Ei, meinst du, spricht er, du Närrin, daß es GOtt geboten habe? Bringt es so listig vor, als wollte er sagen: Es sind so viel hundert Bäume im Garten, und hat dir von keinem verboten zu essen, sollte er denn eben den Baum verboten haben? Sollte er dir einen einzigen Baum verbieten, so er dir so viel hundert gibt? Also führt er das Weib an, daß sie denkt, Adam wird es nicht recht verstanden haben.

11. Das ist die rechte schwarze, ja, die weiße Farbe des Teufels. Denn er ist ein heller, lichter Teufel, der uns nicht mit groben Sünden anficht, sondern mit dem Unglauben. Denn wenn er den Glauben umgestoßen hat, so hat er gewonnen. Der Mensch muß GOttes Wort haben, und daran hangen mit dem Glauben: so bald er ihm nun dasselbe entzücken\footnote{entzücken = entreißen, entrücken.} läßt, so ist keine Hülfe mehr da. Also nimmt hier ihm der Teufel das Wort und den Glauben, daß er wankt, und denkt: Wer weiß, es möchte wohl nicht wahr sein. So bald nun das in Zweifel gestellt wird, daß man denkt, ob es recht oder unrecht sei, so ist es verloren. GOtt will nicht eine wankende Seele haben, die da sage: Ja, wenn es wahr wäre.

12. Darum merke, wie der Teufel thut, daß er nur den Glauben angreift. Die Heiden, Ungläubigen und Unchristen ficht er nicht an, die hangen an ihm wie Schuppen; aber wenn er die ficht, die da GOttes Wort, den Glauben und Geist haben, denen kann er nicht zu; er weiß wohl, daß er nicht gewinnen kann, wenn sie schon straucheln. Sieht wohl, wenn gleich einer in grobe Sünde fällt, daß es damit nicht verloren ist, denn er immer wieder aufstehen kann. Darum denkt er, er muß den Schaden anders thun, und das Hauptgut nehmen. Wenn er es denn dahin gebracht hat, daß er zweifelt, ob das GOttes Wort sei, so ist es gewonnen Spiel.

13. Darum ist es noth, daß wir eben darauf sehen. Paulus [1 Tim. 6, 20.] und Christus [Matth. 7, 15. 16, 6. 12.] haben uns genug gewarnet; dazu die ganze Schrift, darinne keine größere Warnung ist, denn daß man sich hüte vor falscher Lehre. Denn GOtt kann alles für gut halten, wie wir straucheln, allein, daß wir bleiben bei dem reinen, lauteren Wort GOttes, das da sagt: Dies ist recht, dies ist unrecht. Solches weiß der Teufel, darum kriecht er da zum ersten ein; so bald er dasselbe hinweg reißt, so kann ihm der Mensch darnach selbst nicht wehren, er muß fallen in alle Laster; wie du siehst, daß er hier thut, als er sie herum rücken will.

14. Sie will es besser machen, und spricht also: Wir essen nicht von dem Baum mitten im Garten, auf daß wir nicht sterben (wie wir es auch noch bisher immer haben wollen besser machen, denn es GOtt geboten hat); und setzt auch einen Zusatz dazu: Wir sollen es auch nicht anrühren. Damit hebt sie schon an zu wanken. Als das der Teufel sieht, fährt er fort, und reißt sie dahin, daß sie GOtt lästert und straks wider ihn steht. Denn wenn er den Menschen dahin
\quellseite{49}
reißt, daß er zweifelt, bringt er ihn bald dazu, daß er GOtt lästere, und spricht: GOtt hat es nicht geboten, oder wenn er es gleich geboten hat, meint er es nicht recht.

15. Darum fährt er weiter, und sagt zu ihr: Ja, „ihr werdet sein wie GOtt, und wissen“ 2c. Diese Worte haben viel in sich, daß nicht genug davon zu reden ist. Denn es ist alles so abgemalt, wie er fährt mit allen Menschen, die er angreift. Darum hüte dich, daß du nicht meinest, es sei allein mit der Eva so gegangen, noch die Geschichte für eine alte Märe haltest, sondern denke, als wäre es erst die Stunde geschehen. Denn also thut der Satan noch heutiges Tages, von Anfang der Welt her bis zum Ende. So werden wir sehen, und haben es schon vor Augen, wie er mit dem Evangelio handelt, und wie viel Seelen er davon reißen wird und eben also angreifen, daß er von ersten GOttes Wort auf einen Zweifel stelle; darnach fort, daß man es verleugne; und bringt das sonderlich mit ein, wie er hier thut, daß GOtt nicht freundlich sei. Welches die allergiftigste Anfechtung ist, so einem Menschen widerfahren mag, daß er mit solchen Augen GOtt ansieht; welches nichts anders, denn ein Haß GOttes ist, dazu wir mächtig viel Ursach haben, sonderlich wenn wir anheben zu reden von der Versehung, warum GOtt einen erwählt und den andern nicht. Solches treibt der Teufel nun immerdar.

16. Darüber macht er dem Weibe auch böse Gedanken, welches noch viel ärger ist. Denn daß er den Namen und Schein stehen läßt, daß es nicht wider GOtt sei, führt er erst das Weib und uns armen Leute dahin, daß sie meinen, sie sündigen nicht daran, und ist doch im Grunde GOtt verleugnen und hassen. Dazu müssen wir noch immerdar die Erempel, so wir vor Augen haben, ansehen, sonst verstehen wir es nicht. Also fährt er jetzt zu, gibt [dem] Papst, Bischöfen, Pfaffen und Mönchen vor, daß ihr Wesen recht, und nicht wider GOtt sei; so bleiben sie in dem Wahn,\footnote{„Wahn“ von uns gesetzt statt: „Wort“, was uns nicht zu passen scheint. Wegen unserer Correctur vergleiche § 21: „sie stehen auf dem Wahn“ 2c.} daß sie recht fahren; und steht doch GOttes Wort klar da vor ihren Augen, daß sie nicht thun sollen, was sie nicht gewiß sind, daß es GOttes Wort sei. Sagen demnach: Wir müssen wohl, daß GOtt geboten hat, daß man nicht thun solle, was er nicht geboten hat; aber doch möge man wohl etwas mehr und höher GOtt zu Dienste thun, das ihm zu Lobe und Ehren geschehe. Das heißt nun zum ersten, negligere mandatum Dei; darnach, eligere mandatum hominum. Darum ist diese nicht eine alte Anfechtung, sondern währt noch bis an das Ende der Welt.

17. Nun siehe, was weiter geschiehet. Als Eva in dem Wanken steht, und er hat beschlossen, es sei nicht wider GOtt, da hat er gewonnen; der Glaube ist aus und erwürgt, das Wort hat sie verloren. Da spricht der Text: „Und das Weib schaute an“, zum ersten, „daß der Baum fein war, davon zu essen“; zum andern, „und lieblich anzusehen“; zum dritten, „daß es ein lustiger Baum wäre, weil er klug machte.“ Diese drei unseligen Begierden hat sie zuvor überall nicht gehabt, ist nun schon in böse Lust und Liebe gefallen, der sie zuvor keine hatte.

18. Darum, wenn der Glaube und GOttes Wort hin ist, ist nicht zu denken, daß man dann könnte aufhalten die böse Lust und Liebe; der Vorwitz ist da, und eitel sündige, böse Neigung. Zuvor, da sie im Glauben stand, und voll Liebe GOttes im Herzen steckte, sahe sie nicht solches, daß der Baum sonderlich lustig wäre, oder sonderlich klug machte, sondern war ihr einer wie der andere; jetzt aber ist ein Unterschied, daß keiner so schön sei als der. Was GOtt geboten hat, davon will sie nicht; was er aber verboten hat, das will sie, und wollte nun gern klug werden.

\begin{quote}
Und brach der Frucht ab, und aß, und gab ihrem Manne auch davon, und er aß.
\end{quote}

19. Da folgt nun das Werk. Sie hätte nicht gegessen, wäre sie nicht zuvor todt gewesen. Der Glaube war schon dahin, und steckte voll Sünde und böser Lust. Das ist aber der Jammer gar, daß sie Adam auch gibt, und er mit ihr davon isset. Denn es\footnote{Wittenberger und Erlanger: er.} hätte keine Noth gehabt, wo er geblieben wäre, GOtt hätte wohl ein ander Weib geschaffen. „Adam“, spricht Paulus [1 Tim. 2, 14.], „ward nicht verführet, sondern das Weib.“ Aber, daß er auch übertritt, macht die Sünde zumal schwer und greulich. Sie war eine Närrin, leicht zu verführen, mußte es nicht anders; er aber hatte GOttes Wort vor ihm, das mußte er wohl, und sollte sie gestraft haben; so steht er da, sieht zu, und isset auch, verwilligt Muthwillens in des Teufels Rath.
\quellseite{50}

\begin{quote}
B. 7. Da wurden ihrer beider Augen wacker.\footnote{Wittenberger: aufgethan.}
\end{quote}

20. Da folgen nun die groben Sünden. Zum ersten werden ihre Augen wacker, wie der Teufel gesagt hatte, sehen und fühlen, daß sie nacket sind. Da war nun unmöglich zu wehren allen Gliedmaßen, die sie hatten, noch der bösen Lust zu steuern. Sie sahen sich beide an mit böser Lust und unkeuschen Begierden, als sie nacket waren, welche ihnen zuvor unbekannt waren. Sie waren abgefallen, und GOtt ungehorsam worden; da ward hernach auch alles ungehorsam, was am Leibe war, daß sie es nicht zähmen konnten, weder Gedanken, noch Gliedmaßen. Das ist nun an uns geerbt, und bleibt noch also; wie sie gewesen sind, so sind alle ihre Kinder. Wo nicht Glaube und Christen sind, da ist es ungewehrt, daß sie ohne böse Lust und Begierden seien, sonderlich den Gliedmaßen zu wehren, die zum Zorn und Unkeuschheit dienen: Augen, Ohren, Zunge, und alle anderen Gliedmaßen.

21. Also ist nun Adam und Eva mit aller Frucht und Kindern verdammt, stecket in der Sünde und wird keines ausgenommen. Wir sind alle Vater und Mutter ähnlich, und bringen eben dieselbe Seuche und Krankheit mit uns. So geht noch heutiges Tages der Teufel mit uns um, daß er uns führe in die Hauptanfechtung, so den Glauben betrifft; wenn er das Wort und den Glauben gewinnt, so hat er’s alles. Wie wir sehen, wo Pfaffen und Mönche nicht im Glauben leben, wie sie denn fast alle sind, daß kein geiziger, unkeuscher, zorniger Wolf ist, und niemand so voll Laster steckt. Denn sie stehen auf dem Wahn, daß sie wollen der Sache mit Werken helfen, sehen ihren Unglauben und die Krankheit nicht, wissen nicht, woran es ihnen fehlt, fahren zu, und wollen ihnen auch Schürze machen, damit sie sich decken; können aber ihrer Natur nicht wehren, es bricht doch heraus, daß es an Tag kommt, was sie sind, und sie es selbst müssen bekennen. Wenn man ihnen denn sagt, daß es an GOttes Wort und Glauben fehlt, so glauben sie es nicht.

22. Also fahren nun die Zwei zu. Da sie sehen, daß sie nacket sind, fühlen die schändliche Lust in ihrem Fleisch, und können ihnen nicht helfen, gehen sie hin, und machen ihnen Schürze oder Gürte von Blättern geflochten, die sie nun sich her gürteten und den Leib bedeckten. Das bedeutet nun nichts anders, denn daß wir oft gesagt haben, wie dieselbigen Heiligen, so sie den Glauben verloren haben und sehen, daß sie in Sünden stecken, wollen ihnen helfen, sich decken und beschönen mit Werken, die sie erdenken können. Ist aber hier schön ausgedrückt, daß es nicht hilft, wiewohl sie Schürzen machen. Als sie sahen, daß sie nacket waren, und die böse Lust und Liebe fühlten, mochten sie ihr dennoch nicht wehren, noch sie dämpfen; [sie] dachten, sie wollten eine andere Weise angreifen, es half aber nicht; sie konnten wohl einander aus den Augen kommen, daß eines das andere nicht sähe; aber die Lust hörte nichts desto mehr auf.

23. Also geschieht es noch mit Mönchen und Pfaffen. Da haben die Leute, so ihre Sünde fühlten, und nicht los werden konnten, gemeint, sie wollten der Sache helfen mit strengem Leben, und sich ihr Lebenlang in die Klöster verschlossen; wie lange sie aber darinnen sind, fühlen sie, daß es nichts hilft, ja nur ärger wird. Es scheint wohl also, daß die Leute meinen, sie seien fromm und heilig, ist aber nichts, denn daß man sich vor den Leuten deckt; vor GOtt bleiben sie eben Schälke in der Haut, wie vor. Wo uns nicht wieder durch GOttes-Wort geholfen wird, so ist es alles verloren.

24. So haben wir nun den elenden Fall, den Adam und Eva gethan haben, darinne wir alle stecken, also, daß es niemand wehren kann. Denn, ob gleich das weltliche Schwert, item, Vater und Mutter durch dem Werck wehren, so ist doch dem Grunde des Herzens nicht gewehrt, sondern muß allein das Wort wieder kommen, davon wir gefallen sind, und uns aufrichten.

25. Darum ist uns diese Anfechtung zur Warnung geschrieben, daß wir uns vorsehen, und nicht vom Worte lassen reißen, wie sie gethan haben; denn es alles miteinander zu thun ist um das Wort und den Glauben. Als sie nun also vom Teufel betrogen sind und in Sünde und Tod gefallen, und ihren Unglauben und Ungehorsam fühlen in allen Gliedern, am Leibe und Seele, daß sie die Schande decken mußten mit Schürzen und Blättern, folget nun weiter, wie sie GOtt wieder aufgerichtet, und zu Gnaden hat kommen lassen.
\quellseite{51}
V. 8—13. Und sie höreten die Stimme GOttes des HErrn im Garten gehen, da der Tag\footnote{Handglosse: Tag kühle war. Das war um den Abend, wenn die Hitze vergangen ist; bedeutet, daß nach gethaner Sünde das Gewissen Angst leidet, bis daß GOttes gnädige Stimme komme, und wieder kühle und erquicke das Herz; wiewohl sich auch die blöde Natur entsetzt und fleucht vor dem Evangelio, weil es das Kreuz und Sterben lehrt.} kühle worden war, und Adam versteckte\footnote{Handglosse: Adam versteckte. Adam heißt auf Hebräisch Mensch, darum mag man Mensch sagen, wo Adam steht, und wiederum.} sich mit seinem Weibe vor dem Angesichte GOttes des HErrn unter die Bäume im Garten. Und GOtt der HErr rief Adam, und sprach zu ihm: Wo bist du? Und er sprach: Ich hörte deine Stimme im Garten und fürchtete mich; denn ich bin nacket, darum versteckte ich mich. Und er sprach: Wer hat dir’s gesagt, daß du nacket bist? Hast du nicht gegessen von dem Baume, davon ich dir gebot, du solltest nicht davon essen? Da sprach Adam: Das Weib, das du mir gegeben hast, gab mir von dem Baume, daß ich aß. Da sprach GOtt der HErr zum Weibe: Warum hast du das gethan? Das Weib sprach: Die Schlange hat mich also aufgesetzt,\footnote{äußern = betrügen, verführen (illudere).} daß ich aß.

26. Jetzt ist genugsam gesagt, daß die höchste Sünde in diesem Fall gewesen ist, daß Adam und Eva durch den listigen Betrug des Teufels geführt sind vom Worte GOttes. Da ihnen nun das entrückt ward, da war das Leben nimmer da; denn wo das Wort nicht ist, da ist auch das Leben nicht. Da das Leben hin war, fühlten sie die Frucht des Todes, daß sie böse Lust gewonnen hatten.

27. Nun aber siehst du wiederum in diesem Stücke, daß sie GOtt wieder aufrichtet vom Tode in das Leben durch das Wort. Denn da kommt ein anderes GOttes Wort, das sie wieder errettet; aber nicht ehe, denn sie den Tod zuvor fühlen. Und ist der Text wohl zu merken, denn es nicht zu denken ist, als wäre es vergangen Ding, und ginge nimmer also; denn wo es nicht also geht, da wird nimmer kein Christ. Das wollen wir sehen.

28. Da die zwei todt waren, GOttes Wort verloren hatten, und in der Sünde steckten, fühlten sie es dennoch so bald nicht, sondern gingen hin, wie der Text sagt, gürteten und deckten sich zu; fühlten noch nicht, wie übel sie gehandelt hatten; waren doch voll Früchte des Todes und böser Lüste und Liebe, der Kitzel war noch nicht vergangen. Da sie aber höreten die Stimme GOttes, der im Garten ging, sagt Moses, da die Hitze [vor]über und kühl war, um Vesperzeit, da versteckten sie sich, wollten sich nicht lassen sehen; da ward ihnen der Kitzel wohl vertrieben, daß sie solcher Jammer und Noth ankam, daß sie vergaßen, ob sie nacket wären oder nicht.

29. Denn, wie wir noch sehen, wenn es dahin kommt, daß einer verurtheilt wird zum Tode, geht einem solche Angst unter die Augen, daß er nicht weiß, ob er Mann oder Weib sei. So geschah ihnen auch; da sie sich fast gedeckt hatten und Schürzen gemacht: so bald sich GOtt hören läßt, erschrickt das Gewissen, wollten sich gern verbergen und entlaufen, und wird ihnen der Garten zu enge, daß sie nicht wissen, wo sie bleiben sollen. Da fühlten sie erst, was sie gethan hatten, sahen die Sünde, Tod und Hölle vor Augen; da war aller Jammer und Herzeleid. Sie wollten fliehen und kriechen, wo sie ein Loch fänden,\footnote{Erlanger: finden.} da sie vor GOtt bleiben könnten; aber da ist kein Raum. Je mehr sie sich versteckten, je näher GOtt kam, und so nahe, daß er sprach: „Adam, wo bist du?“

30. Das konnten sie nicht leiden, und standen beide zu der Stunde freilich mitten in Todesnöthen, daß GOtt so unfreundlich mit ihnen umgeht, und sie nicht mehr kennen will; als [ob] er sagen wollte: Zuvor warst du mein Kind, nun habe ich dich verloren. Da sie das hören, daß er sich so fremd gegen ihnen stellt, da steckten sie gewiß in der Hölle. Denn also stand ihr Herz: Siehe, GOtt ist dir feind, weil dir aber GOtt feind ist, so sind dir alle Creaturen feind und wider. Da ist ihnen alles zu enge worden. Desgleichen geht es noch täglich, wenn GOtt einen Menschen bekehren will von Sünden, daß er ihn erstlich in solch Schrecken und Angst führt.

31. Nun siehe, was sollte Adam thun? Als er sieht, wie er genarret hat, fährt er zu, will sich beschönen, und narrt erst gröber, daß genug zu sehen ist, wenn man gleich keine anderen Exempel hätte, wie närrisch die Natur ohne GOttes Wort handelt. Glaube das ungeziweifelt: wenn sie hätten alles thun können, was immer zu thun wäre, sie hätten wohl drei Hölle daran gewagt. Und steht der freie Wille da
\quellseite{52}
am höchsten, daß er freilich nimmermehr thun kann, was er niemals nicht gethan hat. Denn da standen sie in der Hölle, daß sie mußten ihr Bestes und Höchstes versuchen, wie sie heraus möchten kommen; aber da ist keine Hülfe noch Rath zu finden. Siehe nur, wie er thut. So fromm ist er nicht, daß er Gott die Ehre könnte geben und sprechen: Ach Herr, ich habe gesündigt! und hätte um Gnade und Hülfe gebeten; sondern so denkt nur das Herz: Ach daß er doch die Sünde nicht wüßte! wollte sie gern schmücken und helfen; darum spricht er: „Ich hörte deine Stimme im Garten, und fürchtete mich“; als sollte er sagen: Ich fühle wohl, daß vor dir nicht gilt noch hilft kein Schürzen und Decken; da bin ich gar bloß. Darum sagt er: Ich verkroch mich. O du Narr! ist das die Weise, vor Gott zu handeln, daß man vor ihm fliehe, so man sollte mit allen Tieren zu ihm kriechen?

32. Da hebt nun Gott mit ihm an, stößt ihn noch tiefer in die Hölle, daß er gar zu Grunde geht. „Wer hat dir gesagt“, spricht er, „daß du nacket bist“, warst du doch zuvor auch nacket? Das war zu tief gefragt, daß er nicht mehr antworten konnte, und war nun beschlossen und stand in der tiefen Hölle, allerdings von Gott verdammt. Denn also schließt er ein Urtheil: Weil du fühlst, daß du nacket bist, und dich vor mir fürchtest und fleuchst, so mußt du wider mich gehandelt haben, und ungehorsam sein, daß du und ich uneins sind. Was thut er aber? Er fährt zu, und will die Schuld nicht ihm selbst geben, sondern auf das Weib, ja auf Gott selbst bringen. Als sollte er sagen: Ach! hättest du mir das Weib nicht gegeben, so wäre ich wohl fromm geblieben. Das ist so viel geredet: Wärest du so klug und so fromm gewesen, solltest du das Weib nicht geschaffen haben. Was ist das anders gesagt, denn so viel: Du hast selbst gesündigt? Also antwortet er der göttlichen Majestät.

33. Nun redet niemand solche Worte gegen Gott, denn ein Gotteslästerer, der da fühlt, daß er verdammt ist. Denn wo ein rechtes Herze ist, das bekennt die Sünde, und gibt Gott je so viel Ehre, daß es ihn nicht straft. Dieser aber fährt zu und straft ihn, gibt ihm die Schuld. Als sollte er sagen: Soll ich ja verdammt sein, so bist du schuldig. Denn Eva war Gottes Werk, das muß es auf sich nehmen. Er sollte aber gesagt haben: Du hast mir ein Weib gegeben, daß sollte ich regiert haben, und nicht ihr, sondern sie mir gehorcht haben; das habe ich nicht gethan. So kehrt er es um, und legt die Schuld auf Gott, macht es jetzt viel ärger denn zuvor, da er vom Baum aß, daß er werth wäre, daß ihn Gott ewiglich hätte in die Hölle gestoßen. Da siehe nun, was der freie Wille und Vernunft thut, wenn es zum Treffen kommt, sonderlich in Todesnöthen, wie fein sie ihr selbst helfen kann, wie sie mit Gott zankt, was sie ihm für Ehre gibt, und wie freundlich sie mit ihm redet, daß sie sich nur je länger je tiefer in die Hölle versenkt.

34. Nun läßt ihn Gott in solcher Noth stecken, und wohl in die Hölle sinken; denn er steht gänzlich auf der Meinung, kann nicht anders fühlen noch denken, denn er sei ewig verdammt; hat keine Hülfe noch Trost.

35. Wo bleibt aber nun das Weib? Sie muß auch hindurch. Darum hebt er abermal an, und fragt: warum sie es gethan habe? schiebt noch den Trost immer auf, läßt sich nicht hören, daß er helfen wolle. Nun ist sie eben so närrisch als Adam, will auch die Schuld nicht tragen. Es ist eine hohe Versuchung, darinne auf das allerheimlichste die Vergebung wird beschuldiget. Die Schlange hat mich aufgesetzet, spricht sie. Als sollte sie auch sagen: Warum hast du die Schlange geschaffen? Da du so ein kluger Gott warst, und solches wußtest, hättest du es doch wohl können unterkommen? Das ist greulich und erschrecklich zu hören.

36. So sind sie nun beide verdammt. Nun aber kommt wieder Trost und Gottes Wort; da steigt Christus vom Himmel herab und hilft; da geht ein ander Wort an, wie folgt:

\begin{quote}
B. 14. 15. Da sprach Gott der Herr zu der Schlange: Weil du solches gethan hast, seiest du verflucht vor allem Vieh, und vor allen Thieren auf dem Felde. Auf deinem Bauche sollst du gehen, und Staub essen dein Lebenlang. Und ich will Feindschaft setzen zwischen dir und dem Weibe, und zwischen deinem Samen und ihrem Samen, derselbe soll dir den Kopf zertreten, und du wirst ihn in die Ferse beißen.
\end{quote}

37. Zum ersten, daß er zu der Schlange spricht: Du sollst auf deinem Bauche gehen, ist schon halber Trost; das lautet bereits, als sollte er sich lassen merken, er wolle sie nicht ewig verdammen. Als sollte er sagen: Ich urtheile euch noch
\quellseite{53}
nicht gar zum Tode, aber es verdrießt mich, daß ihr es gethan habt.

38. Da kommt aber der rechte Trost, da er spricht: „Ich will Feindschaft legen zwischen deinem Samen und des Weibes Samen; er wird dir den Kopf zertreten, du aber wirst ihn in die Ferse beißen.“ Dieser Spruch ist wohl werth, daß man ihn fleißig herausstreiche. Denn darinne ist das Wort des Lebens, davon sie wieder lebendig werden; darinne wird ihnen mit dürren, klaren Worten zugesagt, das Weib soll einen Samen bringen. Weibsname aber heißt, wie die Schrift redet, ein natürlich Kind, vom Weibe geboren; ein Kind, das Fleisch und Blut von der Mutter bringt. Darum will er so sagen: Ich will eine natürliche Frucht schaffen, vom Weibe geboren. Drückt aber nicht aus, ob es ein Mann oder Knäblein sein solle, sondern nur, daß die Mutter von dem mag sagen, das ist mein Kind, und wiederum das Kind sagen könnte, dies ist meine Mutter. Dasselbe soll der Schlange Kopf zertreten; das ist, den Schaden, den sie angerichtet hat, soll er unterdrücken, mit Füßen treten, und zerstören. Da das Adam hörte, kam er wieder aus der Hölle, und ward wieder getröstet.

39. Dies ist nun der Glaube, daß der Same alle Kraft des Teufels sollte nehmen und zertreten, daß sie untergehe. So hat nun Adam diese Worte gefaßt, und so gedacht: Das redet Gott, der leugt nicht; so haben wir ja zu hoffen auf einen Menschen, welcher der Schlange Kopf\footnote{Erlanger: den Schlangen-Kopf.} zertreten soll, die uns getreten hat. Darauf sind sie geblieben und gestorben; dazu auch die Väter, so nach ihnen kommen sind, haben alle darauf gewartet und geglaubt, und immer davon gepredigt, daß eine Frucht kommen würde und der Schlange Haupt zertreten. Und ist in diesem kurzen Spruche eingeschlossen alles, was das Evangelium und unser Glaube in sich hält, als, daß eine Auferstehung und ein ander Leben nach dem Tode kommen soll; item, daß man allein durch den Glauben fromm und selig werde [Röm. 3, 28. Gal. 2, 16.]; darnach, daß an diesem hangt, daß kein Mensch durch seine Werke vor Gott gerecht werde; dazu auch, daß kein Mönch- oder Nonnenstand seliglich ist.

40. Denn der Spruch schließt so viel in sich. Dein Schurz, dein Vermögen wird alles nicht helfen, denn der Teufel hat dich ganz unter sich gebracht; so ist nun keine Hülfe, sondern alles verloren. Soll dir aber geholfen werden, so ist kein ander Mittel, denn der Same des Weibes. Es ist ein starker Spruch, der alles zu Boden schlägt, was anders gepredigt wird; ist schon\footnote{Jenaer und Erlanger: schon.} beschlossen, daß man verzweifeln und verzagen müsse an allem Vermögen, und allein hangen an dem Samen, der es allein thut. Denn wenn ein anderer Weg zur Seligkeit wäre, so müßte der Text falsch sein.

41. Darum siehe, wie das Alte Testament so tapfer von Sachen redet. Da steht, daß Adam ein Christ ist gewesen, schon so lange vor Christi Geburt, denn er eben den Glauben an Christum gehabt hat, den wir haben. Denn die Zeit macht keinen Unterschied des Glaubens: der Glaube ist einerlei vom Anfang der Welt bis an das Ende. Darum hat er eben das empfangen durch seinen Glauben, das ich empfangen habe. Christum hat er nicht mit Augen gesehen, wie auch wir; er hat ihn aber im Worte gehabt, so haben wir ihn auch im Worte. Das ist allein der Unterschied in dem, daß es da sollte geschehen, jetzt aber ist es geschehen. Der Glaube ist eben einerlei; so sind alle Väter, eben wie wir, durch das Wort und Glauben rechtfertiget worden, und auch darinne gestorben. Das ist der Hauptspruch in diesem Texte. Ich will aber hier lassen anstehen, wie das geredet ist: „Die Schlange wird ihn in die Ferse beißen“, denn es wird hernach folgen; item, ob die Schlange zuvor auch getrieben habe, und daß sie jetzt die Erde fresse. Denn wir müssen bei dem Hauptstücke bleiben, da die Macht an liegt.

\begin{quote}
V. 16. Und zum Weibe sprach er: Ich will dir viel Kummer schaffen, wenn du schwanger bist; du sollst deine Kinder mit Kummer gebären, und du sollst dich halten vor deinem Manne, und er soll dein Herr sein.
\end{quote}

42. So fährt Gott nun fort, und greift das Weib an, und straft sie auch. Und ist eine seine Ordnung. Zum ersten fordert er Adam, darnach das Weib, zum dritten die Schlange; das fährt er hier mit der Strafe um. Dem Weibe gibt er ihre Plage, fährt aber säuberlich und schont ihrer, absolvirt sie vom Jammer der Seele, legt die Strafe auf den Leib, wie auch Adam,
\quellseite{54}
und spricht: „Ich will dir viel Kummer schaffen, wenn du schwanger bist“; darnach: „Du sollst deine Kinder mit Kummer gebären“; zum dritten: „Du sollst dich bücken vor deinem Manne, und er soll dein Herr sein.“

43. In diesen dreien Stücken siehst du nichts, denn das den Leib betrifft; die Seele ist schon errettet, und Gottes Kind worden. Darum wendet er die ewige Strafe in eine zeitliche und leibliche, legt hinweg die eiserne Ruthe, und gibt dafür einen Fluchschwanz.

44. Die Strafe geht nun über alle diejenigen, die der Eva Töchter werden; ist nicht zu ihr allein geredet. Redet auch also, als sollten sie alle schwanger werden; welche er aber auszeucht mit Gnaden, die haben ihr Theil. Sonst in der Gemeine hin, was Weiber sind, denen legt er das Herzleid auf, daß sie viel Herzleid haben, wenn\footnote{Erlanger: weil.} sie schwanger sind; nimmt ihnen nicht, daß er zuvor gegeben hat, daß sie sollen fruchtbar sein. Das ist beides Gottes Wort, das niemand wandeln kann. Dazu befiehlt er ihr, daß sie sich vor dem Manne solle bücken, das ist, nicht ihres freien Willens leben; sonst wäre es so zugegangen, daß sie von einander wären gegangen, eines hieher, das andere anderswohin, doch mit Züchten. Jetzt aber kann das Weib ohne den Mann nichts anfangen; wo er ist, muß sie mit, und sich vor ihm bücken. Nun kommt er darnach weiter vom Weibe zu Adam, legt ihm auch eine Strafe auf, und spricht:

\begin{quote}
B. 17—19. Und zum Adam sprach er: Dieweil du hast gehorchet der Stimme deines Weibes, und gessen von dem Baume, davon ich dir gebot, und sprach: Du sollst nicht davon essen; verflucht sei der Acker um deinetwillen, mit Kummer sollst du dich darauf nähren dein Lebelang, Dorn und Disteln soll er dir tragen, und sollst das Kraut auf dem Felde essen. Im Schweiß deines Angesichts sollst du dein Brod essen, bis daß du wieder zur Erde werdest, davon du genommen bist. Denn du bist Staub, und sollst zu Staub werden.
\end{quote}

45. Das ist auch eine sanfte, gnädige Strafe. Ich halte aber, daß der Text allein mache, daß das Land voll Huren und Buben ist, da sonst viel frommer Knaben und Mägdge wären, die sich zusammen in die Ehe begäben. Da\footnote{Erlanger: Da, da.} will niemand hinan. Dem Manne ist nicht auferlegt, daß er Herzleid haben solle mit Kindertragen, das gehört dem Weibe zu; aber daneben ein anderes, daß er Weib und Kind versorgen und ernähren muß; da gehört Mühe und Arbeit zu; da scheuet sich jedermann vor, es will es niemand tragen, und muß doch getragen sein. Denn, nimmst du nicht ein Weib, und issest dein Brod im Schweiß, so nimmt Gott die Strafe, die er dem Weibe aufgelegt hat, und legt sie auf die Seele. Das ist nicht wohl gewechselt. Er will der Seele gnädig sein und helfen, aber den Leib will er wohl plagen. Derhalben, wo Leute sind im Glauben, die schicken sich recht darein und tragen diese Last gern, nehmen Weiber, arbeiten und lassen es ihnen sauer werden. Es ist wohl nicht ein fein Ding nach dem Fleisch, aber ein guter Stand nach der Seele. Noch schreiet die ganze Welt über die Ehe, wie böse Ding es sei. Wem geben sie aber die Schuld? Gott allein, daß er nicht sagt: Du sollst dich auf ein Kissen setzen, im Saue leben, und kein Unglück haben.

46. Darum schließe ich also: Wo man eheliche Leute findet, daß das Weib nicht Unglück hat mit Kindern, denn Mann auch nicht sauer wird, da ist es nicht recht. Rechte eheliche Weib und Mann sollen nicht gute Tage haben, es muß Unglück und Mühe da sein, oder [es] ist vor Gott nicht recht. Gibt dir Gott ein reich Weib oder Mann, daß du gute Tage hast, keine Mühe noch Arbeit da ist, da bist du schon\footnote{Erlanger: schön.} absolvilt von dem Spruch, und stehst nicht wohl. So toll und töricht ist die Welt wider Gott, meint, sie will so ehelich leben, und in solchen Stand treten, darinne sie gute Tage habe, und wohl lebe; so will Gott das Widerspiel haben. Darum, wer ehelich sein will, der muß sich auf den Spruch erwägen, oder des erwägen, daß er zum Teufel fahre. Du wirst es nicht besser machen, denn es Gott gemacht hat; willst du der Welt folgen, und hören, was sie von Gottes Wort und Werken sagt und hält, so mußt du bei ihr bleiben, und zu Gott nicht kommen. Darum schicke dich drein, Gott wird es mit dir nicht ändern, er wolle denn ein Wunderwerk thun.

47. Wenn es nun so geht im ehelichen Stande, in Kummer und Arbeit, so tröste dich deß, und denke also, daß es so sein muß, und
\quellseite{55}
Gott so eingesetzt und haben will, daß ich ehelich soll sein; darum will ich meinen Leib hinnan setzen, und die Mühe auf mich laden in Gottes Namen, willig und fröhlich mich hinein geben. Thust du aber nicht also, und willst es besser machen, so wirst du machen, daß die Seele verloren wird, ob es dem Leibe gleich wohl geht.

43. Derhalben thut Gott der Welt noch eine Gnade, daß er das mehrere Theil wider ihren Willen in ehelichen Stand wirft, und ihnen Mühe und Arbeit auferlegt; wenn sie es nur verstünden. Denn jedermann wollte gern in guten Tagen leben, da trachten wir alle nach; so muß Gott kommen und dasselbe salzen, daß\footnote{Erlanger: da.} es nicht nach unserm Willen gehe.

49. Darum verflucht Gott auch die Erde, daß sie nicht halb Korn trägt, sondern das meiste Theil Dornen und Disteln, welche sonst nicht würden wachsen, da der Mensch nicht\footnote{„nicht“ fehlt in der Erlanger.} arbeiten sollte. Jetzt mengt er es also, daß sie schier das meiste Theil Dornen und Disteln trägt. Darum ist das die Summa, daß er uns also im Zaum will halten, auf daß er uns den Kitzel vertreibe. Das sind die drei Strafen auf die Sünde gelegt: eine greuliche auf die Schlange, und zwei tröstliche auf den Mann und Weib.

50. Also ist\footnote{Wittenberger: sind.} nun zum ersten geschrieben, wie Adam ward geladen und berufen vor Gericht, daß Gott sprach: „Adam, wo bist du?“ Die Worte alle sind kurz, aber fast reich. So viel haben aber\footnote{„aber“ fehlt in der Erlanger.} daraus genommen, die darauf geschrieben haben, daß man niemand solle urtheilen, er werde denn zuvor verhört und überwunden, sintemal Gott Adam (so er doch wohl wußte, daß er sein Gebot übertreten hatte) nicht wollt verdammen, ehe er ihn vor Gericht fordert und Antwort gehört hatte; und damit allen Richtern verboten zu urtheilen, ehe sie den Selbstschuldigen hören. Das soll man denen predigen, die Gewalt haben zu regieren, daß man überall gehen lasse Weise und Form des Rechts. Das lassen wir jetzt fahren.

51. Aber das ist auch hier angezeigt, daß, wenn ein Mensch in Sünde gefallen ist, dennoch Gott ein solcher Mann ist, der nicht sobald folgt mit der Strafe, sondern schiebet sie auf, und stellt sich, als müßte er gar nichts darum. Denn dieweil Adam hier gefallen ist, und da in der Sünde liegt, ist dennoch das rechte Urtheil, so er verwirft hat, noch nicht da. Denn der Text sagt, daß Gott erst auf den Abend kommt und ruft sie. Nun habe ich zuvor gesagt, daß dies alles am sechsten Tage geschehen ist, daß freilich der Mensch nicht lange gestanden, sondern bald um den Mittag gesündigt, und der Herr hernach auf den Abend kommen ist.

52. Das ist alles die Bedeutung, welche alle auf das zukünftige Wesen gerichtet ist, daß in die Welt kommen sollte durch das Evangelium und geistliche Regiment, daß endlich der Mensch kommen muß zu der Erkenntniß, daß er wisse, was ihm fehlt. Weil er sündigt, ist er nicht verzagt; aber hernach, wenn ihn\footnote{Wittenberger: „er“ statt: ihn. In der Erlanger fehlt „ihn“.} das Gewissen trifft, und Gott kommt zu laden, werden die Augen aufgethan, da wird er erst gewahr, wie greuliche, große Sünde er gethan hat. Das will Moses damit anzeigen, daß er sagt, sie haben die Stimme Gottes gehört, da der Tag nun kühl ward und schier vorüber war, das ist, da die böse Lust vorüber kam, da kam Gott, ließ ihnen vorhin Zeit und Raum genug.

53. Item, daß Moses so fleißig beschreibt, daß Gott Adam vor Gericht fordert, und läßt daneben außen, daß er Evam geladen habe. Denn als er hervor ist gekommen, ist sie hernach gegangen. Wo sollte sie hin? Sie war an ihn gebunden, wo er hin kam, mußte sie hinnach. Darum schweigt er davon stille. Die Schlange hat er aber gar nicht genennet, denn sie ist, meine ich, nicht da gewesen. Damit hat, sage ich, der Heilige Geist angezeigt, daß Gott das Amt zu regieren, lehren und zu predigen, der Manns-person befehlt. Denn daß Adam da geladen wird, ist nichts Anderes, denn eine Predigt des Gesetzes, dadurch er erkenne, was er gethan habe, und was er Gott schuldig ist. Die Predigt ist dem Manne, nicht dem Weibe befohlen; wie auch Paulus [1 Tim. 2, 12.] lehrt, so ferne, was christliche Sachen belangt; sonst kann es wohl zuweilen geschehen, daß ein Weib bessern Rath gibt, wie man auch in der Schrift liest. Sonst ist das Amt zu führen, predigen und Gottes Wort lehren, dem Manne befohlen.

54. Daß aber die Schlange nicht gerufen wird, sondern straks das Urtheil über sie geht, ist, daß Gott den Teufel schon verurtheilt hat, daß er
\quellseite{56}
weder das Gesetz noch Evangelium hören noch predigen kann. Darum heißt er die Schlange auf dem Bauche gehen und kriechen; ist schon gar verloren, daß ihm kein Predigen hilft, und keine Hoffnung ist, daß er solle zu Gnaden kommen können.

55. Siehe nun, wie es Gott angreift. Zum ersten ruft er Adam, darnach straft er erstlich die Schlange. Es ist gerade umgekehrt. St. Petrus hat [1. Ep. 4, 17.] also gesagt, daß die Strafe am Hause Gottes anhebe; das Ende aber gehe über die, [so] nicht glauben. Diese Ordnung wird hier auch gehalten, da Adam und Eva auch werden wohl gestraft; aber die Schlange bleibt endlich verdammt. Dahin geht, daß er sagt: „Ich will Feindschaft setzen zwischen dir und dem Weibe“ etc. Da scheidet sich’s in der letzten Strafe, und wird so bald gnädig, daß er die ewige wandelt in eine zeitliche, als ein gnädiger und freundlicher Vater. Die man meint, er werde sie nicht strafen, die straft er am greulichsten; wiederum, welche man meint, daß er werde am härtesten strafen, die straft er am gnädigsten. Es scheint, als habe er den Adam am härtesten angegriffen, und schont doch sein am meisten.

56. So geht sein Werk alles anders denn in der Welt. Die er am ersten angreifen sollte, behält er bis zum letzten. [Er] greift erstlich den Menschen an, daß man meinen möchte, er wäre ärger denn die Schlange. Darnach kehrt er es um, legt die größte Plage auf die Schlange; darnach auf Eva, die geringste auf Adam. Denn dem Weibe gilt es das Leben, dem Manne aber nicht, sondern Mühe und Arbeit. Die Schlange hat ihr Theil, daß sie ewig verdammt ist; das Weib am Leibe; der Mann an der Arbeit, daß er dennoch nicht darf darüber sterben.

57. Das ist Gottes Werk, das sollen wir lernen\footnote{Erlanger: lehren.} kennen. Denn desgleichen handelt er noch, wenn es zum Zagen kommt; da man meint, er sei am gnädigsten, ist er am zornigsten, und wiederum. Also legt er den Gläubigen viel Plage auf, und greift sie hart an; die Strafe geht nur über den Leib und Gut, die Seele errettet er; die Ungläubigen aber verdammt er ewig.

58. Also hat nun die Vermaledeiung gelautet: Verflucht seiest du auf Erden etc. Da ist der Geist verdammt, und keine Gnade mehr. Und ist damit nicht genug, sondern legt auch Feindschaft zwischen der Schlange und des Weibes Samen, der ihm den Kopf zertrete. Das\footnote{Erlanger: Da.} ist das Evangelium gepredigt, dadurch Adam ist wieder lebendig worden, und die geistliche Strafe auf den Leib gewendet wird. Damit ist nun aufgehoben Tod, Sünde und Hölle, alle Kraft und Vermögen des Teufels. Das heißt er den Kopf zertreten, den Schwanz und Bauch läßt er bleiben.

59. Dieser Same ist nun, wie gesagt, Jesus Christus, ein natürliches Kind, vom Weibe geboren [Gal. 4, 4.] und ernährt wie andere, der hat dieser Schlange Haupt zertreten. Darauf hat auch Adam gewartet, aber nicht gewußt, wann es geschehen, oder wie es zugehen sollte. Es sind finstere Worte, aber fast reich. Ist wohl darinne gefaßt, wie es sollte zugehen; es mußte aber der Geist lehren und erklären. So viel haben sie daraus genommen, daß von diesem Weibe, sie sei, wer sie wolle, ein natürliches Kind würde geboren werden, der das ausrichten sollte; wiewohl darinne begriffen ist, daß er mußte von einer Jungfrau kommen [Jes. 7, 14.], weil er ihn zumißt dem Weibe, und ihn allein eines Weibes Samen nennt. Die beide haben es aber nicht also verstanden, wie wir [Cap. 4, 1.] hören werden, da sie sagt: „Nun habe ich den Mann des Herrn überkommen.“

60. Diesen Samen hat nun Gott dem Teufel lassen unterwerfen, daß er meinte, er wollte ihn fressen, also, daß er ihm nahm sein Leben, Ehre, gute Werke, und nichts an Christo blieb, denn Tod, Schande und\footnote{„und“ fehlt in der Wittenberger und in der Erlanger.} Unehre: so gar nahm er ihm alles. Was geschah aber? Der Same war Gott, darum konnte er nicht unterliegen. Der Teufel rang mit einer andern Creatur denn er meinte, versahe sich nicht, daß Gott da wäre, der Herr [des] Teufels und aller Dinge; da vergriff er sich und lief an. Darum gab Gott dies Urtheil über ihn: Weil du den Menschen [1 Tim. 2, 5.], das unschuldige Blut [Matth. 27, 4.], meinen einigen Sohn [Röm. 8, 32.] angegriffen hast, sollst du und der Tod verdammt sein. Du wußtest, daß du nichts an ihm hattest; darum mußt du bezahlen, oder wiederum ihm unterthan sein und zu Füßen liegen [Hebr. 2, 8.]. So muß die Sünde und der\footnote{„der“ fehlt in der Erlanger.} Tod auch unter
\quellseite{57}
than sein, weil sie an ihn gelaufen und kein Recht haben; desgleichen auch Schande, Schmach und Unehre 2c.\footnote{„2c.“ fehlt in der Wittenberger und in der Erlanger.} Da ist es alles gefallen, was an diese Person gelaufen war, und hat den Streit verloren. Er ist ein ewiger König [Jer. 10, 10. 1 Tim. 1, 17. Luc. 1, 33.] und Herr des Lebens [Apost. 3, 15.], Gnade und aller Ehren [Ps. 24, 7. 8.], frißt Schande, Tod, Sünde und Hölle [1 Cor. 15, 54. ff.], es muß gefangen und unterthan sein, oder sich lösen.\footnote{„lösen“ = loslaufen.} Nun kann sich der keines nicht lösen, weil es wider Gott gethan hat; darum hat es alle Kraft und Macht verloren, muß ihm nun unter den Füßen liegen. Wer hätte nun gemeint, daß Moses, der so einfältig daher redet, solche großen Dinge fasse, darinne die Ueberwindung des Todes und alles Herzeleides steht, ja alles, was im Evangelio steht?

61. Nun, wem ist das alles gepredigt? Niemand denn Adam und Eva. Nun geht Gottes Wort nicht vergeblich noch ohne Frucht ab; darum muß es Adam sammt dem Weibe wiedergebracht und geholt haben, und wieder getröstet, daß er darauf gestanden ist, und sich verlassen hat, daß ein Mensch kommen würde, der den Tod und Sünde überwände. Da hat er geglaubt, und ist wieder genesen, und dachte, hier ist wieder Gnade und Friede, Leben und Trost. So siehst du, wie Gott mit so kurzen Worten reden kann, daß man Ein Wort nicht genugsam ausbreiten kann, wenn man gleich lange davon predigt. In den Worten ist begriffen und verfaßt Christi Zukunft von der Jungfrau, Leiden, Sterben und Auferstehen, sein Reich und Evangelium. Wer könnte es so kurz mit kurzen Worten fassen, oder darinne finden? Darum heißt es billig Gottes Wort, daß es anders redet und höher Ding denn Menschenwort.

62. Den andern Theil dieses Spruchs, da er sagt: „Du wirst ihn in die Ferse beißen“, hat St. Paulus genug ausgelegt. Christus, wiewohl er der Schlange ihren Kopf zertreten, und genommen des Todes und Teufels Gewalt und Macht über alle, die an ihn glauben [Hebr. 2, 14.], ist dennoch der Teufel nicht gar todt noch müßig gelassen, daß er nichts schaffe. Was thut er? Er muß noch in die Fersen beißen, das ist, immerdar anfechten. Darum ist des Evangelii Amt nicht [auf] einmal angerichtet, auf einen Augenblick, wie man bisher gepredigt hat, sondern geht so zu: Wenn Gott einen Menschen aufgerichtet hat, läßt er ihn nicht müßig, sondern in steter Uebung gehen, daß er immer zu arbeiten hat. Derhalben, wenn man schon anhebt zu glauben, ist es darum nicht vollkommen; sondern er zertritt den Kopf, so beißt der Teufel ihn in die Ferse, daß er ohne Unterlaß muß im Kampf stehen. Das hat Adam und Eva gehört, darum hat es auch Frucht geschafft.

63. Also ist es auch noch. Wenn wir Christum erkennen, und wissen, daß er Sünde, Teufel\footnote{„Teufel“ fehlt in der Erlanger.} und Tod für uns überwunden hat, ist sie dennoch nicht todt, hört nicht auf, so lange wir auf Erden leben. Das Hauptstück haben wir wohl hinweg, daß der Teufel überwunden;\footnote{Erlanger: daß den Teufel überwinde.} aber weil wir hier auf Erden sind, soll Gott und das Evangelium so viel zu schaffen haben, daß man sich immer mit der Sünde schlage, und dem Teufel, der uns in die Ferse beißt, widerstehe [Eph. 6, 12.]. Der Schlange Schwanz bleibt im Fleisch und Blut, daß wir den Unglauben empfinden im Herzen, Haß, Neid und Geiz, und was mehr für Sünden sind, welche zuvor regiert haben, und der Kopf gemein sind.

64. Das hat St. Paulus oft gelehrt, als, zum Römern, da er [Cap. 6, 12.] spricht: „So lasset nun die Sünde nicht herrschen in eurem sterblichen Leibe, Gehorsam zu leisten seinen Lüsten“ 2c. Wir haben noch böse Lust, die sich in uns regt [Röm. 7, 8.], aber der Geist treibt sie zurück, daß also ein ewiger Streit bleibt in den Christen [Gal. 5, 17.]. Denn es bleibt immer etwas im Fleisch, die Sünde sticht ohne Unterlaß an, daß man sie fühlt, wird aber gedämpft dadurch, daß Christus da ist, und regiert stärker, denn Teufel und Sünde [Luc. 7, 22.]. In denen aber, so den Glauben nicht haben, regiert der Teufel gar, ist noch mit dem Kopf lebendig [Eph. 2, 2.], daß man ihm folgt, und nicht widersteht [Jac. 4, 7.]. Das ist der schöne liebliche Spruch, der hier Adam gegeben wird, durch welchen Gott von ihm nimmt ewige Verdammniß, und gibt ihm die ewige Seligkeit.

65. Womit hat nun Adam solches verdient? Mit großen, mächtigen Sünden. Zum ersten,
\quellseite{58}
daß er fliehen wollte vor Gottes Augen; darnach, daß er sich noch unterstand zu rechtfertigen. Daß man dabei fein sehen kann, was Gott gibt, daß er es frei umsonst gibt, um keines Verdienstes willen, ja, auch um bösen Verdienst, und daß man daran lerne, wie Gott uns thut, daß wir wieder also thun. Er schenkt uns seine Gnade und alle Güter aus lauterer Güte und Liebe: so sollen wir auch unserer Nächsten Götter sein, daß wir auch die ärgsten Feinde lieb haben, und je ärger sie sind, je mehr wir ihnen dienen und Gutes thun sollen [Matth. 5, 44.]. Wie froh, meinst du, daß Adam da gewesen sei? daß ohne Zweifel sein Glaube hervor gebrochen ist, und sich erzeigt hat, und so gedacht: Weil mich Gott der Herr so zu Gnaden annimmt, und mir so reichlich Barmherzigkeit erzeigt, mir, der ich ihn so geschändet und gelästert habe, so will ich auch wieder von Herzen alles Gute thun, auch dem ärgsten Feinde, so ich möchte auf Erden haben. Solche Frucht schafft das Evangelium, wenn es in das Herz kommt, daß der Mensch voll Freuden wird, und jedermann mit Liebe, Lust und fröhlich dient.

66. Als nun die Schlange verflucht, und das Evangelium von des Weibes Samen verheißen ist, folgt dennoch die Strafe, auf den Leib gelegt, beide Adams und Evä, davon genug gesagt ist: daß das Weib Noth und Herzleid haben muß, wenn sie schwanger geht und gebären soll; der Mann Mühe und Arbeit mit der Nahrung, so lange bis der Mensch wieder zur Erde werde, davon er genommen wird. Denn das Stück geht beide auf den Mann und das Weib; wenn gleich das Weib von Kindesnöthen genießet, daß sie entläuft, soll sie dennoch dem Tode nicht entlaufen. So auch, ob der Mann gleich lange arbeitet und sich müht, soll das sein Lohn sein, daß er zuletzt sterbe, und des Unglücks im Leben abkomme.

67. Wiewohl aber das hart geredet ist vor der Welt, so ist es doch, wenn es der Geist ansieht, eine rechte, große Gnade. Denn wo uns die Strafe nicht auf den Hals gelegt wäre, würden wir alle die ärgsten Buben, und niemand fromm bleiben. Darum ist dies alles ein lauter Evangelium, und eben so viel geredet: Ich will euch die Sünde vergeben, und die Seele zu Gnaden bringen, aber dem Leibe zu schaffen geben, daß er nicht zu wild und böse, noch das Fleisch zu stolz werde.

68. Dazu, wenn der Tod nicht wäre, würde die Sünde nimmer untergehen; darum wird eben damit der Sünde endlich gewehrt, und ist sonst kein Rath, ihrer los zu werden. Solche gnädige und heilsame Strafe gibt er uns, daß die Sünde durch den Tod erwürgt werde. Darum sollten wir solches mit Freuden aufnehmen und tragen, als von einem gnädigen Vater; wie denn auch thun die Gläubigen. Denn so gut ist der Vater, daß auch der Tod muß dienen, zu tödten und auszurotten alles Unglück.

69. Derhalben der Tod nun nichts\footnote{Jenaer und Erlanger: nicht.} ist, denn eine lautere Gnade, ja, ein Anfang des Lebens. Denn nachdem er macht, daß die Seele genießet, so muß das leibliche Wesen, was da ist, Krankheit, Fährlichkeit, Mühe und Arbeit, alles dienen zum Besten, daß nicht besser zu wünschen wäre. Wo der Geist ist, findet er so viel böse Lust im Fleisch, das sich nicht will zähmen lassen, daß, wenn er sich lange mit bläuet, [und] doch nichts hilft, muß er selbst wünschen, daß der Leib todt wäre [Röm. 7, 24.]. Daß also der Tod eben zur Arznei gegeben ist, der sich selbst fresse, die Sünde ermürge, und helfe, daß der Geist selig werde. Darum, wie er Adam zuvor mit dem Tode bräuete, und sprach: „Wirst du von dem Baum essen, so mußt du sterben“, so kehrt er es jetzt um, und tröstet ihn damit. Als sollte er sagen: Willst du alles Unglücks los werden, und ewig leben, so stirb. Das sind die mächtigen tröstlichen Sprüche, da unser Evangelium ganz inne verfaßt ist, wiewohl nicht so klar, als es uns gepredigt und ausgebreitet wird. Sie haben den Geist gehabt, wohl reichlicher denn wir, doch haben wir es viel gewisser; denn sie haben nicht können wissen, wer der Christus und Heiland sollte sein, den wir nun kennen. Folgt nun weiter:

\begin{quote}
B. 20. Und Adam hieß sein Weib Heva,\footnote{Handglosse: Heva. Hai heißt Leben, daher kommt Heva oder Hava, Leben oder Lebendige.} darum, daß sie eine Mutter ist aller Lebendigen.
\end{quote}

70. Bisher haben wir gehört, wie Gott der Allmächtige Adam und Eva wieder aufgerichtet hat von ihrem Falle, und wieder gesetzt durch sein göttlich Wort und Evangelium in die Hoffnung des Lebens nach der Seele, daß sie dem Tode entrännen.\footnote{Erlanger: entrinnen.} Nun geht es wieder zum
\quellseite{59}
Leben, bringt sie zusammen, daß sie sich zu Haus setzen und in den Stand, den er ihnen hat aufgelegt. Freilich aber sagt er, wie Adam seinem Weibe einen Namen gibt nach dem Leben, und Ursache, warum er sie also nennt. Denn das Wort Heva heißt hebräisch das Leben. Also sollte er sagen: Er hieß sie das Leben darum, daß, was da leben sollte, von ihr kommen müßte.

71. Da ist erst wieder bestätigt, daß zuvor [V. 16.] geschrieben ist, daß das Weib dem Manne solle unterthan sein; denn so viel Herrschaft wird ihm über sie gegeben, daß sie sich von ihm muß nennen lassen. Daher man die Frauen noch nach den Männern nennt, und nicht wiederum. Das ist nun auch aus gnädigem Willen Gottes geschehen, daß das Weib ohn\footnote{Es bieten zwar alle Ausgaben „im“, doch scheint es uns, daß „on“ (d. i. ohne) gelesen werden sollte, wiewohl nach § 77 das Weib „das Haus regieren soll“. Doch davon ist hier nicht die Rede.} Regiment bliebe, daß sich als eine schwache Creatur nicht regieren kann, noch zum Regiment taugt. Das ist Eine Ursache, warum dies geschrieben ist.

72. Es ist aber auch eine andere dabei, nämlich, daß Adam hebt bereits an einen Schmack zu fühlen vom Leben, weil er dem Weibe den Namen gibt, daß sie Mutter sein solle aller Lebendigen, und redet freilich vom leiblichen Leben. Doch, weil er gehört hat, daß ein Same kommen soll, der die Schlange unterdrücken soll, läßt er sich dünken, daß er forthin werde Kinder haben, die des Samens genießen, versieht sich, daß dies Weib den Samen sollte bringen; er fehlt aber weit. Desgleichen meint sie auch, wiewohl sie es hernach inne wird. Ist aber damit angezeigt, daß sie die Worte Gottes gar gerne gefaßt haben, und mit Freuden empfangen. Darum redet er also davon, und gibt seinem Weibe solche Wahrzeichen und Namen von dem Spruche, und wünscht wohl alle Tage, daß der Same kommen sollte; wie auch die Patriarchen [1 Mos. 49, 18. Ps. 14, 7. Ps. 53, 7. Jes. 45, 8. 64, 1.]. Darum sind die Sprüche Gottes gewiß, aber doch dunkel, also, daß dem Menschen verhalten wird Person\footnote{„Person“ fehlt in der Erlanger.}, Zeit und Weise, wie es gehen soll. Gewiß war es, daß er sollte geboren werden, aber das war dunkel und verborgen, welches Weibes Same er sein würde, welche Zeit, und wie er sollte geboren werden.

73. Also thut GOtt noch immerdar, daß er gewisse Verheißung gibt, hält ihn doch daneben vor Zeit, Person und Weise; wie hernach auch mit Abraham. Der hat die Zusagung, daß von seinem Samen das Kind sollte geboren werden, in welchem alle Welt sollte gesegnet werden [1 Mos. 15, 4. 17, 5. 6.]. Nun konnte er nicht wissen, wie es zugehen müßte; seines Weibes halben war er sicher, aber das Weib mußte er nicht; er meinte wohl, es sollte Sara sein, und harrte sechzehn Jahre, ward aber nichts daraus. Da gedachte sie, sie wäre es nicht, und gab ihm ihre Magd, die trug einen Sohn [1 Mos. 16, 2. 3.], da meinte sie gewiß, er sollte es sein, fehlten aber beide; noch standen sie fest darauf, daß der Same von Abraham kommen würde. Darnach im dreizehnten Jahre kam GOtt wieder, wohl nach zwanzig Jahren, da deutete er erst die Person, daß es Sara sein sollte [1 Mos. 18, 14.]. So führt GOtt seine Werke, daß er uns im Glauben behalte, läßt sein Wort gewiß gehen, will aber nicht, daß man gewisse Zeit oder Person vor[be]stimme, läßt uns wohl darnach denken, aber auch oft fehlen, wie hier Adam, der zu Heva sagt, sie soll eine Mutter der Lebendigen werden, wird aber nichts daraus, denn daß sie eine Mutter bleibt Adams Kinder.

\begin{quote}
V. 21. Und GOtt der HErr machte Adam und seinem Weibe Röcke von Fellen, und zog sie an.
\end{quote}

74. Bisher sind sie gegangen mit Schürzen bedeckt, waren noch nicht so klug, daß sie wußten Kleider zu machen; darum hebt nun GOtt an, zieht ihnen Pelze an, daß sie ganz bedeckt werden. Das ist aber[mal]s ein Stück des Trostes, und ein Zeichen der Barmherzigkeit, daß er sich so freundlich zu ihnen thut, und sie auch versorgt mit Kleidern. Denn die Speise und Nahrung hat er ihnen zuvor gegeben, da er [V. 19.] sagt: „Im Schweiß sollt du dein Brod essen“ etc. Hier sehen wir die väterliche Treue, daß er uns ja nicht läßt, ob wir gleich Sünder sind, sorgt für Futter und Kleider. Wie wir vor Augen sehen, daß der Spruch so stark geht, daß auch die ernährt und gekleidet werden, die keine Sorge dafür haben; als man manch wüstes Kind findet, das nicht denkt noch achtet, was Korn oder die Wolle gilt.

75. Noch steckt der verzweifelte Unglaube so tief in uns, daß wir immerdar Sorge haben, wir werden nicht ernährt. Das macht allein,
\quellseite{60}
daß wir gewiß wissen wollen, wie uns Gott ernähren wolle; also, daß wir das Haus voll Korns, und die Kasten voll Gelds haben, wollen Gott also anbinden an Haus und Kasten: so will er frei und ungebunden sein, weder an Zeit, Person, Stätte, noch dieses oder jenes. Man lasse ihn dafür sorgen, wie er uns ernähren werde, er wird wohl Korn und Geld geben, die Zeit und Maß wohl treffen; daß du nur denkest: Ich will heute arbeiten, werde wohl sehen, woher er es gibt; morgen wieder also; so würdest du inne werden, daß er dich ohne deine Sorge ernähre.

76. Denn er läßt niemand Hungers sterben, der sich auf ihn verläßt [Sprichw. 10, 3. Ps. 34, 12. 13.], wie Christus Matth. 6, 25—34. sagt: „Seid nicht sorgfältig für euer Leben, was ihr essen und trinken werdet, auch nicht für euren Leib, was ihr anziehen werdet. Ist nicht das Leben mehr, denn die Speise? und der Leib mehr, denn die Kleidung? Sehet die Vögel unter dem Himmel an, sie säen nicht, sie ernten nicht, sie sammeln nicht in die Scheuren, und euer himmlischer Vater nähret sie doch. Seid ihr denn nicht viel mehr denn sie? Wer ist unter euch, der seiner Länge eine Elle zusetzen möge, ob er gleich darum sorget? Warum sorget ihr denn für die Kleidung? Schauet die Lilien auf dem Felde, wie sie wachsen; sie arbeiten nicht, auch spinnen sie nicht. So denn Gott das Gras auf dem Felde also kleidet, das doch heute stehet, und morgen in den Ofen geworfen wird, sollte er das nicht viel mehr euch thun? o ihr Kleingläubigen! Darum sollt ihr nicht sorgen und sagen: Was werden wir essen? was werden wir trinken? womit werden wir uns kleiden? Nach solchem allen trachten die Heiden; denn euer himmlischer Vater weiß, daß ihr des alles bedürfet. Trachtet am ersten nach dem Reiche Gottes, und nach seiner Gerechtigkeit, so wird euch solches alles zufallen. Darum sorget nicht für den andern Morgen, denn der morgende Tag wird für das Seine sorgen.“ Als sollte er sagen: Arbeite ein jeglicher des Tages, den er lebt; morgen weiß er nicht, ob er lebe. Lebt er, so arbeite er aber. Was will er denn heute sorgen für morgen?

77. Darum soll man ihm das Sorgen lassen. Die Arbeit und Mühe, die du thust, ist nicht wider den Glauben, sondern ist nütze, das Fleisch zu zähmen; die Sorge aber ist wider Gott. Das Weib soll der Kinder warten, [das] Haus regieren und warten, was Gott mit ihr schaffe. Der Mann deselbigen gleichen arbeite, und sich Gott befehle; der wird ihn nicht lassen, er hat es stark genug verheißen; sintemal ehe Adam oder Eva darauf denkt oder darum sorgt, kommt Gott zuvor, kleidet sie, und schafft ihnen Nahrung. Darum thun wir nicht mehr mit unserm Sorgen, denn daß wir ihn hindern und im Wege liegen.

78. Was nun dies bedeutet, daß Adam sein Weib nennt; item, daß ihnen Gott die Pelze anzieht, will ich müßigen Geistern befehlen. Adam ist ein Vorbild Christi; das Weib seiner Braut, der christlichen Kirche, die er nach seinem Namen nennt; davon sonst gesagt ist. Uns ist genug, daß wir den Glauben aus dem Texte lehren und gründen.

\begin{quote}
V. 22. Und Gott der Herr sprach: Siehe, Adam ist worden als unser einer, und weiß, was gut und böse ist.
\end{quote}

79. Da hat er Adam und Eva abermal ihr Gewissen gekitzelt: denn diese Worte sind hämisch und spöttisch geredet. Als sollte er sagen: Ei, wie fein habt ihr es nun ausgerichtet! Ich habe gesagt: Ihr sollt nicht von dem Baum essen, so habt ihr der Schlange gehorcht, die da sagt, ihr würdet wie Gott sein, und wissen was gut und böse ist. Wie fein ist es nun geschehen? Diese Worte aber, ist nicht zu denken, daß Gott so gesagt habe, als hätte er Lust, also zu spotten im größten Jammer und Elend, sondern das ist damit angezeigt, daß ihm sein Gewissen also gesagt hat, als ihm der Jammer, darein er geraten war, unter Augen geschlagen ist, daß er ihn täglich mußte fühlen. Also geschiehet auch uns noch immer, daß wir in uns müssen schlagen, wenn wir gefallen sind, und unser Herz uns sagt: Ei, wie fein hast du es ausgerichtet!

80. Darum ist es nicht mehr denn so viel gesagt, daß der Mensch sein Lebenlang Reue trägt der Thorheit, die er da begangen hat. Als sollte er sagen: Ja, ich meine, wir sind nun Götter worden! Ach wie sind wir nun so arme, mühselige Menschen worden; das haben mir niemand denn der Schlange zu danken. Das Heulen und Weinen haben sie gehabt, und müssen es alle haben, wollen wir selig werden. Denn den Fall, den Adam gethan hat, müssen wir alle beweinen, klagen und sprechen wie er: Ei, wie
\quellseite{61}
\subsection{Predigt über 1 Mose 3, 22—24}

\subsection{B. III, 141—144}

sein sind wir nun Götter worden! So soll einer gewitzigt werden, der GOtt verläßt und dem Teufel folgt. Das sind, meine ich, die Worte auf das einfältigste gedeutet. Was fühlet er nun weiter?

Nun aber, daß er nicht ausstrecke seine Hand, und breche auch von dem Baume des Lebens, und esse, und lebe ewiglich.

81. Die Worte, sage ich, sind alle so her geredet, wie es Adam fühlt im Herzen, denn da redet GOtt am allermeisten; wenn das Gewissen hört, so hört es recht. Und ist so viel geredet: Adam ist zum Tod geurtheilt, so ist nun ein Baum des Lebens mitten im Garten, darum will ich wehren, daß er nicht auch davon breche und esse, und lebe ewiglich. Wie geht das zu? Also: der Klage nach, davon [§ 80] gesagt ist, daß sie schreien und heulen müssen über die Sünde, daß sie der Teufel so betrogen hat.

82. Also finden und fühlen wir auch, daß der Sache kein Rath ist; wir müssen sterben, da kann niemand schützen noch aufhalten; wir sind zum Tode geurtheilt, so strenge und stark, daß niemand wehren kann; wir müssen hinan, daß wir nimmer kommen können zum Baum des Lebens, daß wir davon brechen und ewig leben. Das geht nun nicht allein über die, die in Sünden sind, sondern auch über die Gläubigen, denn es bleibt noch in ihnen, daß sie die Sünde fühlen, und klagen über den Tod. So ist es kurzum abgesagt mit dem Spruch, daß wir alle sterben müssen; nur todt, todt, und kein Anderes.

B. 23. 24. Da ließ ihn GOtt der HErr aus dem Garten Eden, daß er das Feld baute, davon er genommen ist. Und trieb Adam aus, und lagerte vor den Garten Eden den Cherubim, und ein glänzendes feuriges Schwert, zu bewahren den Weg zu dem Baum des Lebens.

83. Da kommen wir wieder in einen wilden Wald. Eden [Cap. 2, § 19 ff.] ist gesagt von dem Garten oder Paradies, daß dieser Text noch finster ist, und durch niemand erleuchtet, was der Garten sei. Es lautet an etlichen Orten, als sei es nichts denn der ganze Erdboden gewesen, zu der Zeit, da sie noch nicht gefallen waren, da er noch eitel gute Früchte trug, hernach aber verflucht ward. Das ist eine Meinung. Ich traue es aber schwerlich zu erhalten, wiewohl es großen Schein hat, daß viel Worte in diesem und vorigem Capitel schier darauf klingen. Wiederum ist das auch stark, daß GOtt einen sonderlichen Ort gemacht hat, den er genannt hat einen Garten; und sonderlich ist [1 Mos. 2, 8.] ausgedrückt, daß er ihn gepflanzet habe zu Eden gegen dem Morgen; und hier wird es wiederholt, daß er ausgetrieben ist aus dem Garten Eden, und aufs Feld gesetzt, also, daß je das Feld der Garten nicht gewesen ist. Auch wird das Wort Eden mehr hernach kommen, daß es ein Name eines Landes sein muß gegen den Morgen wärts, und nicht der ganze Erdboden.

84. Also zwingt der Text dawider, daß es ein sonderlicher Ort sei, gepflanzt als ein sonderlicher Lustgarten an einem lustigen Orte (denn das Wort Eden heißt auch auf griechisch Lust), daß da der Mensch wohnete vor allen Thieren auf der Erde. Da müssen wir bei bleiben. Wiewohl wir es nicht gewiß treffen können, so ist doch das Sicherste, daß man bei den einfältigen Worten bleibe; wir sind der Sache zu gering. Denn oben [Cap. 2, 8.] hat er gesagt: „GOtt hatte gepflanzet.“ Die Worte müssen wir stehen lassen, daß natürlich gepflanzt sei, wie man sonst pflanzt; also, daß er sagt: Er hat lassen wachsen aus der Erde allerlei; leidet auch nicht, daß man träumt [Cap. 2, § 19], es sei oben über der Erde, nicht ferne vom Monde. Wir sollen GOtt die Ehre geben, so wir es nicht wissen, und nicht thun wie die hohen Schulen, die meinen, es wäre eine Schande, daß sie etwas nicht wissen sollten, oder nichts dazu sagen könnten.

85. So ist nun der Text klar, daß Adam nicht konnte kommen zum Baum des Lebens, daß ihm das Herz sagte, es wäre kein Rath zum Leben; darum ist er daraus gestoßen, und kommt wieder auf das Land und arbeitet. Das ist alles so viel: GOtt hat ihn so weit vom Leben gesetzt, daß das Herz sagt: Es wird nichts Anderes daraus, wir müssen sterben auf Erden.

86. Daß aber GOtt vor der Thür des Gartens gegen dem Morgen lagert den Cherub, und ein glänzendes feuriges Schwert, das sich hin und her wendet, und glänzt wie eine Flamme, hat er darum gethan, sagt er, daß er den Weg verwahrete, daß niemand zum Garten käme; ist allzumal geredet, wie genug gesagt ist, daß der Mensch im Herzen fühlt und die Erfahrung gibt, daß dem Tode nicht zu wehren und keine Hülfe dawider sei.
\quellseite{62}

87. Hier muß aber die Närrin, die Vernunft, ihre Augen blenden, die sich fast bekümmert, was der Cherub und das glänzende Schwert sei. Cherub, was für ein Thier heiße, ist noch heutiges Tages unbewußt; aber bisher ist es dafür gehalten von unsern hohen Schulen, daß es sei der neun Chöre der Engel im Himmel einer. In der hebräischen Sprache findet man nicht mehr davon, denn daß es Flügel habe; es sei für ein Thier, was es wolle. Also lesen wir 2 Mos. 25, 20, daß Moses auf der Lade Gottes zween Cherub machen sollte, die [mit ihren Flügeln den Gnadenstuhl bedecken sollten und]\footnote{Die eingeklammerten Worte sind von uns eingefügt um des Folgenden willen.} ihr Angesicht gegen einander wendeten. Daraus es klar ist, daß die Cherub müssen Flügel haben, wie die Vögel. Was sie aber für Angesichte haben, weiß ich nicht. Darum haben sie gemeint, es seien Engel, wie sie auch die Maler mit Flügeln malen, nach diesem Wort. Der Text schließet nicht, daß der Cherub das Schwert in der Hand habe gehabt und gehalten, sondern gibt nur so viel, daß es vor den Garten gelegt sei neben oder mit dem Cherub.

88. Aber wie dem allen, wollen wir bei dem einfältigen Sinne erstlich bleiben, wie Adam sammt allen Nachkommen verschlossen ist Hülfe, Trost und alle Wege, damit das Leben zu retten oder aufzuhalten wäre. Darum ist der Cherub und das Schwert da, das da blitzt und schreckt, und wenn jemand hinzu gehen wollte, daß es ihn erwürgete; welches er alles bei sich gefühlt hat. Derhalben gefällt mir’s auch nicht übel, daß man dies deutet nach dem geistlichen Verstande, daß es das böse Gewissen sei. Denn Cherubim bedeutet eigentlich das Predigtamt oder Wort, und die mündliche Rede, die man vergleicht einem Fittig; wie auch die Heiden gethan haben, wenn sie das Wort wollten malen. Darum sie auch ihrem Mercurio, der das Reden führt, Fittige angemalt haben, und die Poeten vom Wort sagen, wie es dahin fliegt wie ein Pfeil, und nicht kann widerrufen werden. Also will Cherubim hier heißen, die Predigt des Gesetzes Gottes, das hat neben sich ein glänzendes feuriges Schwert, das ist, ein strenges Urtheil Gottes, das da schrecklich ist dem Gewissen.

89. Ist nun die Meinung: Gott hat Adam ein Wort gesetzt, das war: „Du sollst sterben“; item: „In Schweiß deines Angesichts sollst du dein Brod essen“, und „mit Kummer dich nähren“ etc. Das waren eitel feurige Schwerter, gestrenge Urtheile, daß er wohl sahe und fühlte, daß er es nicht konnte umstoßen, noch da vorüber kommen, mußte daran verzweifeln, wieder zum vorigen Wesen zu kommen auf Erden. Das ist die Deutung, die sich zum schlechten, einfältigen Sinne wohl reimt. Ob man nun das verstehet, daß es leiblich so geschehen sei, lasse ich auch gern zu, denn ich wollte nicht gern, daß man von den Worten wiche. Ob wir nicht alles können ausreden, liegt nichts an; [es] ist genug, daß wir so viel haben, als wir davon gesagt haben. So haben wir das dritte Capitel, wie der Mensch geschaffen, gefallen, und wieder aufgerichtet ist, und das ganze menschliche Leben.

\section{Das vierte Capitel}

V. 1. 2. Und Adam beschlief sein Weib Eva, und sie ward schwanger, und gebar den Kain, und sprach: Ich habe überkommen\footnote{Randglosse: Ueberkommen. Kain heißt, daß man kriegt oder überkommt; Eva meint, er sollte der Same sein, da der Herr von gesagt hatte, daß er der Schlange Kopf zertreten würde.} den Mann des HErrn. Und sie fuhr fort, und gebar Abel, seinen Bruder. Und Abel ward ein Schäfer, Kain aber ward ein Ackermann.

1. Da kommen nun auch Historien und Exempel, darinne wir sehen, daß es so geht, wie es zuvor geschrieben ist. Sonderlich ist hier eine schöne Historie, darinne uns vorgeschrieben wird, was da sei fast mit einander das ganze christliche Leben, von den zweien Brüdern, Kain und Abel. Auf’s erste möchte jemand fragen: Was hat Mosen gelüstet, daß er uns daher schreibt, wie Adam bei seinem Weibe gelegen sei? Es ist darum geschehen, daß man das lesen muß in
\quellseite{63}
aller Welt vor jedermann, wie die Zwey Hochzeit mit einander haben, daß das Gesetz wird angezeigt, daß Moses mannigfaltig hernach anzeigt, daß vor Christi Geburt nicht müßte Jungfrauenstand bleiben, ja, auch verdammt war, darum, daß Gott gesagt hatte, daß vom Weibe sollte der Same kommen. So lange nun der Same nicht kommen war, und nicht ausgedrückt war, welches Weib sein sollte, die den Samen sollte bringen, mußte kein Weib Jungfrau bleiben, sondern alle gewarten, wo Gott den Samen wollte herbringen.

2. So ist hernach geschrieben in Mose: Verflucht sei der, oder die, die nicht Frucht bringt, oder Samen hinter sich läßt. Darum war es eine große Schande, wenn ein Weib nicht fruchtbar war, und ward gehalten für große Strafe und Plage [1 Mos. 16, 1. 4. 29, 31. 30, 1.]. Da werden wir hören viele wunderliche Dinge von den heiligen Weibern, daß sie also haben nöthige Ursache dazu gehabt, daß sie mußten Kinder haben. Wären sonst wohl etliche gewesen, die gerne Jungfrauen geblieben wären, müßten aber alle des Samens gewarten, daß er nicht durch jemand verhindert würde. Sollte nun jemand ausgezogen sein, so mußte es Gott thun durch sonderliches Eingeben des Heiligen Geistes. Wie man liest von Jeremia [Jer. 16, 2.], der es darum gethan hat, daß Gott sprach, er sollte kein Weib noch Kinder haben; sonst dürfte er nicht ohne Ehe geblieben sein.

3. Das hat nun weiter dahin gedient, daß viel Sünden vermieden würden, die sonst geschähen, und ist ohne Zweifel ein züchtiger Leben, auch unter denen, so nicht geglaubt haben, gewesen, denn jetzt. Christus [Matth. 19, 10. 11.] hat es frei gelassen, wir aber haben es zu gemein gemacht, und Gebote darauf geschlagen. Was Jungfrauen sein wollten, sollten allein Christen sein, denn die mögen allein solche Gnade haben; dem andern Haufen ist es nicht gegeben, die sollten unter dem Gebote bleiben, um vieler Sünde und schändliches Wesens willen, vorzukommen. Denn welche diese Freiheit annehmen, so nicht Christen sind, thun es allein darum, daß sie dadurch entfliehen der Strafe, die Gott dem Mann [1 Mos. 3, 17. 19.] auferlegt hat, sich zu nähren im Schweiß seines Angesichts. Was es aber hilft, sieht man wohl; wenn man dem Regen entläuft, so kommt man mitten ins Wasser. Dazu ist es gut gewesen, und wäre noch gut, daß man Knaben und Magd nicht ließe ohne Ehe hingehen. Wer will es aber thun? Wer kann sie sehen, wer Christen oder nicht Christen sind?

4. Item, es ist auch dazu nützlich, daß die Hochzeiten der lieben Väter und Mütter beschrieben sind, daß den künftigen Ketzern das Maul gestopft würde. Denn der Heilige Geist zuvor gesehen hat, daß etliche kommen und aufstehen würden, derer etliche den ehelichen Stand verdammen, etliche verbieten würden [1 Tim. 4, 3.]. Wie zum ersten die Tatianerketzer, die wohl sahen, daß die Leute verderbt waren, die Jungfrauschaft aber\footnote{„aber“ fehlt in der Erlanger.} ein köstlich Ding wäre, darum wollten sie sie\footnote{Wittenberger und Jenaer: sie sich.} so heilig machen, daß sie sollten leben ohne Ehe; gaben vor, wer im ehelichen Stande wäre, der wäre verdammt. Denselben ist hiemit zu begegnen gewesen, daß Adam und Eva, die nun wieder lebendig und heilig waren, dennoch bei einander blieben, und Kinder zeugten natürlicher Weise.

5. Darnach sind aufgetreten, die die Ehe verbieten, das ist Junker Papst und sein Haufe gewesen. Der hat eingesetzt und bestätiget etliche Stände, darinne die Ehe nicht billig noch zuzulassen sei; daßelbe soll ein sonderlich Gottes-Volk sein. Damit hat er die Welt voll geistlicher Leute, das ist, voll Huren und Buben gemacht, und so viel zuwege gebracht, daß Mann und Weib nicht Mann und Weib sei. Da haben sie es hin gebracht, daß schier eine Schande war, daß eine Magd oder Knabe einander zur Ehe genommen haben, als wäre es nicht christlich.

6. Das ist auch hier gewehrt. Es soll keinem Menschen verboten sein, ehelich Leben und Jungfrauschaft; beides soll im neuen Testament frei sein, aber nur den Christen. Im alten Testament war Jungfrauschaft verboten und die Ehe gezwungen; wo aber nicht Christen sind, da ist der Heilige Geist nicht, da denke nur nicht, daß Keuschheit sei; es sei denn, daß ein natürliches Gebrechen da sei.

7. So wird nun Eva schwanger (sagt Moses), und gebiert einen Sohn. Da wird sie froh, meint, sie habe nun genug, und spricht: Das wollte Gott, „nun habe ich überkommen den Mann Gottes“, oder mit Gott. Sie hatte das Wort gefaßt von dem Samen, der den Teufel
\quellseite{64}
sollte zertreten; das schmeckt ihr im Herzen, da lebet sie von, darauf stand alle ihre Begierde. Darum widerfuhr ihr hier ein großer Trost, daß sie einen Sohn überkommt; da dachte sie: Das wird der Mann werden, der wieder zurecht bringe, daß die Schlange verderbt hat. Darum gibt sie ihm den Namen Kain, das heißt „überkommen“. Als sollte sie sagen: Nun ist der mir worden, oder: Nun habe ich das edle Kleinod gewonnen oder gekriegt. Gibt ihm den edelsten Namen, den sie ihm immer geben kann. Bald darnach gebieret sie noch einen Sohn, den nennt sie nicht also; da bekümmert sie sich vor.\footnote{„da bekümmert sie sich vor“ wird wohl heißen: über den macht sie sich Kummer. Die Jenaer Ausgabe bringt die Conjectur am Rande: „da bekümmert sie sich nicht vor“.} Der erste nähm ihr das Mutterherz gar, der war das liebe Kind; dieses aber nimmt sie sich nicht so sehr an, sagt nicht, was aus ihm werden soll, sondern nennt ihn Abel, das heißt Noth, Wehe, Klage und Leid; gleich das Widerspiel. Wird ihr vielleicht wehe gethan haben und sauer worden sein in der Geburt, daß er solchen Namen tragen mußte.

8. Nun siehe, wie sie die zween Söhne aufziehen, daß Abel ein Schäfer wird, Kain aber ein Ackermann, der das Feld baue. Ist wohl zu denken, daß Eva und auch Adam ein Auge auf den ersten Sohn gehabt haben und so aufgezogen, als sollte er es thun; haben ihm auch das Wort gepredigt, das Gott gesagt hatte von dem Samen. Das Evangelium haben die heiligen Väter fleißig getrieben, und ihre Kinder gelehrt. Also hat Adam weltliches und geistliches Amt müssen üben. Darum waren seine beiden Söhne auch Priester, wie wir hören werden.

9. So ist nun der erste Sohn, Kain, so aufgezogen, daß er dachte reich zu werden, und sich zu setzen und zu schicken in das Regiment, wie sich gebührt nach der ersten Geburt; wie Moses hernach [1 Mos. 49, 3.] schreibt, daß dem ersten Sohn allemal zwei Stücke gebühren des Erbes vor den andern, dazu vornehmlich das Regiment und Priestertum, daß er Herr wäre der andern Brüder; also, daß [das] Priesteramt, zu opfern und predigen, dazu auch das geistliche Regiment sein wäre. Das\footnote{Erlanger: da.} hat hier erstlich Kain geerbt, darum ist er gar weit vorgezogen.

\subsection[B. 3—5]{B. 3—5}

Es begab sich aber nach etlichen Tagen, daß Kain dem Herrn Opfer brachte von den Früchten des Feldes. Und Abel brachte auch von den Erstlingen der Schafe und von ihrem Fetten. Und der Herr hielt sich zu Abel und zu seinem Opfer, aber zu Kain und zu seinem Opfer hielt er sich nicht.

10. Da siehe nun Gottes Werk, davon ich oft gesagt habe, wie er regiert, daß er läßt heilige Leute hingehen in ihrem Wahn, und führt es doch gar anders hinaus, denn sie meinen. Da geht der Spruch Salomonis [Sprüchw. 16, 1.], daß ihm der Mensch etwas vorschlägt, aber Gott führt es hinaus. Sie hatten wohl Gottes Wort und den Glauben, wußten aber Person, Zeit und Weise nicht, welche ihm Gott wollte vorbehalten haben. Also thut er auch mit uns. Wiewohl er haben will, daß wir uns gewiß auf ihn verlassen sollen aller Hülfe und Trostes, will er uns doch nicht Zeit, Person, Städte vorstimmen, wie, wo, wann, und wodurch er es schaffen wolle. So folgt nun, wie es sich klar umkehrt, wider Adams Meinung, der auf großer Zuversicht steht, daß dies der rechte Same sein solle.

11. Da schreibt er, wie die Brüder beide geopfert haben. Opfern aber ist eigentlich ein Priesteramt, vor Gott zu treten und ihm zu dienen. Darum, wo ein Opfer geschieht, muß die Person gewißlich ein Priester sein, ob sie gleich ungefalls und unbeschoren ist, wie ohne Zweifel die zwei gewesen sind.

12. Von wem haben es denn diese gehabt, oder wer hat sie es gelehrt, wie sie Gott dienen und opfern sollten? Von niemand denn von Adam; der hat ihnen nichts Anderes gepredigt, denn daß sie des Samens gewarten müßten, der der Schlange Kopf zertreten würde. Solches hat er immerdar getrieben, und ist ihr Evangelium gewesen; daneben haben sie auch gebetet und geopfert.

13. Was haben sie aber geopfert? Kain von Früchten des Feldes; Abel von Schafen und ihrem Fetten. Da sieht man, woher die Opfer im Gesetze kommen sind. Denn fast die Gesetze, die Moses hat, von Adam und der Welt Anfang her entsprungen, und so fort behalten sind und bestätigt; und die Opfer, davon er hernach im dritten Buche schreibt, gemeiniglich von Früchten des Feldes oder vom Fetten sind. So
